% ============================================================
%  2026声学会议论文摘要集 — LaTeX 模板（单篇示例）
%  编译：pdflatex main  （两次编译更新目录页码）
% ============================================================
\documentclass[a4paper,twoside]{ctexart}

\usepackage[top=2.5cm, bottom=2.5cm, left=2.8cm, right=2.8cm]{geometry}
\setlength{\headheight}{14pt}

\usepackage{newtxtext}
\usepackage{newtxmath}
\usepackage{setspace}

\newcommand{\wuhao}{\fontsize{10.5pt}{13.125pt}\selectfont} % 五号，行距1.25倍
\newcommand{\xiaosi}{\fontsize{12pt}{18pt}\selectfont}
\newcommand{\sihao}{\fontsize{14pt}{20pt}\selectfont}
\newcommand{\sanhao}{\fontsize{16pt}{20pt}\selectfont} % 三号，行距1.25倍
\newcommand{\xiaoer}{\fontsize{18pt}{26pt}\selectfont}

\usepackage{fancyhdr}
\pagestyle{fancy}
\fancyhf{}
\fancyhead[C]{\songti\wuhao 2026年全国声学大会}
\fancyfoot[C]{\songti\wuhao\thepage}
\renewcommand{\headrulewidth}{0.4pt}

% 定义论文页面的页脚样式
\fancypagestyle{paperpage}{%
  \fancyhf{}
  \fancyhead[C]{\songti\wuhao 2026年全国声学大会}
  \fancyfoot[C]{\songti\wuhao\thepage}
  \fancyfoot[L]{\rule{3cm}{0.5pt}}
  \renewcommand{\headrulewidth}{0.4pt}
}

\usepackage{indentfirst}
\setlength{\parindent}{2em}
\setlength{\parskip}{0pt plus 1pt}
\setstretch{1.5}

\usepackage{longtable}
\usepackage[hidelinks,colorlinks=false,pdfborder={0 0 0}]{hyperref}

\newcommand{\tocentry}[4]{%
  \begingroup
  \emergencystretch=3em
  \hangindent=5em
  \hangafter=1
  \noindent\makebox[5em][l]{\wuhao #1}%
  {\wuhao\hyperlink{paper:#2}{#3}}%
  \nobreak\leaders\hbox to .5em{\hss.\hss}\hfill
  {\kaishu\wuhao #4}%
  \makebox[3em][r]{\wuhao（\pageref{paper:#2}）}%
  \par\vspace{5pt}%
  \endgroup
}

\newcommand{\tocentrytwo}[4]{%
  \begingroup
  \emergencystretch=3em
  \hangindent=5em
  \hangafter=1
  \noindent\makebox[5em][l]{\wuhao #1}%
  {\wuhao\hyperlink{paper:#2}{#3}}%
  \nobreak\leaders\hbox to .5em{\hss.\hss}\hfill\mbox{}%
  \par
  \noindent\hspace*{5em}%
  \nobreak\leaders\hbox to .5em{\hss.\hss}\hfill
  {\kaishu\wuhao #4}%
  \makebox[3em][r]{\wuhao（\pageref{paper:#2}）}%
  \par\vspace{5pt}%
  \endgroup
}

\newcommand{\printpaper}[9]{%
  \clearpage
  \thispagestyle{paperpage} % 使用论文页面样式
  \hypertarget{paper:#9}{}\label{paper:#9}%
  \noindent{\heiti\xiaosi 专题：#1}\hfill{\songti\wuhao #2}
  \par\vspace{10pt}

  % 标题：黑体三号，居中
  \begin{center}
    {\heiti\sanhao #3}
  \end{center}
  \vspace{0.5\baselineskip}

  % 作者：楷体四号，居中
  \begin{center}
    {\kaishu\sihao #4}
  \end{center}
  \vspace{0.5\baselineskip}

  % 单位：每行一个，前面加序号，去掉点，用括号括起来
  \begin{center}
    {\songti\wuhao #5}
  \end{center}
  \vspace{0.5\baselineskip}

  % 摘要：黑体五号，后空两格，不加冒号
  % 摘要正文：宋体五号，行距1.25倍，段前段后0行
  \begingroup
  \setlength{\parindent}{0pt}
  \setlength{\parskip}{0pt}
  \noindent{\heiti\wuhao 摘要}\hspace{2em}{\songti\wuhao #7}
  \endgroup
  \vspace{\baselineskip}

  % 关键词：黑体五号，左对齐，后空两格
  % 具体关键词：宋体五号，末尾不加标点
  \noindent{\heiti\wuhao 关键词}\hspace{2em}{\songti\wuhao #8}
  \par\vspace{6pt}

  % 页面左下角脚线和通讯作者姓名、邮箱
  \vfill
  \noindent\rule{3cm}{0.5pt}\\[2pt]
  {\songti\wuhao 通讯作者：#6}
}

\newcommand{\ns}[1]{\textsuperscript{#1}}
\newcommand{\ca}{\textsuperscript{*}}

\begin{document}

\thispagestyle{empty}
\vspace*{3cm}
\begin{center}
  {\heiti\fontsize{36pt}{48pt}\selectfont 2026年全国声学大会}
  \vspace{1cm}
  {\kaishu\sanhao 论文摘要集}
  \vfill
  {\songti\wuhao\today}
\end{center}
\clearpage

% ============================================================
%  toc_content.tex — 目录（由 generate_tex.js 自动生成）
%  生成时间：2026-06-26T02:48:27.597Z
%  注意：需二次编译才能正确显示页码
% ============================================================

\begin{center}
  {\heiti\xiaoer 目\quad 录}
\end{center}
\vspace{12pt}

\begin{center}{\heiti\xiaosi A. 物理声学（含声超构材料）}\end{center}
\par\vspace{6pt}
\noindent\makebox[5em][l]{\wuhao A0040}\makebox[\dimexpr\linewidth-5em][s]{\hyperlink{paper:A0040}{\wuhao 非周期结构中的声学传输特性}\dotfill\makebox[2.5em][r]{\wuhao\pageref{paper:A0040}}}
\par
\noindent\hspace*{5em}{\kaishu\wuhao 郭钰宁}
\par\vspace{5pt}
\noindent\makebox[5em][l]{\wuhao A0042}\makebox[\dimexpr\linewidth-5em][s]{\hyperlink{paper:A0042}{\wuhao 基于信号能量的复合材料板声源定位方法研究}\dotfill\makebox[2.5em][r]{\wuhao\pageref{paper:A0042}}}
\par
\noindent\hspace*{5em}{\kaishu\wuhao 马晨宁, 刘金霞, 崔志文}
\par\vspace{5pt}
\noindent\makebox[5em][l]{\wuhao A0073}\makebox[\dimexpr\linewidth-5em][s]{\hyperlink{paper:A0073}{\wuhao 基于双环完美声涡旋的正声学对比度粒子的环形轨道操控}\dotfill\makebox[2.5em][r]{\wuhao\pageref{paper:A0073}}}
\par
\noindent\hspace*{5em}{\kaishu\wuhao 张琪, 郭各朴, 屠娟, 马青玉}
\par\vspace{5pt}
\noindent\makebox[5em][l]{\wuhao A0088}\makebox[\dimexpr\linewidth-5em][s]{\hyperlink{paper:A0088}{\wuhao 拓扑优化Janus型声学超表面的反射波前调控}\dotfill\makebox[2.5em][r]{\wuhao\pageref{paper:A0088}}}
\par
\noindent\hspace*{5em}{\kaishu\wuhao 漆小琳, 刘禧萱}
\par\vspace{5pt}
\noindent\makebox[5em][l]{\wuhao A0108}\makebox[\dimexpr\linewidth-5em][s]{\hyperlink{paper:A0108}{\wuhao 角向-径向二次混合超表面实现声涡旋模式感知}\dotfill\makebox[2.5em][r]{\wuhao\pageref{paper:A0108}}}
\par
\noindent\hspace*{5em}{\kaishu\wuhao 郝国栋, 郭肖晋, 王加慧, 马树青}
\par\vspace{5pt}
\noindent\makebox[5em][l]{\wuhao A0118}\makebox[\dimexpr\linewidth-5em][s]{\hyperlink{paper:A0118}{\wuhao 面向低空噪声抑制问题的超表面波场整形研究}\dotfill\makebox[2.5em][r]{\wuhao\pageref{paper:A0118}}}
\par
\noindent\hspace*{5em}{\kaishu\wuhao 朱萌萌, 何芝华}
\par\vspace{5pt}
\noindent\makebox[5em][l]{\wuhao A0172}\makebox[\dimexpr\linewidth-5em][s]{\hyperlink{paper:A0172}{\wuhao 基于含气泡软媒质的高纯度水下声涡旋束产生}\dotfill\makebox[2.5em][r]{\wuhao\pageref{paper:A0172}}}
\par
\noindent\hspace*{5em}{\kaishu\wuhao 王蓉, 刘京京, 梁彬, 程建春}
\par\vspace{5pt}
\noindent\makebox[5em][l]{\wuhao A0173}\makebox[\dimexpr\linewidth-5em][s]{\hyperlink{paper:A0173}{\wuhao 基于声学时变边界的散射拓扑与频率转换}\dotfill\makebox[2.5em][r]{\wuhao\pageref{paper:A0173}}}
\par
\noindent\hspace*{5em}{\kaishu\wuhao 莫迪}
\par\vspace{5pt}
\noindent\makebox[5em][l]{\wuhao A0176}\makebox[\dimexpr\linewidth-5em][s]{\hyperlink{paper:A0176}{\wuhao 基于幅相解耦超表面的跨介质高保真声全息投影}\dotfill\makebox[2.5em][r]{\wuhao\pageref{paper:A0176}}}
\par
\noindent\hspace*{5em}{\kaishu\wuhao 宋怡璠, 谭杨, 刘京京, 梁彬, 程建春}
\par\vspace{5pt}
\noindent\makebox[5em][l]{\wuhao A0192}\makebox[\dimexpr\linewidth-5em][s]{\hyperlink{paper:A0192}{\wuhao 仿生柔性气动低频水下声源及其亚波长发声机制}\dotfill\makebox[2.5em][r]{\wuhao\pageref{paper:A0192}}}
\par
\noindent\hspace*{5em}{\kaishu\wuhao 张金虎, 张宇}
\par\vspace{5pt}
\noindent\makebox[5em][l]{\wuhao A0219}\makebox[\dimexpr\linewidth-5em][s]{\hyperlink{paper:A0219}{\wuhao 基于复合型缺陷声子晶体的宽频高性能压电能量回收​}\dotfill\makebox[2.5em][r]{\wuhao\pageref{paper:A0219}}}
\par
\noindent\hspace*{5em}{\kaishu\wuhao 冯华, 李云霞}
\par\vspace{5pt}
\noindent\makebox[5em][l]{\wuhao A0222}\makebox[\dimexpr\linewidth-5em][s]{\hyperlink{paper:A0222}{\wuhao 振幅诱导的振荡圆柱尾流模态转变声散射特性研究}\dotfill\makebox[2.5em][r]{\wuhao\pageref{paper:A0222}}}
\par
\noindent\hspace*{5em}{\kaishu\wuhao 李松海, 高东宝, 贾根苏, 韩熙, 贺鹏举, 韩开峰}
\par\vspace{5pt}
\noindent\makebox[5em][l]{\wuhao A0223}\makebox[\dimexpr\linewidth-5em][s]{\hyperlink{paper:A0223}{\wuhao 散射奇异点的时空图像}\dotfill\makebox[2.5em][r]{\wuhao\pageref{paper:A0223}}}
\par
\noindent\hspace*{5em}{\kaishu\wuhao 肖钰}
\par\vspace{5pt}
\noindent\makebox[5em][l]{\wuhao A0228}\makebox[\dimexpr\linewidth-5em][s]{\hyperlink{paper:A0228}{\wuhao 基于高斯贝叶斯优化的声人工结构逆设计}\dotfill\makebox[2.5em][r]{\wuhao\pageref{paper:A0228}}}
\par
\noindent\hspace*{5em}{\kaishu\wuhao 曾云渝, 夏一飞, 陈安, 杨京, 梁彬, 程建春}
\par\vspace{5pt}
\noindent\makebox[5em][l]{\wuhao A0265}\makebox[\dimexpr\linewidth-5em][s]{\hyperlink{paper:A0265}{\wuhao 基于环形稀疏相控阵的多目标操控与组装}\dotfill\makebox[2.5em][r]{\wuhao\pageref{paper:A0265}}}
\par
\noindent\hspace*{5em}{\kaishu\wuhao 蒋智伟, 丁宁, 郭各朴, 马青玉}
\par\vspace{5pt}
\noindent\makebox[5em][l]{\wuhao A0277}\makebox[\dimexpr\linewidth-5em][s]{\hyperlink{paper:A0277}{\wuhao 基于遗传算法的通风隔声超材料设计}\dotfill\makebox[2.5em][r]{\wuhao\pageref{paper:A0277}}}
\par
\noindent\hspace*{5em}{\kaishu\wuhao 胡洁, 王标, 李澔翔}
\par\vspace{5pt}
\noindent\makebox[5em][l]{\wuhao A0280}\makebox[\dimexpr\linewidth-5em][s]{\hyperlink{paper:A0280}{\wuhao 基于有源非厄米声人工结构的奇异点调控}\dotfill\makebox[2.5em][r]{\wuhao\pageref{paper:A0280}}}
\par
\noindent\hspace*{5em}{\kaishu\wuhao 夏一飞, 陈安, 杨京, 梁彬, 程建春}
\par\vspace{5pt}
\noindent\makebox[5em][l]{\wuhao A0328}\makebox[\dimexpr\linewidth-5em][s]{\hyperlink{paper:A0328}{\wuhao 基于双面声学超表面的伪表面波调控研究}\dotfill\makebox[2.5em][r]{\wuhao\pageref{paper:A0328}}}
\par
\noindent\hspace*{5em}{\kaishu\wuhao 李澔翔}
\par\vspace{5pt}
\noindent\hangindent=5em\hangafter=1\makebox[5em][l]{\wuhao A0359}\hyperlink{paper:A0359}{\wuhao Gradient-equivalent medium enables acoustic rainbow capture and acoustic enhancement}
\par\noindent\hspace*{5em}\dotfill\makebox[2.5em][r]{\wuhao\pageref{paper:A0359}}
\par\vspace{2pt}
\noindent\hspace*{5em}{\kaishu\wuhao 任钰琳}
\par\vspace{5pt}
\noindent\makebox[5em][l]{\wuhao A0361}\makebox[\dimexpr\linewidth-5em][s]{\hyperlink{paper:A0361}{\wuhao 基于超声频率分集阵列的电子系统设计与聚焦特性研究}\dotfill\makebox[2.5em][r]{\wuhao\pageref{paper:A0361}}}
\par
\noindent\hspace*{5em}{\kaishu\wuhao 胡洁, 张珈铭, 李澔翔}
\par\vspace{5pt}
\noindent\makebox[5em][l]{\wuhao A0372}\makebox[\dimexpr\linewidth-5em][s]{\hyperlink{paper:A0372}{\wuhao 无对称性下的声学拓扑态研究}\dotfill\makebox[2.5em][r]{\wuhao\pageref{paper:A0372}}}
\par
\noindent\hspace*{5em}{\kaishu\wuhao 高峰}
\par\vspace{5pt}
\noindent\makebox[5em][l]{\wuhao A0394}\makebox[\dimexpr\linewidth-5em][s]{\hyperlink{paper:A0394}{\wuhao 基于超声超透镜的高精度成像及抗干扰性研究}\dotfill\makebox[2.5em][r]{\wuhao\pageref{paper:A0394}}}
\par
\noindent\hspace*{5em}{\kaishu\wuhao 江竣豪, 何佳杰, 陈霑}
\par\vspace{5pt}
\noindent\makebox[5em][l]{\wuhao A0411}\makebox[\dimexpr\linewidth-5em][s]{\hyperlink{paper:A0411}{\wuhao 基于动态可调式超表面的声场扩散特性实时调控}\dotfill\makebox[2.5em][r]{\wuhao\pageref{paper:A0411}}}
\par
\noindent\hspace*{5em}{\kaishu\wuhao 曹桐瑜, 刘京京, 梁彬, 程建春}
\par\vspace{5pt}
\noindent\makebox[5em][l]{\wuhao A0423}\makebox[\dimexpr\linewidth-5em][s]{\hyperlink{paper:A0423}{\wuhao 可重构声学伪表面波拓扑界面态}\dotfill\makebox[2.5em][r]{\wuhao\pageref{paper:A0423}}}
\par
\noindent\hspace*{5em}{\kaishu\wuhao 代若雪, 刘京京, 梁彬, 程建春}
\par\vspace{5pt}
\noindent\makebox[5em][l]{\wuhao A0463}\makebox[\dimexpr\linewidth-5em][s]{\hyperlink{paper:A0463}{\wuhao 锅炉管阵列中的声场分布及管壁声致剪切作用数值研究}\dotfill\makebox[2.5em][r]{\wuhao\pageref{paper:A0463}}}
\par
\noindent\hspace*{5em}{\kaishu\wuhao 高媛毓, 李颢, 周玉, 孙建浩, 姜根山}
\par\vspace{5pt}
\noindent\makebox[5em][l]{\wuhao A0464}\makebox[\dimexpr\linewidth-5em][s]{\hyperlink{paper:A0464}{\wuhao 声激励对扩散火焰燃烧响应与碳烟行为的影响}\dotfill\makebox[2.5em][r]{\wuhao\pageref{paper:A0464}}}
\par
\noindent\hspace*{5em}{\kaishu\wuhao 李颢, 高媛毓, 孙建浩, 周玉, 姜根山}
\par\vspace{5pt}
\noindent\makebox[5em][l]{\wuhao A0465}\makebox[\dimexpr\linewidth-5em][s]{\hyperlink{paper:A0465}{\wuhao 椭圆换热管外声流与温度场耦合机制研究}\dotfill\makebox[2.5em][r]{\wuhao\pageref{paper:A0465}}}
\par
\noindent\hspace*{5em}{\kaishu\wuhao 周玉, 孙建浩, 李颢, 高媛毓, 姜根山}
\par\vspace{5pt}
\noindent\makebox[5em][l]{\wuhao A0466}\makebox[\dimexpr\linewidth-5em][s]{\hyperlink{paper:A0466}{\wuhao 有源声学耦合谐振腔中的 PT 与反 PT 对称性}\dotfill\makebox[2.5em][r]{\wuhao\pageref{paper:A0466}}}
\par
\noindent\hspace*{5em}{\kaishu\wuhao 张名岳, 赵启睿, 杨玉真, 贾晗}
\par\vspace{5pt}
\vspace{6pt}

\begin{center}{\heiti\xiaosi B. 水声物理}\end{center}
\par\vspace{6pt}
\noindent\makebox[5em][l]{\wuhao B0017}\makebox[\dimexpr\linewidth-5em][s]{\hyperlink{paper:B0017}{\wuhao 深海海底混响波束-时延分布不对称性分析}\dotfill\makebox[2.5em][r]{\wuhao\pageref{paper:B0017}}}
\par
\noindent\hspace*{5em}{\kaishu\wuhao 孙一凡, 顾怡鸣, 秦继兴, 吴禹沈, 柳云峰, 张鹏, 殷凡}
\par\vspace{5pt}
\noindent\makebox[5em][l]{\wuhao B0021}\makebox[\dimexpr\linewidth-5em][s]{\hyperlink{paper:B0021}{\wuhao 阵列式微气泡声流体镊子操纵流体空间中的亚微米颗粒}\dotfill\makebox[2.5em][r]{\wuhao\pageref{paper:B0021}}}
\par
\noindent\hspace*{5em}{\kaishu\wuhao 王生庚}
\par\vspace{5pt}
\noindent\makebox[5em][l]{\wuhao B0024}\makebox[\dimexpr\linewidth-5em][s]{\hyperlink{paper:B0024}{\wuhao 水声传播波数积理论中深度格林函数的谱元解法}\dotfill\makebox[2.5em][r]{\wuhao\pageref{paper:B0024}}}
\par
\noindent\hspace*{5em}{\kaishu\wuhao 屠厚旺, 王勇献, 汪洋, 张运祥, 邓子豪}
\par\vspace{5pt}
\noindent\makebox[5em][l]{\wuhao B0025}\makebox[\dimexpr\linewidth-5em][s]{\hyperlink{paper:B0025}{\wuhao 基于物理约束稀疏贝叶斯学习的浅海爆炸声源宽带模态分离方法}\dotfill\makebox[2.5em][r]{\wuhao\pageref{paper:B0025}}}
\par
\noindent\hspace*{5em}{\kaishu\wuhao 张雪冬, 何琪, 宫在晓}
\par\vspace{5pt}
\noindent\makebox[5em][l]{\wuhao B0041}\makebox[\dimexpr\linewidth-5em][s]{\hyperlink{paper:B0041}{\wuhao 面向圆柱壳结构水下噪声截断模型数值预报方法研究}\dotfill\makebox[2.5em][r]{\wuhao\pageref{paper:B0041}}}
\par
\noindent\hspace*{5em}{\kaishu\wuhao 高南沙}
\par\vspace{5pt}
\noindent\makebox[5em][l]{\wuhao B0063}\makebox[\dimexpr\linewidth-5em][s]{\hyperlink{paper:B0063}{\wuhao 深海运动平台声场传播特性建模与分析}\dotfill\makebox[2.5em][r]{\wuhao\pageref{paper:B0063}}}
\par
\noindent\hspace*{5em}{\kaishu\wuhao 张祚祥, 吴金荣, 侯倩男, 蒋方冰, 晏子悦}
\par\vspace{5pt}
\noindent\makebox[5em][l]{\wuhao B0064}\makebox[\dimexpr\linewidth-5em][s]{\hyperlink{paper:B0064}{\wuhao 深海会聚区判别方法研究}\dotfill\makebox[2.5em][r]{\wuhao\pageref{paper:B0064}}}
\par
\noindent\hspace*{5em}{\kaishu\wuhao 任云, 张仁和}
\par\vspace{5pt}
\noindent\makebox[5em][l]{\wuhao B0068}\makebox[\dimexpr\linewidth-5em][s]{\hyperlink{paper:B0068}{\wuhao 南海内孤立波动态声场建模及其对水声探测的影响}\dotfill\makebox[2.5em][r]{\wuhao\pageref{paper:B0068}}}
\par
\noindent\hspace*{5em}{\kaishu\wuhao 荣林泰, 雷波}
\par\vspace{5pt}
\noindent\makebox[5em][l]{\wuhao B0069}\makebox[\dimexpr\linewidth-5em][s]{\hyperlink{paper:B0069}{\wuhao LFKAN-PINN ：融合KAN的物理信息神经网络用于水下声场重构}\dotfill\makebox[2.5em][r]{\wuhao\pageref{paper:B0069}}}
\par
\noindent\hspace*{5em}{\kaishu\wuhao 张帅帅, 赵文珺, 尉宣捷, 胡哲健, 郑广赢}
\par\vspace{5pt}
\noindent\makebox[5em][l]{\wuhao B0086}\makebox[\dimexpr\linewidth-5em][s]{\hyperlink{paper:B0086}{\wuhao 基于Welch谱估计的深海环境噪声数据处理参数优化与评估研究}\dotfill\makebox[2.5em][r]{\wuhao\pageref{paper:B0086}}}
\par
\noindent\hspace*{5em}{\kaishu\wuhao 张小艺, 郭新毅, 宋国丽}
\par\vspace{5pt}
\noindent\makebox[5em][l]{\wuhao B0100}\makebox[\dimexpr\linewidth-5em][s]{\hyperlink{paper:B0100}{\wuhao 海底沉积物剪切波速度和衰减对频率的依赖关系试验研究}\dotfill\makebox[2.5em][r]{\wuhao\pageref{paper:B0100}}}
\par
\noindent\hspace*{5em}{\kaishu\wuhao 张禹豪, 王景强, 李官保, 华清峰}
\par\vspace{5pt}
\noindent\makebox[5em][l]{\wuhao B0102}\makebox[\dimexpr\linewidth-5em][s]{\hyperlink{paper:B0102}{\wuhao 基于分位数回归的海洋环境噪声统计建模方法}\dotfill\makebox[2.5em][r]{\wuhao\pageref{paper:B0102}}}
\par
\noindent\hspace*{5em}{\kaishu\wuhao 宋国丽, 张乾初, 郭新毅, 张小艺, 马力}
\par\vspace{5pt}
\noindent\makebox[5em][l]{\wuhao B0109}\makebox[\dimexpr\linewidth-5em][s]{\hyperlink{paper:B0109}{\wuhao OpenOceanBellhop：声线并行加速与功能扩展实践}\dotfill\makebox[2.5em][r]{\wuhao\pageref{paper:B0109}}}
\par
\noindent\hspace*{5em}{\kaishu\wuhao 钱鹏, 梁羿, 谢世伟, 李整林, 肖鹏, 路一阳}
\par\vspace{5pt}
\noindent\makebox[5em][l]{\wuhao B0111}\makebox[\dimexpr\linewidth-5em][s]{\hyperlink{paper:B0111}{\wuhao 多普勒对声场干涉结构的影响研究}\dotfill\makebox[2.5em][r]{\wuhao\pageref{paper:B0111}}}
\par
\noindent\hspace*{5em}{\kaishu\wuhao 黄军, 李辉, 于亮, 史家辉}
\par\vspace{5pt}
\noindent\makebox[5em][l]{\wuhao B0125}\makebox[\dimexpr\linewidth-5em][s]{\hyperlink{paper:B0125}{\wuhao 基于附着框架气泡群的渔网声学特性增强研究}\dotfill\makebox[2.5em][r]{\wuhao\pageref{paper:B0125}}}
\par
\noindent\hspace*{5em}{\kaishu\wuhao 王越海钦, 程玄, 黄占东, 刘本奇, 龚志雄}
\par\vspace{5pt}
\noindent\makebox[5em][l]{\wuhao B0128}\makebox[\dimexpr\linewidth-5em][s]{\hyperlink{paper:B0128}{\wuhao 一种浅海锋面环境下的声源测距方法}\dotfill\makebox[2.5em][r]{\wuhao\pageref{paper:B0128}}}
\par
\noindent\hspace*{5em}{\kaishu\wuhao 孙勇, 秦继兴}
\par\vspace{5pt}
\noindent\makebox[5em][l]{\wuhao B0129}\makebox[\dimexpr\linewidth-5em][s]{\hyperlink{paper:B0129}{\wuhao 基于声线轨迹复用以及贡献插值的宽带声场快速计算方法}\dotfill\makebox[2.5em][r]{\wuhao\pageref{paper:B0129}}}
\par
\noindent\hspace*{5em}{\kaishu\wuhao 路一阳, 钱鹏, 李整林}
\par\vspace{5pt}
\noindent\makebox[5em][l]{\wuhao B0131}\makebox[\dimexpr\linewidth-5em][s]{\hyperlink{paper:B0131}{\wuhao 深海拖曳阵双基地主动声纳中结构化干扰实测分析}\dotfill\makebox[2.5em][r]{\wuhao\pageref{paper:B0131}}}
\par
\noindent\hspace*{5em}{\kaishu\wuhao 顾怡鸣, 宫在晓, 汪俊, 秦继兴, 李超}
\par\vspace{5pt}
\noindent\makebox[5em][l]{\wuhao B0143}\makebox[\dimexpr\linewidth-5em][s]{\hyperlink{paper:B0143}{\wuhao 水下航行体伴流场低频特性分析}\dotfill\makebox[2.5em][r]{\wuhao\pageref{paper:B0143}}}
\par
\noindent\hspace*{5em}{\kaishu\wuhao 张敬东, 李松, 李学智}
\par\vspace{5pt}
\noindent\makebox[5em][l]{\wuhao B0144}\makebox[\dimexpr\linewidth-5em][s]{\hyperlink{paper:B0144}{\wuhao 基于Bernoulli-Gaussian先验的联合多频离网格匹配场定位方法}\dotfill\makebox[2.5em][r]{\wuhao\pageref{paper:B0144}}}
\par
\noindent\hspace*{5em}{\kaishu\wuhao 李卿基, 韩笑}
\par\vspace{5pt}
\noindent\makebox[5em][l]{\wuhao B0167}\makebox[\dimexpr\linewidth-5em][s]{\hyperlink{paper:B0167}{\wuhao 耦合气泡的声信号研究}\dotfill\makebox[2.5em][r]{\wuhao\pageref{paper:B0167}}}
\par
\noindent\hspace*{5em}{\kaishu\wuhao 屈士琪, 青昕}
\par\vspace{5pt}
\noindent\makebox[5em][l]{\wuhao B0169}\makebox[\dimexpr\linewidth-5em][s]{\hyperlink{paper:B0169}{\wuhao 高阻尼改性非局域增强水下超薄低频宽带吸声超材料}\dotfill\makebox[2.5em][r]{\wuhao\pageref{paper:B0169}}}
\par
\noindent\hspace*{5em}{\kaishu\wuhao 李康乐, 范海燕, 朱一凡, 张辉}
\par\vspace{5pt}
\noindent\makebox[5em][l]{\wuhao B0199}\makebox[\dimexpr\linewidth-5em][s]{\hyperlink{paper:B0199}{\wuhao 基于分阵宽带相位特征提取的深海弱目标深度估计}\dotfill\makebox[2.5em][r]{\wuhao\pageref{paper:B0199}}}
\par
\noindent\hspace*{5em}{\kaishu\wuhao 杨霄宏, 尉宣捷, 郑广赢}
\par\vspace{5pt}
\noindent\makebox[5em][l]{\wuhao B0209}\makebox[\dimexpr\linewidth-5em][s]{\hyperlink{paper:B0209}{\wuhao 有限角域照明条件下深海界面路径族贡献与接收方位结构特性}\dotfill\makebox[2.5em][r]{\wuhao\pageref{paper:B0209}}}
\par
\noindent\hspace*{5em}{\kaishu\wuhao 陈颖璇, 李整林}
\par\vspace{5pt}
\noindent\makebox[5em][l]{\wuhao B0231}\makebox[\dimexpr\linewidth-5em][s]{\hyperlink{paper:B0231}{\wuhao 深海信道水声信号相关特性分析}\dotfill\makebox[2.5em][r]{\wuhao\pageref{paper:B0231}}}
\par
\noindent\hspace*{5em}{\kaishu\wuhao 徐传秀, 张虎, 徐硕硕, 谷雨}
\par\vspace{5pt}
\noindent\makebox[5em][l]{\wuhao B0239}\makebox[\dimexpr\linewidth-5em][s]{\hyperlink{paper:B0239}{\wuhao 深海连续海底条件下甚低频声传播特性研究}\dotfill\makebox[2.5em][r]{\wuhao\pageref{paper:B0239}}}
\par
\noindent\hspace*{5em}{\kaishu\wuhao 梁民帅, 吴涵雨, 江厚萱, 师俊杰, 孙大军}
\par\vspace{5pt}
\noindent\makebox[5em][l]{\wuhao B0275}\makebox[\dimexpr\linewidth-5em][s]{\hyperlink{paper:B0275}{\wuhao 基于相干信号与环境噪声相关性差异的测深方法性能分析}\dotfill\makebox[2.5em][r]{\wuhao\pageref{paper:B0275}}}
\par
\noindent\hspace*{5em}{\kaishu\wuhao 张文轩}
\par\vspace{5pt}
\noindent\makebox[5em][l]{\wuhao B0352}\makebox[\dimexpr\linewidth-5em][s]{\hyperlink{paper:B0352}{\wuhao 基于L-BFGS迭代修正的深海近场目标被动定位方法}\dotfill\makebox[2.5em][r]{\wuhao\pageref{paper:B0352}}}
\par
\noindent\hspace*{5em}{\kaishu\wuhao 侯森}
\par\vspace{5pt}
\noindent\makebox[5em][l]{\wuhao B0366}\makebox[\dimexpr\linewidth-5em][s]{\hyperlink{paper:B0366}{\wuhao 台风过境期间海洋环境噪声的时空频演化特性初步研究}\dotfill\makebox[2.5em][r]{\wuhao\pageref{paper:B0366}}}
\par
\noindent\hspace*{5em}{\kaishu\wuhao 王璟琰, 刘代}
\par\vspace{5pt}
\noindent\makebox[5em][l]{\wuhao B0379}\makebox[\dimexpr\linewidth-5em][s]{\hyperlink{paper:B0379}{\wuhao 起伏海面下基于射线理论的深海表面声道声传播建模}\dotfill\makebox[2.5em][r]{\wuhao\pageref{paper:B0379}}}
\par
\noindent\hspace*{5em}{\kaishu\wuhao 丁新宇, 段睿, 张虎}
\par\vspace{5pt}
\noindent\makebox[5em][l]{\wuhao B0389}\makebox[\dimexpr\linewidth-5em][s]{\hyperlink{paper:B0389}{\wuhao 融合抛物方程推进结构的逐步傅里叶神经算子声场预测方法}\dotfill\makebox[2.5em][r]{\wuhao\pageref{paper:B0389}}}
\par
\noindent\hspace*{5em}{\kaishu\wuhao 高宇翔, 肖鹏, 李整林}
\par\vspace{5pt}
\noindent\makebox[5em][l]{\wuhao B0420}\makebox[\dimexpr\linewidth-5em][s]{\hyperlink{paper:B0420}{\wuhao 内孤立波诱导次表层流场变化的组网式深海声层析观测}\dotfill\makebox[2.5em][r]{\wuhao\pageref{paper:B0420}}}
\par
\noindent\hspace*{5em}{\kaishu\wuhao 刘文凤, 王勇献}
\par\vspace{5pt}
\noindent\makebox[5em][l]{\wuhao B0443}\makebox[\dimexpr\linewidth-5em][s]{\hyperlink{paper:B0443}{\wuhao 水下航行器艏部声呐舱流噪声干扰识别方法}\dotfill\makebox[2.5em][r]{\wuhao\pageref{paper:B0443}}}
\par
\noindent\hspace*{5em}{\kaishu\wuhao 贾香路, 陈林烽}
\par\vspace{5pt}
\noindent\makebox[5em][l]{\wuhao B0448}\makebox[\dimexpr\linewidth-5em][s]{\hyperlink{paper:B0448}{\wuhao 深海目标多普勒特征分析}\dotfill\makebox[2.5em][r]{\wuhao\pageref{paper:B0448}}}
\par
\noindent\hspace*{5em}{\kaishu\wuhao 徐硕硕, 徐传秀, 张虎}
\par\vspace{5pt}
\noindent\makebox[5em][l]{\wuhao B0449}\makebox[\dimexpr\linewidth-5em][s]{\hyperlink{paper:B0449}{\wuhao 超大振幅内孤立波对声传播特性的影响研究}\dotfill\makebox[2.5em][r]{\wuhao\pageref{paper:B0449}}}
\par
\noindent\hspace*{5em}{\kaishu\wuhao 姜彦宇}
\par\vspace{5pt}
\vspace{6pt}

\begin{center}{\heiti\xiaosi C. 水声工程和水声信号处理}\end{center}
\par\vspace{6pt}
\noindent\makebox[5em][l]{\wuhao C0003}\makebox[\dimexpr\linewidth-5em][s]{\hyperlink{paper:C0003}{\wuhao 基于轻量级神经网络的远距离水声信号降噪技术研究}\dotfill\makebox[2.5em][r]{\wuhao\pageref{paper:C0003}}}
\par
\noindent\hspace*{5em}{\kaishu\wuhao 万芯源, 江伟华, 童峰}
\par\vspace{5pt}
\noindent\makebox[5em][l]{\wuhao C0014}\makebox[\dimexpr\linewidth-5em][s]{\hyperlink{paper:C0014}{\wuhao 深度学习方法在低频声传播场预测中的应用}\dotfill\makebox[2.5em][r]{\wuhao\pageref{paper:C0014}}}
\par
\noindent\hspace*{5em}{\kaishu\wuhao 冯肖}
\par\vspace{5pt}
\noindent\makebox[5em][l]{\wuhao C0020}\makebox[\dimexpr\linewidth-5em][s]{\hyperlink{paper:C0020}{\wuhao 基于涡旋相位调制的三维高分辨率水声成像}\dotfill\makebox[2.5em][r]{\wuhao\pageref{paper:C0020}}}
\par
\noindent\hspace*{5em}{\kaishu\wuhao 何先忠}
\par\vspace{5pt}
\noindent\makebox[5em][l]{\wuhao C0022}\makebox[\dimexpr\linewidth-5em][s]{\hyperlink{paper:C0022}{\wuhao 面向水下机动目标跟踪的水声网络高效传输协议}\dotfill\makebox[2.5em][r]{\wuhao\pageref{paper:C0022}}}
\par
\noindent\hspace*{5em}{\kaishu\wuhao 王超}
\par\vspace{5pt}
\noindent\makebox[5em][l]{\wuhao C0045}\makebox[\dimexpr\linewidth-5em][s]{\hyperlink{paper:C0045}{\wuhao 结合水平面阵实现深海声影区声源无模糊定位的方法}\dotfill\makebox[2.5em][r]{\wuhao\pageref{paper:C0045}}}
\par
\noindent\hspace*{5em}{\kaishu\wuhao 陈飞彤, 郭良浩, 陈德喜, 周雨飞, 牟学铭}
\par\vspace{5pt}
\noindent\makebox[5em][l]{\wuhao C0048}\makebox[\dimexpr\linewidth-5em][s]{\hyperlink{paper:C0048}{\wuhao 水下多目标特征辅助稳健纯方位跟踪算法研究}\dotfill\makebox[2.5em][r]{\wuhao\pageref{paper:C0048}}}
\par
\noindent\hspace*{5em}{\kaishu\wuhao 张淑娴, 唐胜雨, 侯翔昊}
\par\vspace{5pt}
\noindent\makebox[5em][l]{\wuhao C0049}\makebox[\dimexpr\linewidth-5em][s]{\hyperlink{paper:C0049}{\wuhao 南海主温跃层水声信道特性的深度依赖性研究}\dotfill\makebox[2.5em][r]{\wuhao\pageref{paper:C0049}}}
\par
\noindent\hspace*{5em}{\kaishu\wuhao 毕文旭, 陈鹏, 王威, 史阳}
\par\vspace{5pt}
\noindent\makebox[5em][l]{\wuhao C0051}\makebox[\dimexpr\linewidth-5em][s]{\hyperlink{paper:C0051}{\wuhao 基于希尔伯特-黄变换和数据增强的水声目标识别方法}\dotfill\makebox[2.5em][r]{\wuhao\pageref{paper:C0051}}}
\par
\noindent\hspace*{5em}{\kaishu\wuhao 姚琦海, 赵子杰, 汪勇, 杨益新}
\par\vspace{5pt}
\noindent\makebox[5em][l]{\wuhao C0052}\makebox[\dimexpr\linewidth-5em][s]{\hyperlink{paper:C0052}{\wuhao 复杂水声环境下抑制多途反射干扰的声源体积速度高精度测量方法}\dotfill\makebox[2.5em][r]{\wuhao\pageref{paper:C0052}}}
\par
\noindent\hspace*{5em}{\kaishu\wuhao 康鲁迪, 刘鹏, 刘碧龙}
\par\vspace{5pt}
\noindent\makebox[5em][l]{\wuhao C0055}\makebox[\dimexpr\linewidth-5em][s]{\hyperlink{paper:C0055}{\wuhao 基于 FxLMS 的矢量传感器振动噪声控制方法}\dotfill\makebox[2.5em][r]{\wuhao\pageref{paper:C0055}}}
\par
\noindent\hspace*{5em}{\kaishu\wuhao 邵睿涵, 温佳豪, 刘嘉文, 鹿一龙, 李智}
\par\vspace{5pt}
\noindent\makebox[5em][l]{\wuhao C0061}\makebox[\dimexpr\linewidth-5em][s]{\hyperlink{paper:C0061}{\wuhao 基于复合储能的水下脉冲声发射机设计方法}\dotfill\makebox[2.5em][r]{\wuhao\pageref{paper:C0061}}}
\par
\noindent\hspace*{5em}{\kaishu\wuhao 卢地华}
\par\vspace{5pt}
\noindent\makebox[5em][l]{\wuhao C0062}\makebox[\dimexpr\linewidth-5em][s]{\hyperlink{paper:C0062}{\wuhao 基于环境解耦与对比学习的跨环境水声目标识别方法}\dotfill\makebox[2.5em][r]{\wuhao\pageref{paper:C0062}}}
\par
\noindent\hspace*{5em}{\kaishu\wuhao 冉春晴, 刘一宁, 李整林}
\par\vspace{5pt}
\noindent\makebox[5em][l]{\wuhao C0078}\makebox[\dimexpr\linewidth-5em][s]{\hyperlink{paper:C0078}{\wuhao 水声通信定位一体化平台设计及试验}\dotfill\makebox[2.5em][r]{\wuhao\pageref{paper:C0078}}}
\par
\noindent\hspace*{5em}{\kaishu\wuhao 张腾飞}
\par\vspace{5pt}
\noindent\makebox[5em][l]{\wuhao C0079}\makebox[\dimexpr\linewidth-5em][s]{\hyperlink{paper:C0079}{\wuhao 不同发射波形下的主动声纳探测效能量化方法}\dotfill\makebox[2.5em][r]{\wuhao\pageref{paper:C0079}}}
\par
\noindent\hspace*{5em}{\kaishu\wuhao 赵猛, 钟何平, 唐劲松, 吴浩然}
\par\vspace{5pt}
\noindent\hangindent=5em\hangafter=1\makebox[5em][l]{\wuhao C0083}\hyperlink{paper:C0083}{\wuhao 基于Transformer与物理信息神经网络的水下机动目标二维联合定位}
\par\noindent\hspace*{5em}\dotfill\makebox[2.5em][r]{\wuhao\pageref{paper:C0083}}
\par\vspace{2pt}
\noindent\hspace*{5em}{\kaishu\wuhao 蒙子坤, 张文, 刘硕}
\par\vspace{5pt}
\noindent\hangindent=5em\hangafter=1\makebox[5em][l]{\wuhao C0084}\hyperlink{paper:C0084}{\wuhao 基于空间-通道双重注意力联合网络的主被动声呐特征级关联探测方法}
\par\noindent\hspace*{5em}\dotfill\makebox[2.5em][r]{\wuhao\pageref{paper:C0084}}
\par\vspace{2pt}
\noindent\hspace*{5em}{\kaishu\wuhao 宁玺翔, 李楠松}
\par\vspace{5pt}
\noindent\makebox[5em][l]{\wuhao C0089}\makebox[\dimexpr\linewidth-5em][s]{\hyperlink{paper:C0089}{\wuhao 面向水声信道的图像语义通信系统设计与实现}\dotfill\makebox[2.5em][r]{\wuhao\pageref{paper:C0089}}}
\par
\noindent\hspace*{5em}{\kaishu\wuhao 高鹏, 葛威, 张嘉恒, 陈迈, 贾亦真}
\par\vspace{5pt}
\noindent\makebox[5em][l]{\wuhao C0097}\makebox[\dimexpr\linewidth-5em][s]{\hyperlink{paper:C0097}{\wuhao 用于时间分集MIMO声呐成像的混合波束形成方法}\dotfill\makebox[2.5em][r]{\wuhao\pageref{paper:C0097}}}
\par
\noindent\hspace*{5em}{\kaishu\wuhao 赵泽恺, 刘雄厚, 杨益新}
\par\vspace{5pt}
\noindent\makebox[5em][l]{\wuhao C0098}\makebox[\dimexpr\linewidth-5em][s]{\hyperlink{paper:C0098}{\wuhao 深海多垂直阵列分布式特征融合的角度可分辨性退化与增益限制}\dotfill\makebox[2.5em][r]{\wuhao\pageref{paper:C0098}}}
\par
\noindent\hspace*{5em}{\kaishu\wuhao 刘鹏超, 郭新毅}
\par\vspace{5pt}
\noindent\makebox[5em][l]{\wuhao C0099}\makebox[\dimexpr\linewidth-5em][s]{\hyperlink{paper:C0099}{\wuhao 基于单矢量水听器的目标深度分类方法}\dotfill\makebox[2.5em][r]{\wuhao\pageref{paper:C0099}}}
\par
\noindent\hspace*{5em}{\kaishu\wuhao 何佩佩, 时洁, 夏峙}
\par\vspace{5pt}
\noindent\makebox[5em][l]{\wuhao C0106}\makebox[\dimexpr\linewidth-5em][s]{\hyperlink{paper:C0106}{\wuhao 基于稀疏贝叶斯学习的非均匀网格频率差分波达方向估计方法}\dotfill\makebox[2.5em][r]{\wuhao\pageref{paper:C0106}}}
\par
\noindent\hspace*{5em}{\kaishu\wuhao 刘晓燕, 牛海强, 袁泽, 王海斌}
\par\vspace{5pt}
\noindent\makebox[5em][l]{\wuhao C0107}\makebox[\dimexpr\linewidth-5em][s]{\hyperlink{paper:C0107}{\wuhao 一种基于数据库的被动目标特征匹配方法}\dotfill\makebox[2.5em][r]{\wuhao\pageref{paper:C0107}}}
\par
\noindent\hspace*{5em}{\kaishu\wuhao 程亚乔, 张虎}
\par\vspace{5pt}
\noindent\makebox[5em][l]{\wuhao C0116}\makebox[\dimexpr\linewidth-5em][s]{\hyperlink{paper:C0116}{\wuhao 基于加权双曲调频组合波形的水声感知通信一体化方法}\dotfill\makebox[2.5em][r]{\wuhao\pageref{paper:C0116}}}
\par
\noindent\hspace*{5em}{\kaishu\wuhao 许尚尚, 商志刚}
\par\vspace{5pt}
\noindent\makebox[5em][l]{\wuhao C0119}\makebox[\dimexpr\linewidth-5em][s]{\hyperlink{paper:C0119}{\wuhao 无人潜航器目标的主动声透射多维特征融合检测方法}\dotfill\makebox[2.5em][r]{\wuhao\pageref{paper:C0119}}}
\par
\noindent\hspace*{5em}{\kaishu\wuhao 何兆阳, 雷波}
\par\vspace{5pt}
\noindent\makebox[5em][l]{\wuhao C0122}\makebox[\dimexpr\linewidth-5em][s]{\hyperlink{paper:C0122}{\wuhao 基于压缩感知的拖船噪声稀疏抑制方法及实测验证}\dotfill\makebox[2.5em][r]{\wuhao\pageref{paper:C0122}}}
\par
\noindent\hspace*{5em}{\kaishu\wuhao 黄正超, 李整林, 肖鹏}
\par\vspace{5pt}
\noindent\makebox[5em][l]{\wuhao C0126}\makebox[\dimexpr\linewidth-5em][s]{\hyperlink{paper:C0126}{\wuhao 基于拖线阵的网络传输可靠性研究}\dotfill\makebox[2.5em][r]{\wuhao\pageref{paper:C0126}}}
\par
\noindent\hspace*{5em}{\kaishu\wuhao 王玥}
\par\vspace{5pt}
\noindent\makebox[5em][l]{\wuhao C0127}\makebox[\dimexpr\linewidth-5em][s]{\hyperlink{paper:C0127}{\wuhao 几何不确定性条件下AUV拖曳阵纯方位跟踪的动态误设下界研究}\dotfill\makebox[2.5em][r]{\wuhao\pageref{paper:C0127}}}
\par
\noindent\hspace*{5em}{\kaishu\wuhao 方子哲, 李整林, 刘涛}
\par\vspace{5pt}
\noindent\makebox[5em][l]{\wuhao C0132}\makebox[\dimexpr\linewidth-5em][s]{\hyperlink{paper:C0132}{\wuhao 辅助引导目标跟踪方法}\dotfill\makebox[2.5em][r]{\wuhao\pageref{paper:C0132}}}
\par
\noindent\hspace*{5em}{\kaishu\wuhao 安妍妍, 张虎}
\par\vspace{5pt}
\noindent\makebox[5em][l]{\wuhao C0133}\makebox[\dimexpr\linewidth-5em][s]{\hyperlink{paper:C0133}{\wuhao 联合多脉冲回波图像的主动目标运动轨迹可视化检测技术}\dotfill\makebox[2.5em][r]{\wuhao\pageref{paper:C0133}}}
\par
\noindent\hspace*{5em}{\kaishu\wuhao 徐彬, 邹丽娜, 张虎}
\par\vspace{5pt}
\noindent\makebox[5em][l]{\wuhao C0134}\makebox[\dimexpr\linewidth-5em][s]{\hyperlink{paper:C0134}{\wuhao 基于声场垂直波数滤波的水声目标深度估计方法}\dotfill\makebox[2.5em][r]{\wuhao\pageref{paper:C0134}}}
\par
\noindent\hspace*{5em}{\kaishu\wuhao 刘坤, 张虎}
\par\vspace{5pt}
\noindent\makebox[5em][l]{\wuhao C0136}\makebox[\dimexpr\linewidth-5em][s]{\hyperlink{paper:C0136}{\wuhao 基于傅里叶神经算子的失配环境声场预测与水下声源定位方法​}\dotfill\makebox[2.5em][r]{\wuhao\pageref{paper:C0136}}}
\par
\noindent\hspace*{5em}{\kaishu\wuhao 姚奇声, 刘一宁, 查若思, 肖鹏, 李整林}
\par\vspace{5pt}
\noindent\makebox[5em][l]{\wuhao C0142}\makebox[\dimexpr\linewidth-5em][s]{\hyperlink{paper:C0142}{\wuhao 一种考虑水声通信时延与丢包概率的被动协同跟踪方法}\dotfill\makebox[2.5em][r]{\wuhao\pageref{paper:C0142}}}
\par
\noindent\hspace*{5em}{\kaishu\wuhao 赵金虎, 冯西安}
\par\vspace{5pt}
\noindent\makebox[5em][l]{\wuhao C0145}\makebox[\dimexpr\linewidth-5em][s]{\hyperlink{paper:C0145}{\wuhao 基于特征辅助与自适应噪声估计的水下主动声呐多目标跟踪方法}\dotfill\makebox[2.5em][r]{\wuhao\pageref{paper:C0145}}}
\par
\noindent\hspace*{5em}{\kaishu\wuhao 张玉博, 王梦瑶, 张淑娴, 侯翔昊}
\par\vspace{5pt}
\noindent\makebox[5em][l]{\wuhao C0148}\makebox[\dimexpr\linewidth-5em][s]{\hyperlink{paper:C0148}{\wuhao 融合频域亚采样与黎曼距离度量的浅海主动声纳目标深度估计}\dotfill\makebox[2.5em][r]{\wuhao\pageref{paper:C0148}}}
\par
\noindent\hspace*{5em}{\kaishu\wuhao 郭小玮, 郑广赢, 杜栓平}
\par\vspace{5pt}
\noindent\makebox[5em][l]{\wuhao C0149}\makebox[\dimexpr\linewidth-5em][s]{\hyperlink{paper:C0149}{\wuhao 《面向南海复杂环境的声学通信机组网节点可靠性研究》}\dotfill\makebox[2.5em][r]{\wuhao\pageref{paper:C0149}}}
\par
\noindent\hspace*{5em}{\kaishu\wuhao 顾健}
\par\vspace{5pt}
\noindent\makebox[5em][l]{\wuhao C0154}\makebox[\dimexpr\linewidth-5em][s]{\hyperlink{paper:C0154}{\wuhao 深海低频宽带声源潜标设计与试验}\dotfill\makebox[2.5em][r]{\wuhao\pageref{paper:C0154}}}
\par
\noindent\hspace*{5em}{\kaishu\wuhao 陈琦, 郝浩琦, 马振, 高少波, 应晓伟, 郝运博, 徐杨帆}
\par\vspace{5pt}
\noindent\makebox[5em][l]{\wuhao C0159}\makebox[\dimexpr\linewidth-5em][s]{\hyperlink{paper:C0159}{\wuhao 复杂海洋波导声传播与水平横向相干性实验研究}\dotfill\makebox[2.5em][r]{\wuhao\pageref{paper:C0159}}}
\par
\noindent\hspace*{5em}{\kaishu\wuhao 曾启天, 李整林}
\par\vspace{5pt}
\noindent\makebox[5em][l]{\wuhao C0160}\makebox[\dimexpr\linewidth-5em][s]{\hyperlink{paper:C0160}{\wuhao 基于紧凑型最坏情况性能优化的鲁棒时域宽带波束形成方法}\dotfill\makebox[2.5em][r]{\wuhao\pageref{paper:C0160}}}
\par
\noindent\hspace*{5em}{\kaishu\wuhao 邓宇汐, 安良, 叶扬}
\par\vspace{5pt}
\noindent\hangindent=5em\hangafter=1\makebox[5em][l]{\wuhao C0161}\hyperlink{paper:C0161}{\wuhao 频移阵列数据：提升水声波束形成的波束分辨率及实现同方位脉冲声源分离}
\par\noindent\hspace*{5em}\dotfill\makebox[2.5em][r]{\wuhao\pageref{paper:C0161}}
\par\vspace{2pt}
\noindent\hspace*{5em}{\kaishu\wuhao 刘昌鹏, 周士弘, 戚聿波}
\par\vspace{5pt}
\noindent\makebox[5em][l]{\wuhao C0162}\makebox[\dimexpr\linewidth-5em][s]{\hyperlink{paper:C0162}{\wuhao 温度调制的水下小目标识别多模态算法}\dotfill\makebox[2.5em][r]{\wuhao\pageref{paper:C0162}}}
\par
\noindent\hspace*{5em}{\kaishu\wuhao 张博宇, 李书东, 彭旭名, 宋忠长, 张宇}
\par\vspace{5pt}
\noindent\makebox[5em][l]{\wuhao C0166}\makebox[\dimexpr\linewidth-5em][s]{\hyperlink{paper:C0166}{\wuhao 高速仿生隐蔽水声通信及其应用}\dotfill\makebox[2.5em][r]{\wuhao\pageref{paper:C0166}}}
\par
\noindent\hspace*{5em}{\kaishu\wuhao 黄子豪, 宋忠长, 张宇}
\par\vspace{5pt}
\noindent\makebox[5em][l]{\wuhao C0171}\makebox[\dimexpr\linewidth-5em][s]{\hyperlink{paper:C0171}{\wuhao 面向海洋环境噪声失配的水声目标识别方法研究}\dotfill\makebox[2.5em][r]{\wuhao\pageref{paper:C0171}}}
\par
\noindent\hspace*{5em}{\kaishu\wuhao 王培兵, 曾向阳}
\par\vspace{5pt}
\noindent\makebox[5em][l]{\wuhao C0184}\makebox[\dimexpr\linewidth-5em][s]{\hyperlink{paper:C0184}{\wuhao 基于量测复用的TPMBM的被动声纳目标跟踪方法}\dotfill\makebox[2.5em][r]{\wuhao\pageref{paper:C0184}}}
\par
\noindent\hspace*{5em}{\kaishu\wuhao 桂廉炜}
\par\vspace{5pt}
\noindent\makebox[5em][l]{\wuhao C0194}\makebox[\dimexpr\linewidth-5em][s]{\hyperlink{paper:C0194}{\wuhao 置信度驱动双环扩展卡尔曼滤波的水声通信波束跟踪方法}\dotfill\makebox[2.5em][r]{\wuhao\pageref{paper:C0194}}}
\par
\noindent\hspace*{5em}{\kaishu\wuhao 崔英桥}
\par\vspace{5pt}
\noindent\makebox[5em][l]{\wuhao C0195}\makebox[\dimexpr\linewidth-5em][s]{\hyperlink{paper:C0195}{\wuhao 基于传播状态先验的深海近底观测信息多距离适配回传方法}\dotfill\makebox[2.5em][r]{\wuhao\pageref{paper:C0195}}}
\par
\noindent\hspace*{5em}{\kaishu\wuhao 孙佳蓓, 雷波}
\par\vspace{5pt}
\noindent\makebox[5em][l]{\wuhao C0207}\makebox[\dimexpr\linewidth-5em][s]{\hyperlink{paper:C0207}{\wuhao 基于信道解耦合的目标径向尺度特征提取方法}\dotfill\makebox[2.5em][r]{\wuhao\pageref{paper:C0207}}}
\par
\noindent\hspace*{5em}{\kaishu\wuhao 王贺, 王栋梁, 谷新禹}
\par\vspace{5pt}
\noindent\makebox[5em][l]{\wuhao C0208}\makebox[\dimexpr\linewidth-5em][s]{\hyperlink{paper:C0208}{\wuhao 复杂海洋环境下基于环境噪声的声学层析}\dotfill\makebox[2.5em][r]{\wuhao\pageref{paper:C0208}}}
\par
\noindent\hspace*{5em}{\kaishu\wuhao 马天宇}
\par\vspace{5pt}
\noindent\makebox[5em][l]{\wuhao C0210}\makebox[\dimexpr\linewidth-5em][s]{\hyperlink{paper:C0210}{\wuhao 基于深度嵌入式多特征融合网络的水声目标识别方法}\dotfill\makebox[2.5em][r]{\wuhao\pageref{paper:C0210}}}
\par
\noindent\hspace*{5em}{\kaishu\wuhao 宋文涵, 郭子豪}
\par\vspace{5pt}
\noindent\makebox[5em][l]{\wuhao C0213}\makebox[\dimexpr\linewidth-5em][s]{\hyperlink{paper:C0213}{\wuhao 基于开源GPU架构的高频主动声纳检测功能组件化设计与实现}\dotfill\makebox[2.5em][r]{\wuhao\pageref{paper:C0213}}}
\par
\noindent\hspace*{5em}{\kaishu\wuhao 王志鹏, 雷波}
\par\vspace{5pt}
\noindent\makebox[5em][l]{\wuhao C0215}\makebox[\dimexpr\linewidth-5em][s]{\hyperlink{paper:C0215}{\wuhao 基于联合网络关联域适应的水下声学船舶识别方法}\dotfill\makebox[2.5em][r]{\wuhao\pageref{paper:C0215}}}
\par
\noindent\hspace*{5em}{\kaishu\wuhao 丁旺盼, 楼成淦, 黄清龙}
\par\vspace{5pt}
\noindent\makebox[5em][l]{\wuhao C0217}\makebox[\dimexpr\linewidth-5em][s]{\hyperlink{paper:C0217}{\wuhao 掠射角对深海水声通信信道中海面散射簇时延扩展的影响}\dotfill\makebox[2.5em][r]{\wuhao\pageref{paper:C0217}}}
\par
\noindent\hspace*{5em}{\kaishu\wuhao 田焕, 渠苏娜, 皇甫韬宇, 赵实}
\par\vspace{5pt}
\noindent\makebox[5em][l]{\wuhao C0218}\makebox[\dimexpr\linewidth-5em][s]{\hyperlink{paper:C0218}{\wuhao 基于复数卷积网络与块托普利兹先验的矢量阵列无网格DOA估计}\dotfill\makebox[2.5em][r]{\wuhao\pageref{paper:C0218}}}
\par
\noindent\hspace*{5em}{\kaishu\wuhao 林婧, 徐灵基, 陈鑫哲}
\par\vspace{5pt}
\noindent\makebox[5em][l]{\wuhao C0221}\makebox[\dimexpr\linewidth-5em][s]{\hyperlink{paper:C0221}{\wuhao 分数阶卷积公式在主动声呐探测中的应用}\dotfill\makebox[2.5em][r]{\wuhao\pageref{paper:C0221}}}
\par
\noindent\hspace*{5em}{\kaishu\wuhao 王成, 王方勇, 谷新禹}
\par\vspace{5pt}
\noindent\makebox[5em][l]{\wuhao C0225}\makebox[\dimexpr\linewidth-5em][s]{\hyperlink{paper:C0225}{\wuhao 基于扩展卡尔曼滤波的运动声源线谱增强方法}\dotfill\makebox[2.5em][r]{\wuhao\pageref{paper:C0225}}}
\par
\noindent\hspace*{5em}{\kaishu\wuhao 张羽飞, 赵梅}
\par\vspace{5pt}
\noindent\makebox[5em][l]{\wuhao C0229}\makebox[\dimexpr\linewidth-5em][s]{\hyperlink{paper:C0229}{\wuhao 深度学习的深海声影区单水听器俯仰角估计方法}\dotfill\makebox[2.5em][r]{\wuhao\pageref{paper:C0229}}}
\par
\noindent\hspace*{5em}{\kaishu\wuhao 周雨飞, 郭良浩, 章伟裕}
\par\vspace{5pt}
\noindent\makebox[5em][l]{\wuhao C0237}\makebox[\dimexpr\linewidth-5em][s]{\hyperlink{paper:C0237}{\wuhao 水下声学多芯线拖缆与对接缆对接悬挂方法}\dotfill\makebox[2.5em][r]{\wuhao\pageref{paper:C0237}}}
\par
\noindent\hspace*{5em}{\kaishu\wuhao 王晓飞, 张峰海, 胡建强, 陈玲, 江民振, 汪超}
\par\vspace{5pt}
\noindent\makebox[5em][l]{\wuhao C0242}\makebox[\dimexpr\linewidth-5em][s]{\hyperlink{paper:C0242}{\wuhao 融合量测来源判别与空关联建模的 Student-t 鲁棒目标跟踪方法}\dotfill\makebox[2.5em][r]{\wuhao\pageref{paper:C0242}}}
\par
\noindent\hspace*{5em}{\kaishu\wuhao 叶诗旗, 马晓川}
\par\vspace{5pt}
\noindent\makebox[5em][l]{\wuhao C0247}\makebox[\dimexpr\linewidth-5em][s]{\hyperlink{paper:C0247}{\wuhao 面向深远海动平台的声场快速预测与主被动声纳仿真方法}\dotfill\makebox[2.5em][r]{\wuhao\pageref{paper:C0247}}}
\par
\noindent\hspace*{5em}{\kaishu\wuhao 宋紫豪, 商志刚}
\par\vspace{5pt}
\noindent\makebox[5em][l]{\wuhao C0251}\makebox[\dimexpr\linewidth-5em][s]{\hyperlink{paper:C0251}{\wuhao 基于深度展开稀疏STAP的多基地主动声呐混响抑制方法研究}\dotfill\makebox[2.5em][r]{\wuhao\pageref{paper:C0251}}}
\par
\noindent\hspace*{5em}{\kaishu\wuhao 郑润帆, 李整林, 肖鹏}
\par\vspace{5pt}
\noindent\makebox[5em][l]{\wuhao C0254}\makebox[\dimexpr\linewidth-5em][s]{\hyperlink{paper:C0254}{\wuhao 发射系统幅频畸变下的舰船辐射噪声连续谱重构与谱级修正方法}\dotfill\makebox[2.5em][r]{\wuhao\pageref{paper:C0254}}}
\par
\noindent\hspace*{5em}{\kaishu\wuhao 李金, 陆志成, 马晓川}
\par\vspace{5pt}
\noindent\makebox[5em][l]{\wuhao C0255}\makebox[\dimexpr\linewidth-5em][s]{\hyperlink{paper:C0255}{\wuhao 基于稀疏贝叶斯估计的深海匹配场定位方法}\dotfill\makebox[2.5em][r]{\wuhao\pageref{paper:C0255}}}
\par
\noindent\hspace*{5em}{\kaishu\wuhao 胡佳华, 李整林, 肖鹏}
\par\vspace{5pt}
\noindent\makebox[5em][l]{\wuhao C0258}\makebox[\dimexpr\linewidth-5em][s]{\hyperlink{paper:C0258}{\wuhao 以机器学习精度为标尺的大型商船水下辐射噪声源级建模}\dotfill\makebox[2.5em][r]{\wuhao\pageref{paper:C0258}}}
\par
\noindent\hspace*{5em}{\kaishu\wuhao 黄拓, 江鹏飞, 张纹畅, 林建恒}
\par\vspace{5pt}
\noindent\makebox[5em][l]{\wuhao C0260}\makebox[\dimexpr\linewidth-5em][s]{\hyperlink{paper:C0260}{\wuhao 基于目标速度的单枚被动定向浮标定位跟踪}\dotfill\makebox[2.5em][r]{\wuhao\pageref{paper:C0260}}}
\par
\noindent\hspace*{5em}{\kaishu\wuhao 刘江涛, 乔斌, 张虎}
\par\vspace{5pt}
\noindent\makebox[5em][l]{\wuhao C0261}\makebox[\dimexpr\linewidth-5em][s]{\hyperlink{paper:C0261}{\wuhao 基于改进SELD的深海脉冲信号垂直到达角估计方法}\dotfill\makebox[2.5em][r]{\wuhao\pageref{paper:C0261}}}
\par
\noindent\hspace*{5em}{\kaishu\wuhao 盛宇婧, 安良, 蔡嘉宁, 丁子擎}
\par\vspace{5pt}
\noindent\makebox[5em][l]{\wuhao C0262}\makebox[\dimexpr\linewidth-5em][s]{\hyperlink{paper:C0262}{\wuhao 基于特征空间稳健Capon与R-L解卷积的水声高分辨空间谱估计}\dotfill\makebox[2.5em][r]{\wuhao\pageref{paper:C0262}}}
\par
\noindent\hspace*{5em}{\kaishu\wuhao 梁硕, 张虎}
\par\vspace{5pt}
\noindent\makebox[5em][l]{\wuhao C0263}\makebox[\dimexpr\linewidth-5em][s]{\hyperlink{paper:C0263}{\wuhao 基于超小波的矢量传感器脉冲干扰超分辨率检测}\dotfill\makebox[2.5em][r]{\wuhao\pageref{paper:C0263}}}
\par
\noindent\hspace*{5em}{\kaishu\wuhao 聂鹏龙, 张虎, 乔斌}
\par\vspace{5pt}
\noindent\makebox[5em][l]{\wuhao C0266}\makebox[\dimexpr\linewidth-5em][s]{\hyperlink{paper:C0266}{\wuhao 基于二维混合加权的密排双线阵空间谱估计方法研究}\dotfill\makebox[2.5em][r]{\wuhao\pageref{paper:C0266}}}
\par
\noindent\hspace*{5em}{\kaishu\wuhao 杨威, 武建国, 聂鹏龙}
\par\vspace{5pt}
\noindent\makebox[5em][l]{\wuhao C0267}\makebox[\dimexpr\linewidth-5em][s]{\hyperlink{paper:C0267}{\wuhao 面向水下无人航行器辐射噪声的门控谱及空间联合检测方法}\dotfill\makebox[2.5em][r]{\wuhao\pageref{paper:C0267}}}
\par
\noindent\hspace*{5em}{\kaishu\wuhao 周煌卿, 张亚静, 史阳}
\par\vspace{5pt}
\noindent\makebox[5em][l]{\wuhao C0268}\makebox[\dimexpr\linewidth-5em][s]{\hyperlink{paper:C0268}{\wuhao 基于加权邻域相对熵的浅海主动声呐弱目标检测算法}\dotfill\makebox[2.5em][r]{\wuhao\pageref{paper:C0268}}}
\par
\noindent\hspace*{5em}{\kaishu\wuhao 程垦, 程嘉玺, 王方勇}
\par\vspace{5pt}
\noindent\makebox[5em][l]{\wuhao C0269}\makebox[\dimexpr\linewidth-5em][s]{\hyperlink{paper:C0269}{\wuhao 基于深海拖线阵声波垂直到达角补偿的测向误差修正方法研究}\dotfill\makebox[2.5em][r]{\wuhao\pageref{paper:C0269}}}
\par
\noindent\hspace*{5em}{\kaishu\wuhao 武建国, 谷雨, 张虎}
\par\vspace{5pt}
\noindent\makebox[5em][l]{\wuhao C0270}\makebox[\dimexpr\linewidth-5em][s]{\hyperlink{paper:C0270}{\wuhao 基于全相位卷积窗的水声单频信号能量泄漏抑制方法}\dotfill\makebox[2.5em][r]{\wuhao\pageref{paper:C0270}}}
\par
\noindent\hspace*{5em}{\kaishu\wuhao 程嘉玺, 王方勇, 程垦}
\par\vspace{5pt}
\noindent\makebox[5em][l]{\wuhao C0284}\makebox[\dimexpr\linewidth-5em][s]{\hyperlink{paper:C0284}{\wuhao 基于模糊函数域形态学重建的WVD交叉项抑制方法}\dotfill\makebox[2.5em][r]{\wuhao\pageref{paper:C0284}}}
\par
\noindent\hspace*{5em}{\kaishu\wuhao 胡选富}
\par\vspace{5pt}
\noindent\makebox[5em][l]{\wuhao C0291}\makebox[\dimexpr\linewidth-5em][s]{\hyperlink{paper:C0291}{\wuhao 面向类别不均衡场景的原型协同适配小样本水声目标识别方法}\dotfill\makebox[2.5em][r]{\wuhao\pageref{paper:C0291}}}
\par
\noindent\hspace*{5em}{\kaishu\wuhao 黄骏龙, 徐灵基, 李整林}
\par\vspace{5pt}
\noindent\makebox[5em][l]{\wuhao C0294}\makebox[\dimexpr\linewidth-5em][s]{\hyperlink{paper:C0294}{\wuhao 基于声场干涉结构的深海目标深度估计及嵌入式实现方法}\dotfill\makebox[2.5em][r]{\wuhao\pageref{paper:C0294}}}
\par
\noindent\hspace*{5em}{\kaishu\wuhao 陈弘淇, 雷波}
\par\vspace{5pt}
\noindent\makebox[5em][l]{\wuhao C0297}\makebox[\dimexpr\linewidth-5em][s]{\hyperlink{paper:C0297}{\wuhao 基于非均匀搜索策略的大规模阵列近场源离格定位方法}\dotfill\makebox[2.5em][r]{\wuhao\pageref{paper:C0297}}}
\par
\noindent\hspace*{5em}{\kaishu\wuhao 王域, 黄健丽, 宫在晓, 汪俊, 牛海强}
\par\vspace{5pt}
\noindent\makebox[5em][l]{\wuhao C0303}\makebox[\dimexpr\linewidth-5em][s]{\hyperlink{paper:C0303}{\wuhao 柔性拖线阵三维阵形稳健重构算法}\dotfill\makebox[2.5em][r]{\wuhao\pageref{paper:C0303}}}
\par
\noindent\hspace*{5em}{\kaishu\wuhao 王冬, 宫在晓, 王域, 汪俊}
\par\vspace{5pt}
\noindent\makebox[5em][l]{\wuhao C0308}\makebox[\dimexpr\linewidth-5em][s]{\hyperlink{paper:C0308}{\wuhao 面向OFDM体制的水声通信定位一体化技术研究}\dotfill\makebox[2.5em][r]{\wuhao\pageref{paper:C0308}}}
\par
\noindent\hspace*{5em}{\kaishu\wuhao 曲星昊, 乔钢, 商志刚, 许尚尚}
\par\vspace{5pt}
\noindent\makebox[5em][l]{\wuhao C0311}\makebox[\dimexpr\linewidth-5em][s]{\hyperlink{paper:C0311}{\wuhao 基于深度学习的短时间历程深海声源深度估计方法研究}\dotfill\makebox[2.5em][r]{\wuhao\pageref{paper:C0311}}}
\par
\noindent\hspace*{5em}{\kaishu\wuhao 梁民帅, 白一杉, 孙大军, 师俊杰, 张晓妍}
\par\vspace{5pt}
\noindent\makebox[5em][l]{\wuhao C0313}\makebox[\dimexpr\linewidth-5em][s]{\hyperlink{paper:C0313}{\wuhao 基于声场高阶量的小型化阵列高分辨方位估计算法研究}\dotfill\makebox[2.5em][r]{\wuhao\pageref{paper:C0313}}}
\par
\noindent\hspace*{5em}{\kaishu\wuhao 吴迪, 王世博, 王希晨, 乔小瑞}
\par\vspace{5pt}
\noindent\makebox[5em][l]{\wuhao C0315}\makebox[\dimexpr\linewidth-5em][s]{\hyperlink{paper:C0315}{\wuhao 面向被动探测场景的海洋环境参数时空分辨率阈值分析方法}\dotfill\makebox[2.5em][r]{\wuhao\pageref{paper:C0315}}}
\par
\noindent\hspace*{5em}{\kaishu\wuhao 谢世伟, 钱鹏, 肖鹏, 李整林}
\par\vspace{5pt}
\noindent\makebox[5em][l]{\wuhao C0316}\makebox[\dimexpr\linewidth-5em][s]{\hyperlink{paper:C0316}{\wuhao 一种改进型灰度关联算法}\dotfill\makebox[2.5em][r]{\wuhao\pageref{paper:C0316}}}
\par
\noindent\hspace*{5em}{\kaishu\wuhao 申屠云奇, 张虎}
\par\vspace{5pt}
\noindent\makebox[5em][l]{\wuhao C0317}\makebox[\dimexpr\linewidth-5em][s]{\hyperlink{paper:C0317}{\wuhao 异步多用户水声通信的多址干扰抑制技术研究}\dotfill\makebox[2.5em][r]{\wuhao\pageref{paper:C0317}}}
\par
\noindent\hspace*{5em}{\kaishu\wuhao 熊省军, 李栋, 冯玮, 朱小辉}
\par\vspace{5pt}
\noindent\makebox[5em][l]{\wuhao C0319}\makebox[\dimexpr\linewidth-5em][s]{\hyperlink{paper:C0319}{\wuhao 波浪滑翔艇场景下水声分层广播通信研究}\dotfill\makebox[2.5em][r]{\wuhao\pageref{paper:C0319}}}
\par
\noindent\hspace*{5em}{\kaishu\wuhao 全威威, 商志刚}
\par\vspace{5pt}
\noindent\makebox[5em][l]{\wuhao C0326}\makebox[\dimexpr\linewidth-5em][s]{\hyperlink{paper:C0326}{\wuhao 一种基于连续谱峰值特征的水声目标分类识别方法}\dotfill\makebox[2.5em][r]{\wuhao\pageref{paper:C0326}}}
\par
\noindent\hspace*{5em}{\kaishu\wuhao 李晓彬, 叶思懋, 谷新禹}
\par\vspace{5pt}
\noindent\makebox[5em][l]{\wuhao C0334}\makebox[\dimexpr\linewidth-5em][s]{\hyperlink{paper:C0334}{\wuhao 基于人机协同学习的少样本水声目标识别方法}\dotfill\makebox[2.5em][r]{\wuhao\pageref{paper:C0334}}}
\par
\noindent\hspace*{5em}{\kaishu\wuhao 朱梓硕, 潘湘林, 李忠杰, 赵安康, 赵亮, 张博研, 计博源, 甘霖}
\par\vspace{5pt}
\noindent\makebox[5em][l]{\wuhao C0337}\makebox[\dimexpr\linewidth-5em][s]{\hyperlink{paper:C0337}{\wuhao 基于高阶谱特征的水声目标辐射噪声线谱检测方法研究}\dotfill\makebox[2.5em][r]{\wuhao\pageref{paper:C0337}}}
\par
\noindent\hspace*{5em}{\kaishu\wuhao 宋楠楠, 梅贝宁, 邓宇汐, 范青, 朱启轩, 安良}
\par\vspace{5pt}
\noindent\makebox[5em][l]{\wuhao C0339}\makebox[\dimexpr\linewidth-5em][s]{\hyperlink{paper:C0339}{\wuhao 北极双声道波导宽带频散特征提取与无源测距方法研究}\dotfill\makebox[2.5em][r]{\wuhao\pageref{paper:C0339}}}
\par
\noindent\hspace*{5em}{\kaishu\wuhao 胡思源, 徐灵基, 文洪涛, 刘思健}
\par\vspace{5pt}
\noindent\hangindent=5em\hangafter=1\makebox[5em][l]{\wuhao C0344}\hyperlink{paper:C0344}{\wuhao 基于短距离主动源激发的冰层Quasi-Scholte模态频散曲线提取方法}
\par\noindent\hspace*{5em}\dotfill\makebox[2.5em][r]{\wuhao\pageref{paper:C0344}}
\par\vspace{2pt}
\noindent\hspace*{5em}{\kaishu\wuhao 杨锐, 马丁一, 张宇翔}
\par\vspace{5pt}
\noindent\makebox[5em][l]{\wuhao C0348}\makebox[\dimexpr\linewidth-5em][s]{\hyperlink{paper:C0348}{\wuhao 基于海底声纳图像的底质分类方法研究}\dotfill\makebox[2.5em][r]{\wuhao\pageref{paper:C0348}}}
\par
\noindent\hspace*{5em}{\kaishu\wuhao 赵廷, 陈明, 朱可心}
\par\vspace{5pt}
\noindent\makebox[5em][l]{\wuhao C0351}\makebox[\dimexpr\linewidth-5em][s]{\hyperlink{paper:C0351}{\wuhao 面相中心式组网的水声通信网络协议设计}\dotfill\makebox[2.5em][r]{\wuhao\pageref{paper:C0351}}}
\par
\noindent\hspace*{5em}{\kaishu\wuhao 王桢铎, 温梦华, 王超, 谢哲}
\par\vspace{5pt}
\noindent\makebox[5em][l]{\wuhao C0353}\makebox[\dimexpr\linewidth-5em][s]{\hyperlink{paper:C0353}{\wuhao 基于稀疏重构处理的协方差矩阵拟合性能分析}\dotfill\makebox[2.5em][r]{\wuhao\pageref{paper:C0353}}}
\par
\noindent\hspace*{5em}{\kaishu\wuhao 王宣, 赵欧南}
\par\vspace{5pt}
\noindent\makebox[5em][l]{\wuhao C0358}\makebox[\dimexpr\linewidth-5em][s]{\hyperlink{paper:C0358}{\wuhao 多源异构水声信号的大数据融合与自适应降噪算法}\dotfill\makebox[2.5em][r]{\wuhao\pageref{paper:C0358}}}
\par
\noindent\hspace*{5em}{\kaishu\wuhao 刘璐}
\par\vspace{5pt}
\noindent\makebox[5em][l]{\wuhao C0360}\makebox[\dimexpr\linewidth-5em][s]{\hyperlink{paper:C0360}{\wuhao 基于涡旋声呐的水下掩埋目标探测研究}\dotfill\makebox[2.5em][r]{\wuhao\pageref{paper:C0360}}}
\par
\noindent\hspace*{5em}{\kaishu\wuhao 陈利伟, 范军, 龚志雄}
\par\vspace{5pt}
\noindent\makebox[5em][l]{\wuhao C0363}\makebox[\dimexpr\linewidth-5em][s]{\hyperlink{paper:C0363}{\wuhao 面向被动声纳的降维快速变分稀疏贝叶斯测向框架}\dotfill\makebox[2.5em][r]{\wuhao\pageref{paper:C0363}}}
\par
\noindent\hspace*{5em}{\kaishu\wuhao 赵欧南, 王宣}
\par\vspace{5pt}
\noindent\makebox[5em][l]{\wuhao C0365}\makebox[\dimexpr\linewidth-5em][s]{\hyperlink{paper:C0365}{\wuhao 结合频谱特征的波束历程图多目标分辨与跟踪​}\dotfill\makebox[2.5em][r]{\wuhao\pageref{paper:C0365}}}
\par
\noindent\hspace*{5em}{\kaishu\wuhao 殷凡, 宫在晓, 顾怡鸣, 王域, 汪俊}
\par\vspace{5pt}
\noindent\makebox[5em][l]{\wuhao C0373}\makebox[\dimexpr\linewidth-5em][s]{\hyperlink{paper:C0373}{\wuhao DQN-TED:面向水声传感器网络的自适应早停高效拓扑发现算法}\dotfill\makebox[2.5em][r]{\wuhao\pageref{paper:C0373}}}
\par
\noindent\hspace*{5em}{\kaishu\wuhao 赵倩, 孙瑶, 葛威}
\par\vspace{5pt}
\noindent\makebox[5em][l]{\wuhao C0374}\makebox[\dimexpr\linewidth-5em][s]{\hyperlink{paper:C0374}{\wuhao 基于人工特征与语义特征空间联合增强的被动水声多目标识别方法}\dotfill\makebox[2.5em][r]{\wuhao\pageref{paper:C0374}}}
\par
\noindent\hspace*{5em}{\kaishu\wuhao 肖自远, 韩一娜, 刘昱韬}
\par\vspace{5pt}
\noindent\makebox[5em][l]{\wuhao C0377}\makebox[\dimexpr\linewidth-5em][s]{\hyperlink{paper:C0377}{\wuhao 基于波形扩散模型的水声目标微弱动态频率线恢复方法}\dotfill\makebox[2.5em][r]{\wuhao\pageref{paper:C0377}}}
\par
\noindent\hspace*{5em}{\kaishu\wuhao 肖自远, 韩一娜, 刘昱韬}
\par\vspace{5pt}
\noindent\makebox[5em][l]{\wuhao C0388}\makebox[\dimexpr\linewidth-5em][s]{\hyperlink{paper:C0388}{\wuhao 融合声学物理特征的船舶目标识别梅尔谱优化方法}\dotfill\makebox[2.5em][r]{\wuhao\pageref{paper:C0388}}}
\par
\noindent\hspace*{5em}{\kaishu\wuhao 李虹希, 郎子乂, 朱宏娜}
\par\vspace{5pt}
\noindent\hangindent=5em\hangafter=1\makebox[5em][l]{\wuhao C0398}\hyperlink{paper:C0398}{\wuhao 基于多约束联合筛选的四波束 ADCP 非等间隔 CRT 速度解模糊方法}
\par\noindent\hspace*{5em}\dotfill\makebox[2.5em][r]{\wuhao\pageref{paper:C0398}}
\par\vspace{2pt}
\noindent\hspace*{5em}{\kaishu\wuhao 曹文婧, 杨永寿, 胡长和, 程何小}
\par\vspace{5pt}
\noindent\makebox[5em][l]{\wuhao C0399}\makebox[\dimexpr\linewidth-5em][s]{\hyperlink{paper:C0399}{\wuhao 基于深度强化学习的时延容忍水声通信网络路由协议研究}\dotfill\makebox[2.5em][r]{\wuhao\pageref{paper:C0399}}}
\par
\noindent\hspace*{5em}{\kaishu\wuhao 葛晨昊, 王晓林, 罗栋栋}
\par\vspace{5pt}
\noindent\makebox[5em][l]{\wuhao C0403}\makebox[\dimexpr\linewidth-5em][s]{\hyperlink{paper:C0403}{\wuhao 利用多线谱相位信息的运动目标径向速度估计方法}\dotfill\makebox[2.5em][r]{\wuhao\pageref{paper:C0403}}}
\par
\noindent\hspace*{5em}{\kaishu\wuhao 张柳, 王方勇, 陈孝森, 修明明}
\par\vspace{5pt}
\noindent\makebox[5em][l]{\wuhao C0419}\makebox[\dimexpr\linewidth-5em][s]{\hyperlink{paper:C0419}{\wuhao 基于节点加权EKF的双滑翔机分布式被动定位方法}\dotfill\makebox[2.5em][r]{\wuhao\pageref{paper:C0419}}}
\par
\noindent\hspace*{5em}{\kaishu\wuhao 周子豪, 史阳, 许晨天, 刘孟龙}
\par\vspace{5pt}
\noindent\makebox[5em][l]{\wuhao C0427}\makebox[\dimexpr\linewidth-5em][s]{\hyperlink{paper:C0427}{\wuhao 利用多途匹配的深海瞬态声源深度属性分类方法}\dotfill\makebox[2.5em][r]{\wuhao\pageref{paper:C0427}}}
\par
\noindent\hspace*{5em}{\kaishu\wuhao 耿彬, 李恒光, 范泽亚}
\par\vspace{5pt}
\noindent\hangindent=5em\hangafter=1\makebox[5em][l]{\wuhao C0430}\hyperlink{paper:C0430}{\wuhao 强混响下基于宽带仿射一致性与距离偏差自校正的高速水下目标检测方法}
\par\noindent\hspace*{5em}\dotfill\makebox[2.5em][r]{\wuhao\pageref{paper:C0430}}
\par\vspace{2pt}
\noindent\hspace*{5em}{\kaishu\wuhao 李志欣}
\par\vspace{5pt}
\noindent\makebox[5em][l]{\wuhao C0454}\makebox[\dimexpr\linewidth-5em][s]{\hyperlink{paper:C0454}{\wuhao 基于对比学习和多维度量融合的跨域开集水声目标识别方法}\dotfill\makebox[2.5em][r]{\wuhao\pageref{paper:C0454}}}
\par
\noindent\hspace*{5em}{\kaishu\wuhao 黎志坚, 刘一宁, 黄骏龙}
\par\vspace{5pt}
\noindent\makebox[5em][l]{\wuhao C0455}\makebox[\dimexpr\linewidth-5em][s]{\hyperlink{paper:C0455}{\wuhao 基于TDOA-DOA观测信息融合的水下目标稳健分布式定位方法}\dotfill\makebox[2.5em][r]{\wuhao\pageref{paper:C0455}}}
\par
\noindent\hspace*{5em}{\kaishu\wuhao 梅贝宁, 安良, 朱启轩}
\par\vspace{5pt}
\vspace{6pt}

\begin{center}{\heiti\xiaosi D. 功率超声}\end{center}
\par\vspace{6pt}
\noindent\makebox[5em][l]{\wuhao D0018}\makebox[\dimexpr\linewidth-5em][s]{\hyperlink{paper:D0018}{\wuhao 时间延迟对两个相互作用气泡运动的影响}\dotfill\makebox[2.5em][r]{\wuhao\pageref{paper:D0018}}}
\par
\noindent\hspace*{5em}{\kaishu\wuhao 张玲玲}
\par\vspace{5pt}
\noindent\makebox[5em][l]{\wuhao D0043}\makebox[\dimexpr\linewidth-5em][s]{\hyperlink{paper:D0043}{\wuhao 工业超声清洗机的空化控制研究}\dotfill\makebox[2.5em][r]{\wuhao\pageref{paper:D0043}}}
\par
\noindent\hspace*{5em}{\kaishu\wuhao 刘恩泽}
\par\vspace{5pt}
\noindent\makebox[5em][l]{\wuhao D0080}\makebox[\dimexpr\linewidth-5em][s]{\hyperlink{paper:D0080}{\wuhao 基于声黑洞与模式转换设计的杯型超声辐射器研究}\dotfill\makebox[2.5em][r]{\wuhao\pageref{paper:D0080}}}
\par
\noindent\hspace*{5em}{\kaishu\wuhao 陈诚, 董宜雷, 徐春龙, 郭建中, 林书玉}
\par\vspace{5pt}
\noindent\makebox[5em][l]{\wuhao D0152}\makebox[\dimexpr\linewidth-5em][s]{\hyperlink{paper:D0152}{\wuhao 超声凝聚强化颗粒物控制：从实验室到工业应用}\dotfill\makebox[2.5em][r]{\wuhao\pageref{paper:D0152}}}
\par
\noindent\hspace*{5em}{\kaishu\wuhao 刘鹏展, Ng Bing Feng}
\par\vspace{5pt}
\noindent\makebox[5em][l]{\wuhao D0197}\makebox[\dimexpr\linewidth-5em][s]{\hyperlink{paper:D0197}{\wuhao 近刚性壁面半椭球形空化泡团动力学演化行为研究}\dotfill\makebox[2.5em][r]{\wuhao\pageref{paper:D0197}}}
\par
\noindent\hspace*{5em}{\kaishu\wuhao 郭霞杰, 吴玉婷, 雷照康, 李秀如, 崔致远, 王成会}
\par\vspace{5pt}
\noindent\makebox[5em][l]{\wuhao D0224}\makebox[\dimexpr\linewidth-5em][s]{\hyperlink{paper:D0224}{\wuhao 毛细管中空气柱存在调节气泡模态振动分布研究}\dotfill\makebox[2.5em][r]{\wuhao\pageref{paper:D0224}}}
\par
\noindent\hspace*{5em}{\kaishu\wuhao 吴玉婷, 李秀如, 冯康艺, 王成会}
\par\vspace{5pt}
\noindent\makebox[5em][l]{\wuhao D0232}\makebox[\dimexpr\linewidth-5em][s]{\hyperlink{paper:D0232}{\wuhao 近共振气泡与糖颗粒的动态自组装及持续驱动}\dotfill\makebox[2.5em][r]{\wuhao\pageref{paper:D0232}}}
\par
\noindent\hspace*{5em}{\kaishu\wuhao 马佳昱, 王成会}
\par\vspace{5pt}
\noindent\makebox[5em][l]{\wuhao D0287}\makebox[\dimexpr\linewidth-5em][s]{\hyperlink{paper:D0287}{\wuhao 磁致伸缩换能器多层线圈中的高频铜损}\dotfill\makebox[2.5em][r]{\wuhao\pageref{paper:D0287}}}
\par
\noindent\hspace*{5em}{\kaishu\wuhao 周磊, 贺西平, 陈恬雅, 廖迈伟}
\par\vspace{5pt}
\noindent\makebox[5em][l]{\wuhao D0288}\makebox[\dimexpr\linewidth-5em][s]{\hyperlink{paper:D0288}{\wuhao Terfenol-D环的磁导率和电磁能损耗研究}\dotfill\makebox[2.5em][r]{\wuhao\pageref{paper:D0288}}}
\par
\noindent\hspace*{5em}{\kaishu\wuhao 陈恬雅, 贺西平, 周磊, 廖迈伟}
\par\vspace{5pt}
\noindent\makebox[5em][l]{\wuhao D0314}\makebox[\dimexpr\linewidth-5em][s]{\hyperlink{paper:D0314}{\wuhao 超声波滚焊特性分析与研究}\dotfill\makebox[2.5em][r]{\wuhao\pageref{paper:D0314}}}
\par
\noindent\hspace*{5em}{\kaishu\wuhao 姚震, 林超, 符涵, 刘云龙}
\par\vspace{5pt}
\noindent\makebox[5em][l]{\wuhao D0347}\makebox[\dimexpr\linewidth-5em][s]{\hyperlink{paper:D0347}{\wuhao 氟硅烷改性聚醚型聚氨酯水声透声材料的制备和性能分析}\dotfill\makebox[2.5em][r]{\wuhao\pageref{paper:D0347}}}
\par
\noindent\hspace*{5em}{\kaishu\wuhao 林海景, 沈壮志}
\par\vspace{5pt}
\noindent\makebox[5em][l]{\wuhao D0376}\makebox[\dimexpr\linewidth-5em][s]{\hyperlink{paper:D0376}{\wuhao 超声诱导液相剥离制备高性能碲化铋基热电薄膜}\dotfill\makebox[2.5em][r]{\wuhao\pageref{paper:D0376}}}
\par
\noindent\hspace*{5em}{\kaishu\wuhao 李思慧, 邵逸之, 孙可, 陈丹, 逯瑶, 龚涛}
\par\vspace{5pt}
\noindent\makebox[5em][l]{\wuhao D0384}\makebox[\dimexpr\linewidth-5em][s]{\hyperlink{paper:D0384}{\wuhao 具有“弯曲-声黑洞”辐射梁的纵弯超声辐射器}\dotfill\makebox[2.5em][r]{\wuhao\pageref{paper:D0384}}}
\par
\noindent\hspace*{5em}{\kaishu\wuhao 江昊, 章超, 叶文旭, 林书玉}
\par\vspace{5pt}
\noindent\makebox[5em][l]{\wuhao D0385}\makebox[\dimexpr\linewidth-5em][s]{\hyperlink{paper:D0385}{\wuhao 基于声黑洞结构的多模态复频超声换能器}\dotfill\makebox[2.5em][r]{\wuhao\pageref{paper:D0385}}}
\par
\noindent\hspace*{5em}{\kaishu\wuhao 叶文旭, 林书玉}
\par\vspace{5pt}
\noindent\makebox[5em][l]{\wuhao D0386}\makebox[\dimexpr\linewidth-5em][s]{\hyperlink{paper:D0386}{\wuhao 双声黑洞圆盘环形激励纵弯换能器}\dotfill\makebox[2.5em][r]{\wuhao\pageref{paper:D0386}}}
\par
\noindent\hspace*{5em}{\kaishu\wuhao 陈慧琴, 王怡, 任文博, 江昊, 林书玉}
\par\vspace{5pt}
\noindent\makebox[5em][l]{\wuhao D0392}\makebox[\dimexpr\linewidth-5em][s]{\hyperlink{paper:D0392}{\wuhao 基于声黑洞结构的三叠片弯曲振动超声切割刀}\dotfill\makebox[2.5em][r]{\wuhao\pageref{paper:D0392}}}
\par
\noindent\hspace*{5em}{\kaishu\wuhao 谭逸轩, 林书玉}
\par\vspace{5pt}
\noindent\makebox[5em][l]{\wuhao D0393}\makebox[\dimexpr\linewidth-5em][s]{\hyperlink{paper:D0393}{\wuhao 基于声黑洞结构的三叠片式超声换能器设计与性能优化}\dotfill\makebox[2.5em][r]{\wuhao\pageref{paper:D0393}}}
\par
\noindent\hspace*{5em}{\kaishu\wuhao 唐一璠, 任刘娜, 林书玉}
\par\vspace{5pt}
\noindent\makebox[5em][l]{\wuhao D0395}\makebox[\dimexpr\linewidth-5em][s]{\hyperlink{paper:D0395}{\wuhao 基于声黑洞结构的环形弯张换能器}\dotfill\makebox[2.5em][r]{\wuhao\pageref{paper:D0395}}}
\par
\noindent\hspace*{5em}{\kaishu\wuhao 章超, 江昊, 任文博, 林书玉}
\par\vspace{5pt}
\noindent\makebox[5em][l]{\wuhao D0412}\makebox[\dimexpr\linewidth-5em][s]{\hyperlink{paper:D0412}{\wuhao 高频驱动下非线性声场演变及焦点声压特征的数值模拟研究}\dotfill\makebox[2.5em][r]{\wuhao\pageref{paper:D0412}}}
\par
\noindent\hspace*{5em}{\kaishu\wuhao 罗兰, 宋丹, 熊久鹏, 李发琪}
\par\vspace{5pt}
\noindent\makebox[5em][l]{\wuhao D0421}\makebox[\dimexpr\linewidth-5em][s]{\hyperlink{paper:D0421}{\wuhao 面向功率超声应用的可调节的电荷泵超声发射芯片}\dotfill\makebox[2.5em][r]{\wuhao\pageref{paper:D0421}}}
\par
\noindent\hspace*{5em}{\kaishu\wuhao 张绪, 刘一苇, 李照希, 李迪}
\par\vspace{5pt}
\noindent\makebox[5em][l]{\wuhao D0436}\makebox[\dimexpr\linewidth-5em][s]{\hyperlink{paper:D0436}{\wuhao 超声辅助制备单分散百里香酚微乳液：物理和化学稳定性}\dotfill\makebox[2.5em][r]{\wuhao\pageref{paper:D0436}}}
\par
\noindent\hspace*{5em}{\kaishu\wuhao 武志林, 黄雪清, 李婧婧, 吴彩君}
\par\vspace{5pt}
\noindent\makebox[5em][l]{\wuhao D0438}\makebox[\dimexpr\linewidth-5em][s]{\hyperlink{paper:D0438}{\wuhao 超声辅助合成磷酸锌及其在水性防锈涂料中的防腐性能研究}\dotfill\makebox[2.5em][r]{\wuhao\pageref{paper:D0438}}}
\par
\noindent\hspace*{5em}{\kaishu\wuhao 武志林, 吴彩君}
\par\vspace{5pt}
\noindent\makebox[5em][l]{\wuhao D0439}\makebox[\dimexpr\linewidth-5em][s]{\hyperlink{paper:D0439}{\wuhao 限域流超声臭氧化反应器中有机物降解的传质强化与臭氧转化机制}\dotfill\makebox[2.5em][r]{\wuhao\pageref{paper:D0439}}}
\par
\noindent\hspace*{5em}{\kaishu\wuhao 武志林, 杨司涵}
\par\vspace{5pt}
\vspace{6pt}

\begin{center}{\heiti\xiaosi E. 检测声学}\end{center}
\par\vspace{6pt}
\noindent\makebox[5em][l]{\wuhao E0039}\makebox[\dimexpr\linewidth-5em][s]{\hyperlink{paper:E0039}{\wuhao 基于OTPA方法的烟机工作噪声贡献量研究}\dotfill\makebox[2.5em][r]{\wuhao\pageref{paper:E0039}}}
\par
\noindent\hspace*{5em}{\kaishu\wuhao 李杨铭}
\par\vspace{5pt}
\noindent\makebox[5em][l]{\wuhao E0138}\makebox[\dimexpr\linewidth-5em][s]{\hyperlink{paper:E0138}{\wuhao 基于三维立体阵列与自适应融合的无人机声学探测方法}\dotfill\makebox[2.5em][r]{\wuhao\pageref{paper:E0138}}}
\par
\noindent\hspace*{5em}{\kaishu\wuhao 刘然则}
\par\vspace{5pt}
\noindent\makebox[5em][l]{\wuhao E0147}\makebox[\dimexpr\linewidth-5em][s]{\hyperlink{paper:E0147}{\wuhao 基于稀疏阵列的超声平面波成像方法研究}\dotfill\makebox[2.5em][r]{\wuhao\pageref{paper:E0147}}}
\par
\noindent\hspace*{5em}{\kaishu\wuhao 张辉, 任科, 张海燕, 张辉, 张辉, 徐吉超}
\par\vspace{5pt}
\noindent\makebox[5em][l]{\wuhao E0164}\makebox[\dimexpr\linewidth-5em][s]{\hyperlink{paper:E0164}{\wuhao 基于扇形传感器簇阵列的板状结构损伤定位方法研究}\dotfill\makebox[2.5em][r]{\wuhao\pageref{paper:E0164}}}
\par
\noindent\hspace*{5em}{\kaishu\wuhao 韩悦, 刘金霞, 崔志文}
\par\vspace{5pt}
\noindent\makebox[5em][l]{\wuhao E0211}\makebox[\dimexpr\linewidth-5em][s]{\hyperlink{paper:E0211}{\wuhao 基于超声红外热成像技术的发动机叶片裂纹检测}\dotfill\makebox[2.5em][r]{\wuhao\pageref{paper:E0211}}}
\par
\noindent\hspace*{5em}{\kaishu\wuhao 张凯, 张辉}
\par\vspace{5pt}
\noindent\hangindent=5em\hangafter=1\makebox[5em][l]{\wuhao E0241}\hyperlink{paper:E0241}{\wuhao 基于压电晶片双谐振模式实现非线性超声导波高频基波激发和低频静态分量接收的方法}
\par\noindent\hspace*{5em}\dotfill\makebox[2.5em][r]{\wuhao\pageref{paper:E0241}}
\par\vspace{2pt}
\noindent\hspace*{5em}{\kaishu\wuhao 万翔}
\par\vspace{5pt}
\noindent\makebox[5em][l]{\wuhao E0283}\makebox[\dimexpr\linewidth-5em][s]{\hyperlink{paper:E0283}{\wuhao 基于声学超构透镜增强空耦超声的锂电池状态无损检测研究}\dotfill\makebox[2.5em][r]{\wuhao\pageref{paper:E0283}}}
\par
\noindent\hspace*{5em}{\kaishu\wuhao 靳先帝, 赵伟祥, 李鸿超, 郭淳钰, 魏彤, 徐杰, 简小华}
\par\vspace{5pt}
\noindent\makebox[5em][l]{\wuhao E0289}\makebox[\dimexpr\linewidth-5em][s]{\hyperlink{paper:E0289}{\wuhao 空耦超声-热波成像融合的高精度非接触原位检测方法}\dotfill\makebox[2.5em][r]{\wuhao\pageref{paper:E0289}}}
\par
\noindent\hspace*{5em}{\kaishu\wuhao 罗志涛, 颜迪, 靳树衫, 张辉}
\par\vspace{5pt}
\noindent\makebox[5em][l]{\wuhao E0327}\makebox[\dimexpr\linewidth-5em][s]{\hyperlink{paper:E0327}{\wuhao 基于超声导波同向共轴混频的CFRP板分层定位}\dotfill\makebox[2.5em][r]{\wuhao\pageref{paper:E0327}}}
\par
\noindent\hspace*{5em}{\kaishu\wuhao 孙茂循, 李龙飞, 刘禹恒, 刘宏业}
\par\vspace{5pt}
\noindent\makebox[5em][l]{\wuhao E0437}\makebox[\dimexpr\linewidth-5em][s]{\hyperlink{paper:E0437}{\wuhao 壁面入射声波传播速度研究}\dotfill\makebox[2.5em][r]{\wuhao\pageref{paper:E0437}}}
\par
\noindent\hspace*{5em}{\kaishu\wuhao 卢亮亮, 李言钦}
\par\vspace{5pt}
\noindent\makebox[5em][l]{\wuhao E0461}\makebox[\dimexpr\linewidth-5em][s]{\hyperlink{paper:E0461}{\wuhao 未知温度场下火电锅炉水冷壁多源泄漏声学定位方法}\dotfill\makebox[2.5em][r]{\wuhao\pageref{paper:E0461}}}
\par
\noindent\hspace*{5em}{\kaishu\wuhao 孙建浩, 周玉, 高媛毓, 李颢, 姜根山}
\par\vspace{5pt}
\vspace{6pt}

\begin{center}{\heiti\xiaosi G.生物医学超声}\end{center}
\par\vspace{6pt}
\noindent\hangindent=5em\hangafter=1\makebox[5em][l]{\wuhao G0004}\hyperlink{paper:G0004}{\wuhao 光声影像组学特征和临床因素的整合用于识别乳腺癌患者 Ki-67 状态}
\par\noindent\hspace*{5em}\dotfill\makebox[2.5em][r]{\wuhao\pageref{paper:G0004}}
\par\vspace{2pt}
\noindent\hspace*{5em}{\kaishu\wuhao 王梦云, 郭乐杭}
\par\vspace{5pt}
\noindent\hangindent=5em\hangafter=1\makebox[5em][l]{\wuhao G0005}\hyperlink{paper:G0005}{\wuhao 深度学习与多模态光声/超声成像相结合早期区分乳腺癌中的管腔型和非管腔型肿瘤}
\par\noindent\hspace*{5em}\dotfill\makebox[2.5em][r]{\wuhao\pageref{paper:G0005}}
\par\vspace{2pt}
\noindent\hspace*{5em}{\kaishu\wuhao 王梦云, 郭乐杭}
\par\vspace{5pt}
\noindent\hangindent=5em\hangafter=1\makebox[5em][l]{\wuhao G0006}\hyperlink{paper:G0006}{\wuhao Deep Learning Enhanced Multimodal Photoacoustic/Ultrasound Imaging for Noninvasive Prediction of Ki 67 Expression in Breast Cancer}
\par\noindent\hspace*{5em}\dotfill\makebox[2.5em][r]{\wuhao\pageref{paper:G0006}}
\par\vspace{2pt}
\noindent\hspace*{5em}{\kaishu\wuhao 王梦云, 郭乐杭}
\par\vspace{5pt}
\noindent\hangindent=5em\hangafter=1\makebox[5em][l]{\wuhao G0007}\hyperlink{paper:G0007}{\wuhao 超声衰减分析和受控衰减参数在非酒精性脂肪性肝病患者中肝脂肪变性诊断和分级的比较研究}
\par\noindent\hspace*{5em}\dotfill\makebox[2.5em][r]{\wuhao\pageref{paper:G0007}}
\par\vspace{2pt}
\noindent\hspace*{5em}{\kaishu\wuhao 王梦云, 郭乐杭}
\par\vspace{5pt}
\noindent\hangindent=5em\hangafter=1\makebox[5em][l]{\wuhao G0008}\hyperlink{paper:G0008}{\wuhao Delta Radiomics Based on Longitudinal Ultrasound for Early Prediction of Response to Neoadjuvant Chemotherapy in Breast Cancer}
\par\noindent\hspace*{5em}\dotfill\makebox[2.5em][r]{\wuhao\pageref{paper:G0008}}
\par\vspace{2pt}
\noindent\hspace*{5em}{\kaishu\wuhao 王梦云, 郭乐杭}
\par\vspace{5pt}
\noindent\hangindent=5em\hangafter=1\makebox[5em][l]{\wuhao G0009}\hyperlink{paper:G0009}{\wuhao 深度学习与多模态光声/超声成像结合以术前预测乳腺癌患者Ki-67表达}
\par\noindent\hspace*{5em}\dotfill\makebox[2.5em][r]{\wuhao\pageref{paper:G0009}}
\par\vspace{2pt}
\noindent\hspace*{5em}{\kaishu\wuhao 郭乐杭, 王梦云}
\par\vspace{5pt}
\noindent\hangindent=5em\hangafter=1\makebox[5em][l]{\wuhao G0010}\hyperlink{paper:G0010}{\wuhao Interpretable Prediction of Ki-67 Expression in Breast Cancer through Integrated Photoacoustic Imaging Radiomics and Clinical Parameters}
\par\noindent\hspace*{5em}\dotfill\makebox[2.5em][r]{\wuhao\pageref{paper:G0010}}
\par\vspace{2pt}
\noindent\hspace*{5em}{\kaishu\wuhao 郭乐杭, 王梦云}
\par\vspace{5pt}
\noindent\hangindent=5em\hangafter=1\makebox[5em][l]{\wuhao G0011}\hyperlink{paper:G0011}{\wuhao Deep learning combined with photoacoustic imaging to differentiate between luminal and non-luminal tumors in early breast cancer}
\par\noindent\hspace*{5em}\dotfill\makebox[2.5em][r]{\wuhao\pageref{paper:G0011}}
\par\vspace{2pt}
\noindent\hspace*{5em}{\kaishu\wuhao 郭乐杭, 王梦云}
\par\vspace{5pt}
\noindent\hangindent=5em\hangafter=1\makebox[5em][l]{\wuhao G0012}\hyperlink{paper:G0012}{\wuhao 基于纵向超声 Delta 放射组学的乳腺癌新辅助化疗肿瘤反应早期分层预测}
\par\noindent\hspace*{5em}\dotfill\makebox[2.5em][r]{\wuhao\pageref{paper:G0012}}
\par\vspace{2pt}
\noindent\hspace*{5em}{\kaishu\wuhao 郭乐杭, 王梦云}
\par\vspace{5pt}
\noindent\hangindent=5em\hangafter=1\makebox[5em][l]{\wuhao G0013}\hyperlink{paper:G0013}{\wuhao Ultrasound Attenuation Analysis and Controlled Attenuation Parameter for Liver Steatosis Diagnosis in Non-alcoholic Fatty Liver Disease: A Comparative Study}
\par\noindent\hspace*{5em}\dotfill\makebox[2.5em][r]{\wuhao\pageref{paper:G0013}}
\par\vspace{2pt}
\noindent\hspace*{5em}{\kaishu\wuhao 郭乐杭, 王梦云}
\par\vspace{5pt}
\noindent\makebox[5em][l]{\wuhao G0034}\makebox[\dimexpr\linewidth-5em][s]{\hyperlink{paper:G0034}{\wuhao 脊柱外科手术中椎骨超声精准定位和位姿实时监测技术研究}\dotfill\makebox[2.5em][r]{\wuhao\pageref{paper:G0034}}}
\par
\noindent\hspace*{5em}{\kaishu\wuhao 蒋希哲, 陈玉宣, 周红生, 刘逍逸}
\par\vspace{5pt}
\noindent\makebox[5em][l]{\wuhao G0060}\makebox[\dimexpr\linewidth-5em][s]{\hyperlink{paper:G0060}{\wuhao 基于超快超声的 B-SMI-ULM 一体化跨尺度介观血管成像平台}\dotfill\makebox[2.5em][r]{\wuhao\pageref{paper:G0060}}}
\par
\noindent\hspace*{5em}{\kaishu\wuhao 龚新越, 杜秉泽, 张国峰, 朱逸斐, 薛洪惠, 屠娟}
\par\vspace{5pt}
\noindent\makebox[5em][l]{\wuhao G0071}\makebox[\dimexpr\linewidth-5em][s]{\hyperlink{paper:G0071}{\wuhao 硅基纳米颗粒调控超声空化及其在癌症治疗中的应用}\dotfill\makebox[2.5em][r]{\wuhao\pageref{paper:G0071}}}
\par
\noindent\hspace*{5em}{\kaishu\wuhao 张小曼, 曹文武}
\par\vspace{5pt}
\noindent\makebox[5em][l]{\wuhao G0075}\makebox[\dimexpr\linewidth-5em][s]{\hyperlink{paper:G0075}{\wuhao 基于超声射频信号的肌肉内脂肪浸润声衰减表征方法研究}\dotfill\makebox[2.5em][r]{\wuhao\pageref{paper:G0075}}}
\par
\noindent\hspace*{5em}{\kaishu\wuhao 魏轲, 屠娟, 章东}
\par\vspace{5pt}
\noindent\makebox[5em][l]{\wuhao G0076}\makebox[\dimexpr\linewidth-5em][s]{\hyperlink{paper:G0076}{\wuhao 低强度脉冲超声改善小鼠糖尿病肾病的效果及机制研究}\dotfill\makebox[2.5em][r]{\wuhao\pageref{paper:G0076}}}
\par
\noindent\hspace*{5em}{\kaishu\wuhao 王欣, 唐量}
\par\vspace{5pt}
\noindent\makebox[5em][l]{\wuhao G0101}\makebox[\dimexpr\linewidth-5em][s]{\hyperlink{paper:G0101}{\wuhao 深度学习的大鼠脑功率多普勒微血管成像研究}\dotfill\makebox[2.5em][r]{\wuhao\pageref{paper:G0101}}}
\par
\noindent\hspace*{5em}{\kaishu\wuhao 王睿菲, 安志武, 柳雪丽, 吴文焘, 张秀梅}
\par\vspace{5pt}
\noindent\makebox[5em][l]{\wuhao G0103}\makebox[\dimexpr\linewidth-5em][s]{\hyperlink{paper:G0103}{\wuhao 用于在体光声微血管成像的深度补偿型针状光束}\dotfill\makebox[2.5em][r]{\wuhao\pageref{paper:G0103}}}
\par
\noindent\hspace*{5em}{\kaishu\wuhao 杨一帆, 陶超, 刘晓峻}
\par\vspace{5pt}
\noindent\makebox[5em][l]{\wuhao G0104}\makebox[\dimexpr\linewidth-5em][s]{\hyperlink{paper:G0104}{\wuhao 自动聚焦光学分辨率光声显微镜}\dotfill\makebox[2.5em][r]{\wuhao\pageref{paper:G0104}}}
\par
\noindent\hspace*{5em}{\kaishu\wuhao 左天翔, 陶超, 刘晓峻}
\par\vspace{5pt}
\noindent\makebox[5em][l]{\wuhao G0155}\makebox[\dimexpr\linewidth-5em][s]{\hyperlink{paper:G0155}{\wuhao 跟腱断裂康复监测的可穿戴超声定量评估研究}\dotfill\makebox[2.5em][r]{\wuhao\pageref{paper:G0155}}}
\par
\noindent\hspace*{5em}{\kaishu\wuhao 柳雪丽, 吴文焘, 王睿菲, 安志武}
\par\vspace{5pt}
\noindent\makebox[5em][l]{\wuhao G0156}\makebox[\dimexpr\linewidth-5em][s]{\hyperlink{paper:G0156}{\wuhao 基于多物理场耦合的超声驱动双微泡系统血管生物力学效应仿真研究}\dotfill\makebox[2.5em][r]{\wuhao\pageref{paper:G0156}}}
\par
\noindent\hspace*{5em}{\kaishu\wuhao 上官哲韬, 纪镇祥, 刘怡然, 杨若妍, 索鼎杰}
\par\vspace{5pt}
\noindent\makebox[5em][l]{\wuhao G0170}\makebox[\dimexpr\linewidth-5em][s]{\hyperlink{paper:G0170}{\wuhao 40 Hz脉冲超声刺激改善AD小鼠的认知功能障碍}\dotfill\makebox[2.5em][r]{\wuhao\pageref{paper:G0170}}}
\par
\noindent\hspace*{5em}{\kaishu\wuhao 易沙沙, 牛丽丽}
\par\vspace{5pt}
\noindent\makebox[5em][l]{\wuhao G0193}\makebox[\dimexpr\linewidth-5em][s]{\hyperlink{paper:G0193}{\wuhao 基于双曲型聚声换能器的机械旋转式超声针灸斜入射方法研究}\dotfill\makebox[2.5em][r]{\wuhao\pageref{paper:G0193}}}
\par
\noindent\hspace*{5em}{\kaishu\wuhao 周姗, 刘欣宇, 凤飞龙, 李锦}
\par\vspace{5pt}
\noindent\makebox[5em][l]{\wuhao G0201}\makebox[\dimexpr\linewidth-5em][s]{\hyperlink{paper:G0201}{\wuhao 菲涅尔波带片频控调焦技术在穴位超声理疗中的应用研究}\dotfill\makebox[2.5em][r]{\wuhao\pageref{paper:G0201}}}
\par
\noindent\hspace*{5em}{\kaishu\wuhao 刘欣宇, 周姗, 凤飞龙, 李锦}
\par\vspace{5pt}
\noindent\makebox[5em][l]{\wuhao G0205}\makebox[\dimexpr\linewidth-5em][s]{\hyperlink{paper:G0205}{\wuhao 低频超声联合微泡在不同黏弹性介质中的仿真研究}\dotfill\makebox[2.5em][r]{\wuhao\pageref{paper:G0205}}}
\par
\noindent\hspace*{5em}{\kaishu\wuhao 沈智勇, 余慧, 范璐瑶}
\par\vspace{5pt}
\noindent\makebox[5em][l]{\wuhao G0271}\makebox[\dimexpr\linewidth-5em][s]{\hyperlink{paper:G0271}{\wuhao 基于深度学习的超分辨超声成像运动矫正研究及应用}\dotfill\makebox[2.5em][r]{\wuhao\pageref{paper:G0271}}}
\par
\noindent\hspace*{5em}{\kaishu\wuhao 高乐, 潘玥, 张志强, 邱维宝}
\par\vspace{5pt}
\noindent\makebox[5em][l]{\wuhao G0281}\makebox[\dimexpr\linewidth-5em][s]{\hyperlink{paper:G0281}{\wuhao 微型穿戴式超声联合载药水凝胶经皮促渗研究}\dotfill\makebox[2.5em][r]{\wuhao\pageref{paper:G0281}}}
\par
\noindent\hspace*{5em}{\kaishu\wuhao 靳先帝, 郭淳钰, 李鸿超, 魏彤, 赵子枫, 吕锦涵, 徐杰, 简小华}
\par\vspace{5pt}
\noindent\makebox[5em][l]{\wuhao G0285}\makebox[\dimexpr\linewidth-5em][s]{\hyperlink{paper:G0285}{\wuhao 面向情绪障碍干预的非侵入式声学脑-迷走神经接口}\dotfill\makebox[2.5em][r]{\wuhao\pageref{paper:G0285}}}
\par
\noindent\hspace*{5em}{\kaishu\wuhao 朱一跃, Niu Lili}
\par\vspace{5pt}
\noindent\makebox[5em][l]{\wuhao G0286}\makebox[\dimexpr\linewidth-5em][s]{\hyperlink{paper:G0286}{\wuhao 基于超声定位显微成像骨架化分析框架的癌症转移微血管重塑表征}\dotfill\makebox[2.5em][r]{\wuhao\pageref{paper:G0286}}}
\par
\noindent\hspace*{5em}{\kaishu\wuhao 杜秉泽, 姜岭寅, 吴意赟, 丁波, 屠娟}
\par\vspace{5pt}
\noindent\makebox[5em][l]{\wuhao G0302}\makebox[\dimexpr\linewidth-5em][s]{\hyperlink{paper:G0302}{\wuhao 负向声辐射力的物理机制、结构设计与调控策略}\dotfill\makebox[2.5em][r]{\wuhao\pageref{paper:G0302}}}
\par
\noindent\hspace*{5em}{\kaishu\wuhao 宫门阳}
\par\vspace{5pt}
\noindent\makebox[5em][l]{\wuhao G0309}\makebox[\dimexpr\linewidth-5em][s]{\hyperlink{paper:G0309}{\wuhao 心-机接口：脑-机接口的声学并行范式}\dotfill\makebox[2.5em][r]{\wuhao\pageref{paper:G0309}}}
\par
\noindent\hspace*{5em}{\kaishu\wuhao 高翔, 郭建中}
\par\vspace{5pt}
\noindent\makebox[5em][l]{\wuhao G0321}\makebox[\dimexpr\linewidth-5em][s]{\hyperlink{paper:G0321}{\wuhao 面向体外膜肺氧合循环场景的超声流量测量系统研究}\dotfill\makebox[2.5em][r]{\wuhao\pageref{paper:G0321}}}
\par
\noindent\hspace*{5em}{\kaishu\wuhao 魏彤, 简小华, 徐杰, 靳先帝, 吕锦涵}
\par\vspace{5pt}
\noindent\makebox[5em][l]{\wuhao G0340}\makebox[\dimexpr\linewidth-5em][s]{\hyperlink{paper:G0340}{\wuhao 超声空化效应增强化疗对卵巢癌细胞的凋亡作用}\dotfill\makebox[2.5em][r]{\wuhao\pageref{paper:G0340}}}
\par
\noindent\hspace*{5em}{\kaishu\wuhao 沈智勇, 余慧, 范璐瑶}
\par\vspace{5pt}
\noindent\makebox[5em][l]{\wuhao G0350}\makebox[\dimexpr\linewidth-5em][s]{\hyperlink{paper:G0350}{\wuhao 高强度聚焦超声热疗中热电偶粘性热效应的建模分析与修正方法研究}\dotfill\makebox[2.5em][r]{\wuhao\pageref{paper:G0350}}}
\par
\noindent\hspace*{5em}{\kaishu\wuhao 丁墨涵}
\par\vspace{5pt}
\noindent\makebox[5em][l]{\wuhao G0354}\makebox[\dimexpr\linewidth-5em][s]{\hyperlink{paper:G0354}{\wuhao 相干非线性超分辨率超声毛细血管成像}\dotfill\makebox[2.5em][r]{\wuhao\pageref{paper:G0354}}}
\par
\noindent\hspace*{5em}{\kaishu\wuhao 程双毅, 徐龙, 许凯亮, 他得安}
\par\vspace{5pt}
\noindent\makebox[5em][l]{\wuhao G0357}\makebox[\dimexpr\linewidth-5em][s]{\hyperlink{paper:G0357}{\wuhao 基于超材料匹配层的经颅超声透射增强研究}\dotfill\makebox[2.5em][r]{\wuhao\pageref{paper:G0357}}}
\par
\noindent\hspace*{5em}{\kaishu\wuhao 豆乔, 田野}
\par\vspace{5pt}
\noindent\makebox[5em][l]{\wuhao G0367}\makebox[\dimexpr\linewidth-5em][s]{\hyperlink{paper:G0367}{\wuhao 超声引导药物递送}\dotfill\makebox[2.5em][r]{\wuhao\pageref{paper:G0367}}}
\par
\noindent\hspace*{5em}{\kaishu\wuhao 刘芷麟}
\par\vspace{5pt}
\noindent\makebox[5em][l]{\wuhao G0369}\makebox[\dimexpr\linewidth-5em][s]{\hyperlink{paper:G0369}{\wuhao 独立于听觉干扰的小鼠视觉环路超声调控研究}\dotfill\makebox[2.5em][r]{\wuhao\pageref{paper:G0369}}}
\par
\noindent\hspace*{5em}{\kaishu\wuhao 何佳儒, 朱杰君, 王新新, 陈子浩, 杨晋, 袁振, 徐宏治, 孙雷, 丘志海}
\par\vspace{5pt}
\noindent\makebox[5em][l]{\wuhao G0396}\makebox[\dimexpr\linewidth-5em][s]{\hyperlink{paper:G0396}{\wuhao 微纳粒子的声涡旋操控及其介导超声治疗研究}\dotfill\makebox[2.5em][r]{\wuhao\pageref{paper:G0396}}}
\par
\noindent\hspace*{5em}{\kaishu\wuhao 郭世放, 胡心茹, 贾万琳, 牙振, 李燕, 万明习}
\par\vspace{5pt}
\noindent\hangindent=5em\hangafter=1\makebox[5em][l]{\wuhao G0397}\hyperlink{paper:G0397}{\wuhao 基于低-高机械指数序贯超声刺激微泡空化的肝癌治疗及超声多模态评价}
\par\noindent\hspace*{5em}\dotfill\makebox[2.5em][r]{\wuhao\pageref{paper:G0397}}
\par\vspace{2pt}
\noindent\hspace*{5em}{\kaishu\wuhao 杨含怡, 张莹, 王斐倩, 张洁, 吴铮, 李小静, 万明习, 宗瑜瑾}
\par\vspace{5pt}
\noindent\makebox[5em][l]{\wuhao G0401}\makebox[\dimexpr\linewidth-5em][s]{\hyperlink{paper:G0401}{\wuhao 超分辨显微成像技术鉴别诊断肝结节病与肝细胞癌的初步研究}\dotfill\makebox[2.5em][r]{\wuhao\pageref{paper:G0401}}}
\par
\noindent\hspace*{5em}{\kaishu\wuhao 张洁, 吴铮, 李小静, 芦幸琪, 于景彤, 王斐倩}
\par\vspace{5pt}
\noindent\hangindent=5em\hangafter=1\makebox[5em][l]{\wuhao G0404}\hyperlink{paper:G0404}{\wuhao 基于高频率聚焦超声热效应调控TRPV1通路治疗糖尿病大鼠神经伤害炎症的研究​}
\par\noindent\hspace*{5em}\dotfill\makebox[2.5em][r]{\wuhao\pageref{paper:G0404}}
\par\vspace{2pt}
\noindent\hspace*{5em}{\kaishu\wuhao 丁予婷, 郭乐杭}
\par\vspace{5pt}
\noindent\makebox[5em][l]{\wuhao G0417}\makebox[\dimexpr\linewidth-5em][s]{\hyperlink{paper:G0417}{\wuhao 基于聚焦超声激发组织机械响应的声能量反演研究}\dotfill\makebox[2.5em][r]{\wuhao\pageref{paper:G0417}}}
\par
\noindent\hspace*{5em}{\kaishu\wuhao 朱小仡, 宋丹, 熊久鹏, 李发琪}
\par\vspace{5pt}
\noindent\makebox[5em][l]{\wuhao G0425}\makebox[\dimexpr\linewidth-5em][s]{\hyperlink{paper:G0425}{\wuhao 基于超声RF信号的下肢样组织声速与密度级联重建}\dotfill\makebox[2.5em][r]{\wuhao\pageref{paper:G0425}}}
\par
\noindent\hspace*{5em}{\kaishu\wuhao 金义博}
\par\vspace{5pt}
\noindent\makebox[5em][l]{\wuhao G0426}\makebox[\dimexpr\linewidth-5em][s]{\hyperlink{paper:G0426}{\wuhao 基于多对称轴相干因子的分层介质声速-厚度联合估计算法研究}\dotfill\makebox[2.5em][r]{\wuhao\pageref{paper:G0426}}}
\par
\noindent\hspace*{5em}{\kaishu\wuhao 刘欣澳, 吕倩, 郭建中}
\par\vspace{5pt}
\noindent\makebox[5em][l]{\wuhao G0428}\makebox[\dimexpr\linewidth-5em][s]{\hyperlink{paper:G0428}{\wuhao 基于隐式神经网络的骨超声全波形反演成像}\dotfill\makebox[2.5em][r]{\wuhao\pageref{paper:G0428}}}
\par
\noindent\hspace*{5em}{\kaishu\wuhao 吕冰泽}
\par\vspace{5pt}
\noindent\hangindent=5em\hangafter=1\makebox[5em][l]{\wuhao G0434}\hyperlink{paper:G0434}{\wuhao 基于超声定位显微镜微血管异质性精准评估胶质母细胞瘤病理恶性分层}
\par\noindent\hspace*{5em}\dotfill\makebox[2.5em][r]{\wuhao\pageref{paper:G0434}}
\par\vspace{2pt}
\noindent\hspace*{5em}{\kaishu\wuhao 胡星}
\par\vspace{5pt}
\noindent\makebox[5em][l]{\wuhao G0441}\makebox[\dimexpr\linewidth-5em][s]{\hyperlink{paper:G0441}{\wuhao 基于迁移学习物理信息神经网络的剪切波弹性快速定量反演方法}\dotfill\makebox[2.5em][r]{\wuhao\pageref{paper:G0441}}}
\par
\noindent\hspace*{5em}{\kaishu\wuhao 刘佳琦, 黄凯, 熊久鹏, 李发琪}
\par\vspace{5pt}
\noindent\makebox[5em][l]{\wuhao G0460}\makebox[\dimexpr\linewidth-5em][s]{\hyperlink{paper:G0460}{\wuhao 基于多角度平面波成像的组织声速估计方法研究}\dotfill\makebox[2.5em][r]{\wuhao\pageref{paper:G0460}}}
\par
\noindent\hspace*{5em}{\kaishu\wuhao 司玉琪, 吕倩, 郭建中}
\par\vspace{5pt}
\vspace{6pt}

\begin{center}{\heiti\xiaosi H.微声学（包括声表面波、薄膜体声波、微纳器件及其相关的系统与材料）}\end{center}
\par\vspace{6pt}
\noindent\hangindent=5em\hangafter=1\makebox[5em][l]{\wuhao H0053}\hyperlink{paper:H0053}{\wuhao 用于光纤EFPI水听器信号解调的滤波补偿与幅值反馈双波长正交相位方法}
\par\noindent\hspace*{5em}\dotfill\makebox[2.5em][r]{\wuhao\pageref{paper:H0053}}
\par\vspace{2pt}
\noindent\hspace*{5em}{\kaishu\wuhao 陈小军, 葛辉良, 刘盛春}
\par\vspace{5pt}
\noindent\makebox[5em][l]{\wuhao H0139}\makebox[\dimexpr\linewidth-5em][s]{\hyperlink{paper:H0139}{\wuhao 基于铁磁共振效应的高灵敏度声表面波磁传感器}\dotfill\makebox[2.5em][r]{\wuhao\pageref{paper:H0139}}}
\par
\noindent\hspace*{5em}{\kaishu\wuhao 罗为}
\par\vspace{5pt}
\noindent\makebox[5em][l]{\wuhao H0259}\makebox[\dimexpr\linewidth-5em][s]{\hyperlink{paper:H0259}{\wuhao NaX分子筛修饰声表面波传感器的快速高灵敏一氧化碳检测}\dotfill\makebox[2.5em][r]{\wuhao\pageref{paper:H0259}}}
\par
\noindent\hspace*{5em}{\kaishu\wuhao 杨启明, 金晶, 胡安宇, 薛蓄峰, 梁勇, 王文}
\par\vspace{5pt}
\noindent\makebox[5em][l]{\wuhao H0295}\makebox[\dimexpr\linewidth-5em][s]{\hyperlink{paper:H0295}{\wuhao 基于声表面波的无线无源宽温域温度传感系统}\dotfill\makebox[2.5em][r]{\wuhao\pageref{paper:H0295}}}
\par
\noindent\hspace*{5em}{\kaishu\wuhao 孙铭辰, 王文}
\par\vspace{5pt}
\noindent\makebox[5em][l]{\wuhao H0298}\makebox[\dimexpr\linewidth-5em][s]{\hyperlink{paper:H0298}{\wuhao 声表面波氢气传感器恒温设计和控制方法}\dotfill\makebox[2.5em][r]{\wuhao\pageref{paper:H0298}}}
\par
\noindent\hspace*{5em}{\kaishu\wuhao 黄丽洁, 王文, 程利娜, 崔柏乐}
\par\vspace{5pt}
\noindent\makebox[5em][l]{\wuhao H0299}\makebox[\dimexpr\linewidth-5em][s]{\hyperlink{paper:H0299}{\wuhao 基于声表面波色散延迟线的实时宽带频谱分析技术}\dotfill\makebox[2.5em][r]{\wuhao\pageref{paper:H0299}}}
\par
\noindent\hspace*{5em}{\kaishu\wuhao 户媛媛, 刘梦伟}
\par\vspace{5pt}
\noindent\makebox[5em][l]{\wuhao H0306}\makebox[\dimexpr\linewidth-5em][s]{\hyperlink{paper:H0306}{\wuhao 基于POI异质结构的悬臂梁式SAW振动传感器优化设计}\dotfill\makebox[2.5em][r]{\wuhao\pageref{paper:H0306}}}
\par
\noindent\hspace*{5em}{\kaishu\wuhao 于冰冰, 王文}
\par\vspace{5pt}
\noindent\makebox[5em][l]{\wuhao H0322}\makebox[\dimexpr\linewidth-5em][s]{\hyperlink{paper:H0322}{\wuhao 基于耦合微泡振动的选择性颗粒操控}\dotfill\makebox[2.5em][r]{\wuhao\pageref{paper:H0322}}}
\par
\noindent\hspace*{5em}{\kaishu\wuhao 李昕珈, 全浩, 刘秀芳, 孟龙}
\par\vspace{5pt}
\noindent\makebox[5em][l]{\wuhao H0355}\makebox[\dimexpr\linewidth-5em][s]{\hyperlink{paper:H0355}{\wuhao 基于SAW高频谐振器的高灵敏度扭矩传感器研究}\dotfill\makebox[2.5em][r]{\wuhao\pageref{paper:H0355}}}
\par
\noindent\hspace*{5em}{\kaishu\wuhao 田亚会, 杨智超}
\par\vspace{5pt}
\noindent\makebox[5em][l]{\wuhao H0408}\makebox[\dimexpr\linewidth-5em][s]{\hyperlink{paper:H0408}{\wuhao 面向精准测距与孔径感知的全集成PMUT超声收发系统}\dotfill\makebox[2.5em][r]{\wuhao\pageref{paper:H0408}}}
\par
\noindent\hspace*{5em}{\kaishu\wuhao 许牧原, 李迪, 李照希, 杨银堂}
\par\vspace{5pt}
\vspace{6pt}

\begin{center}{\heiti\xiaosi I. 环境声学}\end{center}
\par\vspace{6pt}
\noindent\makebox[5em][l]{\wuhao I0029}\makebox[\dimexpr\linewidth-5em][s]{\hyperlink{paper:I0029}{\wuhao 室内声场特性对调幅声烦恼度的影响及形成机理研究}\dotfill\makebox[2.5em][r]{\wuhao\pageref{paper:I0029}}}
\par
\noindent\hspace*{5em}{\kaishu\wuhao 赵焕绮, 陈克安}
\par\vspace{5pt}
\noindent\makebox[5em][l]{\wuhao I0032}\makebox[\dimexpr\linewidth-5em][s]{\hyperlink{paper:I0032}{\wuhao 揭示建筑环境设计中的声景机制：一种多模态知识图谱方法}\dotfill\makebox[2.5em][r]{\wuhao\pageref{paper:I0032}}}
\par
\noindent\hspace*{5em}{\kaishu\wuhao 秦金辉}
\par\vspace{5pt}
\noindent\makebox[5em][l]{\wuhao I0033}\makebox[\dimexpr\linewidth-5em][s]{\hyperlink{paper:I0033}{\wuhao 交通导向型发展区域声环境中声学指标对声景感知和生理指标的影响}\dotfill\makebox[2.5em][r]{\wuhao\pageref{paper:I0033}}}
\par
\noindent\hspace*{5em}{\kaishu\wuhao 秦金辉, 马英}
\par\vspace{5pt}
\noindent\makebox[5em][l]{\wuhao I0046}\makebox[\dimexpr\linewidth-5em][s]{\hyperlink{paper:I0046}{\wuhao GIS支持下的城市三位建筑格局对道路交通噪声的影响}\dotfill\makebox[2.5em][r]{\wuhao\pageref{paper:I0046}}}
\par
\noindent\hspace*{5em}{\kaishu\wuhao 李静文, 王亚婷, 张韬, 谯博文}
\par\vspace{5pt}
\noindent\makebox[5em][l]{\wuhao I0056}\makebox[\dimexpr\linewidth-5em][s]{\hyperlink{paper:I0056}{\wuhao 城中村步行街噪声对临街长租公寓的影响研究——以深圳南头古城为例}\dotfill\makebox[2.5em][r]{\wuhao\pageref{paper:I0056}}}
\par
\noindent\hspace*{5em}{\kaishu\wuhao 庄锐}
\par\vspace{5pt}
\noindent\makebox[5em][l]{\wuhao I0058}\makebox[\dimexpr\linewidth-5em][s]{\hyperlink{paper:I0058}{\wuhao 面部表情分析在城市开放空间交通噪声感知中的应用}\dotfill\makebox[2.5em][r]{\wuhao\pageref{paper:I0058}}}
\par
\noindent\hspace*{5em}{\kaishu\wuhao 胡雪筠, 孟琪, 方照川}
\par\vspace{5pt}
\noindent\makebox[5em][l]{\wuhao I0114}\makebox[\dimexpr\linewidth-5em][s]{\hyperlink{paper:I0114}{\wuhao 深圳市工业噪声污染特征分析及防治对策研究}\dotfill\makebox[2.5em][r]{\wuhao\pageref{paper:I0114}}}
\par
\noindent\hspace*{5em}{\kaishu\wuhao 许沛晨, 杨娜}
\par\vspace{5pt}
\noindent\makebox[5em][l]{\wuhao I0157}\makebox[\dimexpr\linewidth-5em][s]{\hyperlink{paper:I0157}{\wuhao 在蓝色空间声景中慢性稳态噪声暴露者的恢复性研究}\dotfill\makebox[2.5em][r]{\wuhao\pageref{paper:I0157}}}
\par
\noindent\hspace*{5em}{\kaishu\wuhao 霍光大, 秦金辉}
\par\vspace{5pt}
\noindent\makebox[5em][l]{\wuhao I0165}\makebox[\dimexpr\linewidth-5em][s]{\hyperlink{paper:I0165}{\wuhao 视听交互作用下居住环境声景感知及其生理效应}\dotfill\makebox[2.5em][r]{\wuhao\pageref{paper:I0165}}}
\par
\noindent\hspace*{5em}{\kaishu\wuhao 黄议霆, 宋阳}
\par\vspace{5pt}
\noindent\makebox[5em][l]{\wuhao I0189}\makebox[\dimexpr\linewidth-5em][s]{\hyperlink{paper:I0189}{\wuhao 空间站声环境感知理论的实证大数据扎根理论建构}\dotfill\makebox[2.5em][r]{\wuhao\pageref{paper:I0189}}}
\par
\noindent\hspace*{5em}{\kaishu\wuhao 张焱, 余磊, 康健}
\par\vspace{5pt}
\noindent\makebox[5em][l]{\wuhao I0216}\makebox[\dimexpr\linewidth-5em][s]{\hyperlink{paper:I0216}{\wuhao 低空飞行器噪声传播预测模型构建及验证}\dotfill\makebox[2.5em][r]{\wuhao\pageref{paper:I0216}}}
\par
\noindent\hspace*{5em}{\kaishu\wuhao 纪艳春, 何芝华}
\par\vspace{5pt}
\noindent\makebox[5em][l]{\wuhao I0226}\makebox[\dimexpr\linewidth-5em][s]{\hyperlink{paper:I0226}{\wuhao 年龄相关性听力损失对城市声景感知的影响：声源差异与心理声学机制}\dotfill\makebox[2.5em][r]{\wuhao\pageref{paper:I0226}}}
\par
\noindent\hspace*{5em}{\kaishu\wuhao 谭梦姣, 梁林达}
\par\vspace{5pt}
\noindent\hangindent=5em\hangafter=1\makebox[5em][l]{\wuhao I0236}\hyperlink{paper:I0236}{\wuhao 面向精细化监测的环境噪声声源智能识别技术研究——基于手持式监测终端}
\par\noindent\hspace*{5em}\dotfill\makebox[2.5em][r]{\wuhao\pageref{paper:I0236}}
\par\vspace{2pt}
\noindent\hspace*{5em}{\kaishu\wuhao 钱亚勤}
\par\vspace{5pt}
\noindent\makebox[5em][l]{\wuhao I0273}\makebox[\dimexpr\linewidth-5em][s]{\hyperlink{paper:I0273}{\wuhao 城市宁静空间感知的影响因素及噪声阈值研究}\dotfill\makebox[2.5em][r]{\wuhao\pageref{paper:I0273}}}
\par
\noindent\hspace*{5em}{\kaishu\wuhao 舒珊, 王梦雪, 冉梦轩}
\par\vspace{5pt}
\noindent\makebox[5em][l]{\wuhao I0310}\makebox[\dimexpr\linewidth-5em][s]{\hyperlink{paper:I0310}{\wuhao 城市声景增强的掩蔽声自动选择：一种面向决策的价值型离线学习方法}\dotfill\makebox[2.5em][r]{\wuhao\pageref{paper:I0310}}}
\par
\noindent\hspace*{5em}{\kaishu\wuhao 黄榆棋, 宋阳}
\par\vspace{5pt}
\noindent\makebox[5em][l]{\wuhao I0323}\makebox[\dimexpr\linewidth-5em][s]{\hyperlink{paper:I0323}{\wuhao 基于虚拟现实与声学听觉化的无人机噪声对公众情绪响应的影响研究}\dotfill\makebox[2.5em][r]{\wuhao\pageref{paper:I0323}}}
\par
\noindent\hspace*{5em}{\kaishu\wuhao 郑希桐}
\par\vspace{5pt}
\noindent\makebox[5em][l]{\wuhao I0329}\makebox[\dimexpr\linewidth-5em][s]{\hyperlink{paper:I0329}{\wuhao 不同交通流情境下临干道口袋公园声景干预的对比研究}\dotfill\makebox[2.5em][r]{\wuhao\pageref{paper:I0329}}}
\par
\noindent\hspace*{5em}{\kaishu\wuhao 马建军, 路晓东, 祝培生}
\par\vspace{5pt}
\noindent\makebox[5em][l]{\wuhao I0330}\makebox[\dimexpr\linewidth-5em][s]{\hyperlink{paper:I0330}{\wuhao 地铁TOD建筑中二次辐射噪声多维负面感受评价与预测模型研究}\dotfill\makebox[2.5em][r]{\wuhao\pageref{paper:I0330}}}
\par
\noindent\hspace*{5em}{\kaishu\wuhao 王翘楚}
\par\vspace{5pt}
\noindent\makebox[5em][l]{\wuhao I0380}\makebox[\dimexpr\linewidth-5em][s]{\hyperlink{paper:I0380}{\wuhao 基于听觉特征融合和注意力机制的环境声识别方法}\dotfill\makebox[2.5em][r]{\wuhao\pageref{paper:I0380}}}
\par
\noindent\hspace*{5em}{\kaishu\wuhao 方子琪, 曾向阳}
\par\vspace{5pt}
\noindent\makebox[5em][l]{\wuhao I0416}\makebox[\dimexpr\linewidth-5em][s]{\hyperlink{paper:I0416}{\wuhao 城市轨道交通噪声控制措施经济评估方法研究}\dotfill\makebox[2.5em][r]{\wuhao\pageref{paper:I0416}}}
\par
\noindent\hspace*{5em}{\kaishu\wuhao 张献英, 侯小晴}
\par\vspace{5pt}
\noindent\makebox[5em][l]{\wuhao I0422}\makebox[\dimexpr\linewidth-5em][s]{\hyperlink{paper:I0422}{\wuhao 校园声景观对大学生心理压力恢复效益的影响研究}\dotfill\makebox[2.5em][r]{\wuhao\pageref{paper:I0422}}}
\par
\noindent\hspace*{5em}{\kaishu\wuhao 单菁菁, 周新欣, 姜兴华, 徐赞}
\par\vspace{5pt}
\noindent\makebox[5em][l]{\wuhao I0432}\makebox[\dimexpr\linewidth-5em][s]{\hyperlink{paper:I0432}{\wuhao 深圳城中村户外声景感知特征及其驱动机制研究}\dotfill\makebox[2.5em][r]{\wuhao\pageref{paper:I0432}}}
\par
\noindent\hspace*{5em}{\kaishu\wuhao 陈珺瑶, 余磊}
\par\vspace{5pt}
\noindent\makebox[5em][l]{\wuhao I0444}\makebox[\dimexpr\linewidth-5em][s]{\hyperlink{paper:I0444}{\wuhao 历史文化遗产声景演变与感知影响研究——以永乐宫为例}\dotfill\makebox[2.5em][r]{\wuhao\pageref{paper:I0444}}}
\par
\noindent\hspace*{5em}{\kaishu\wuhao 晋美俊, 任超冉}
\par\vspace{5pt}
\vspace{6pt}

\begin{center}{\heiti\xiaosi J. 噪声与振动控制}\end{center}
\par\vspace{6pt}
\noindent\hangindent=5em\hangafter=1\makebox[5em][l]{\wuhao J0015}\hyperlink{paper:J0015}{\wuhao Self-reactive, in-situ generated PUA aerogels: Synergistic low-frequency noise attenuation via metamaterial-like structures and dynamically variation of acoustic impedance}
\par\noindent\hspace*{5em}\dotfill\makebox[2.5em][r]{\wuhao\pageref{paper:J0015}}
\par\vspace{2pt}
\noindent\hspace*{5em}{\kaishu\wuhao 郑婉星}
\par\vspace{5pt}
\noindent\makebox[5em][l]{\wuhao J0016}\makebox[\dimexpr\linewidth-5em][s]{\hyperlink{paper:J0016}{\wuhao 基于运动预测驱动的自适应主动降噪方法研究}\dotfill\makebox[2.5em][r]{\wuhao\pageref{paper:J0016}}}
\par
\noindent\hspace*{5em}{\kaishu\wuhao 王博}
\par\vspace{5pt}
\noindent\makebox[5em][l]{\wuhao J0027}\makebox[\dimexpr\linewidth-5em][s]{\hyperlink{paper:J0027}{\wuhao 超声焊接柔顺控制系统的研究}\dotfill\makebox[2.5em][r]{\wuhao\pageref{paper:J0027}}}
\par
\noindent\hspace*{5em}{\kaishu\wuhao 姚震, 符涵, 林超}
\par\vspace{5pt}
\noindent\makebox[5em][l]{\wuhao J0054}\makebox[\dimexpr\linewidth-5em][s]{\hyperlink{paper:J0054}{\wuhao 基于地震学方法的地铁振动信号特征识别与运行轨迹反演研究}\dotfill\makebox[2.5em][r]{\wuhao\pageref{paper:J0054}}}
\par
\noindent\hspace*{5em}{\kaishu\wuhao 张宇翔, 李骐宇, 马丁一, 卞忻然, 王世饶, 边创}
\par\vspace{5pt}
\noindent\makebox[5em][l]{\wuhao J0081}\makebox[\dimexpr\linewidth-5em][s]{\hyperlink{paper:J0081}{\wuhao 基于力作动器的柱壳水下线谱噪声自适应主动控制方法研究}\dotfill\makebox[2.5em][r]{\wuhao\pageref{paper:J0081}}}
\par
\noindent\hspace*{5em}{\kaishu\wuhao 郭星宇, 吴闯, 袁佳伟, 肖妍, 商德江}
\par\vspace{5pt}
\noindent\makebox[5em][l]{\wuhao J0096}\makebox[\dimexpr\linewidth-5em][s]{\hyperlink{paper:J0096}{\wuhao 四模态亥姆霍兹共振超表面的混响声通信优化研究​}\dotfill\makebox[2.5em][r]{\wuhao\pageref{paper:J0096}}}
\par
\noindent\hspace*{5em}{\kaishu\wuhao 刘子粲, 何芝华}
\par\vspace{5pt}
\noindent\hangindent=5em\hangafter=1\makebox[5em][l]{\wuhao J0112}\hyperlink{paper:J0112}{\wuhao 基于粗精 q 移位频率估计的 Fx-CFQSE-LMS 窄带主动噪声控制算法}
\par\noindent\hspace*{5em}\dotfill\makebox[2.5em][r]{\wuhao\pageref{paper:J0112}}
\par\vspace{2pt}
\noindent\hspace*{5em}{\kaishu\wuhao 孙佳艺, 陈静, 高伟, 葛婷婷}
\par\vspace{5pt}
\noindent\makebox[5em][l]{\wuhao J0124}\makebox[\dimexpr\linewidth-5em][s]{\hyperlink{paper:J0124}{\wuhao 一种基于FxLMP的变步长算法}\dotfill\makebox[2.5em][r]{\wuhao\pageref{paper:J0124}}}
\par
\noindent\hspace*{5em}{\kaishu\wuhao 李清雨, 陈静, 高伟, 孙笑}
\par\vspace{5pt}
\noindent\makebox[5em][l]{\wuhao J0141}\makebox[\dimexpr\linewidth-5em][s]{\hyperlink{paper:J0141}{\wuhao 基于局域共振原理的双向可嵌套复合结构宽带声衰减特性研究}\dotfill\makebox[2.5em][r]{\wuhao\pageref{paper:J0141}}}
\par
\noindent\hspace*{5em}{\kaishu\wuhao 陈国帅, 杨金水, 陈蓝宝}
\par\vspace{5pt}
\noindent\makebox[5em][l]{\wuhao J0150}\makebox[\dimexpr\linewidth-5em][s]{\hyperlink{paper:J0150}{\wuhao 基于声源方位识别的前馈有源降噪耳机半自适应控制方法研究}\dotfill\makebox[2.5em][r]{\wuhao\pageref{paper:J0150}}}
\par
\noindent\hspace*{5em}{\kaishu\wuhao 李岩, 安峰岩}
\par\vspace{5pt}
\noindent\makebox[5em][l]{\wuhao J0153}\makebox[\dimexpr\linewidth-5em][s]{\hyperlink{paper:J0153}{\wuhao 深圳市噪声溯源技术研究与示范应用}\dotfill\makebox[2.5em][r]{\wuhao\pageref{paper:J0153}}}
\par
\noindent\hspace*{5em}{\kaishu\wuhao 杨娜, 许沛晨}
\par\vspace{5pt}
\noindent\makebox[5em][l]{\wuhao J0163}\makebox[\dimexpr\linewidth-5em][s]{\hyperlink{paper:J0163}{\wuhao 13677660396}\dotfill\makebox[2.5em][r]{\wuhao\pageref{paper:J0163}}}
\par
\noindent\hspace*{5em}{\kaishu\wuhao 罗磊}
\par\vspace{5pt}
\noindent\makebox[5em][l]{\wuhao J0175}\makebox[\dimexpr\linewidth-5em][s]{\hyperlink{paper:J0175}{\wuhao 受算盘启发的可重构数字超材料用于全偏振振动抑制与逻辑波动控制}\dotfill\makebox[2.5em][r]{\wuhao\pageref{paper:J0175}}}
\par
\noindent\hspace*{5em}{\kaishu\wuhao 李宪铎, 范海燕, 朱一凡, 张辉}
\par\vspace{5pt}
\noindent\makebox[5em][l]{\wuhao J0177}\makebox[\dimexpr\linewidth-5em][s]{\hyperlink{paper:J0177}{\wuhao 使用卷曲共振结构提升封闭箱扬声器低频辐射性能的研究}\dotfill\makebox[2.5em][r]{\wuhao\pageref{paper:J0177}}}
\par
\noindent\hspace*{5em}{\kaishu\wuhao 张宇科, 陈若彦, 陶建成}
\par\vspace{5pt}
\noindent\makebox[5em][l]{\wuhao J0180}\makebox[\dimexpr\linewidth-5em][s]{\hyperlink{paper:J0180}{\wuhao 湍流边界层激励下复合材料层合板随机振动与声辐射的辛波传播分析}\dotfill\makebox[2.5em][r]{\wuhao\pageref{paper:J0180}}}
\par
\noindent\hspace*{5em}{\kaishu\wuhao 严猛, 蔡文}
\par\vspace{5pt}
\noindent\makebox[5em][l]{\wuhao J0190}\makebox[\dimexpr\linewidth-5em][s]{\hyperlink{paper:J0190}{\wuhao 基于主被动复合吸声的二维自由场构建与自适应控制方法研究}\dotfill\makebox[2.5em][r]{\wuhao\pageref{paper:J0190}}}
\par
\noindent\hspace*{5em}{\kaishu\wuhao 赵立赢, 安峰岩}
\par\vspace{5pt}
\noindent\makebox[5em][l]{\wuhao J0191}\makebox[\dimexpr\linewidth-5em][s]{\hyperlink{paper:J0191}{\wuhao 基于声强目标函数的管路有源噪声控制方法研究​}\dotfill\makebox[2.5em][r]{\wuhao\pageref{paper:J0191}}}
\par
\noindent\hspace*{5em}{\kaishu\wuhao 孟令辉, 安峰岩}
\par\vspace{5pt}
\noindent\makebox[5em][l]{\wuhao J0198}\makebox[\dimexpr\linewidth-5em][s]{\hyperlink{paper:J0198}{\wuhao 用于全向宽带隔声的通风超表面}\dotfill\makebox[2.5em][r]{\wuhao\pageref{paper:J0198}}}
\par
\noindent\hspace*{5em}{\kaishu\wuhao 朱一寰, 毛东兴, 王旭, 李勇}
\par\vspace{5pt}
\noindent\makebox[5em][l]{\wuhao J0203}\makebox[\dimexpr\linewidth-5em][s]{\hyperlink{paper:J0203}{\wuhao 基于颗粒协同运动破坏机制的高效分层颗粒阻尼器}\dotfill\makebox[2.5em][r]{\wuhao\pageref{paper:J0203}}}
\par
\noindent\hspace*{5em}{\kaishu\wuhao 赵圣强, 刘京京, 梁彬, 程建春}
\par\vspace{5pt}
\noindent\makebox[5em][l]{\wuhao J0227}\makebox[\dimexpr\linewidth-5em][s]{\hyperlink{paper:J0227}{\wuhao 基于头枕传声器的远程传声器法在人头平移时的声压估计误差}\dotfill\makebox[2.5em][r]{\wuhao\pageref{paper:J0227}}}
\par
\noindent\hspace*{5em}{\kaishu\wuhao 吴天昊, 陶建成}
\par\vspace{5pt}
\noindent\makebox[5em][l]{\wuhao J0234}\makebox[\dimexpr\linewidth-5em][s]{\hyperlink{paper:J0234}{\wuhao 汽车路噪有源降噪中面向主观感知的代价函数设计}\dotfill\makebox[2.5em][r]{\wuhao\pageref{paper:J0234}}}
\par
\noindent\hspace*{5em}{\kaishu\wuhao 陈栋宇, 匡家旗, 陶建成}
\par\vspace{5pt}
\noindent\makebox[5em][l]{\wuhao J0292}\makebox[\dimexpr\linewidth-5em][s]{\hyperlink{paper:J0292}{\wuhao 道路固废制备吸隔声声屏障材料技术研究}\dotfill\makebox[2.5em][r]{\wuhao\pageref{paper:J0292}}}
\par
\noindent\hspace*{5em}{\kaishu\wuhao 路可欣, 袁旻忞}
\par\vspace{5pt}
\noindent\makebox[5em][l]{\wuhao J0296}\makebox[\dimexpr\linewidth-5em][s]{\hyperlink{paper:J0296}{\wuhao 小尺寸低频宽带微穿孔板复合吸声体设计与研究}\dotfill\makebox[2.5em][r]{\wuhao\pageref{paper:J0296}}}
\par
\noindent\hspace*{5em}{\kaishu\wuhao 方闯闯, 刘秀娟}
\par\vspace{5pt}
\noindent\makebox[5em][l]{\wuhao J0364}\makebox[\dimexpr\linewidth-5em][s]{\hyperlink{paper:J0364}{\wuhao 一种连续可调的宽带通风隔音屏障}\dotfill\makebox[2.5em][r]{\wuhao\pageref{paper:J0364}}}
\par
\noindent\hspace*{5em}{\kaishu\wuhao 梁善军}
\par\vspace{5pt}
\noindent\makebox[5em][l]{\wuhao J0402}\makebox[\dimexpr\linewidth-5em][s]{\hyperlink{paper:J0402}{\wuhao 声学黑洞效应在基阵平台中的减振应用前景}\dotfill\makebox[2.5em][r]{\wuhao\pageref{paper:J0402}}}
\par
\noindent\hspace*{5em}{\kaishu\wuhao 罗栋栋, 葛晨昊, 楼万翔, 祝献}
\par\vspace{5pt}
\noindent\makebox[5em][l]{\wuhao J0407}\makebox[\dimexpr\linewidth-5em][s]{\hyperlink{paper:J0407}{\wuhao 吸附效应下突破因果最小厚度约束的超宽带声吸收}\dotfill\makebox[2.5em][r]{\wuhao\pageref{paper:J0407}}}
\par
\noindent\hspace*{5em}{\kaishu\wuhao 梅中建, 葛传昊, 石同阳, 李勇, 杨军}
\par\vspace{5pt}
\noindent\makebox[5em][l]{\wuhao J0450}\makebox[\dimexpr\linewidth-5em][s]{\hyperlink{paper:J0450}{\wuhao 含多孔夹层的三层圆柱壳弹性波传播特性研究}\dotfill\makebox[2.5em][r]{\wuhao\pageref{paper:J0450}}}
\par
\noindent\hspace*{5em}{\kaishu\wuhao 丁星月, 米泽宁, 石同阳}
\par\vspace{5pt}
\noindent\makebox[5em][l]{\wuhao J0456}\makebox[\dimexpr\linewidth-5em][s]{\hyperlink{paper:J0456}{\wuhao 基于SONAH的混合传感器不规则阵列近场声全息方法}\dotfill\makebox[2.5em][r]{\wuhao\pageref{paper:J0456}}}
\par
\noindent\hspace*{5em}{\kaishu\wuhao 谭盈, 崔贞鸣, 石同阳}
\par\vspace{5pt}
\vspace{6pt}

\begin{center}{\heiti\xiaosi K. 通信声学与音频信号处理（含声频工程）}\end{center}
\par\vspace{6pt}
\noindent\makebox[5em][l]{\wuhao K0113}\makebox[\dimexpr\linewidth-5em][s]{\hyperlink{paper:K0113}{\wuhao 基于双通道串扰消除系统的最优控制点位置选择}\dotfill\makebox[2.5em][r]{\wuhao\pageref{paper:K0113}}}
\par
\noindent\hspace*{5em}{\kaishu\wuhao 黄洋, 王雪珊, 赵斯培, 邹海山, 卢晶}
\par\vspace{5pt}
\noindent\makebox[5em][l]{\wuhao K0117}\makebox[\dimexpr\linewidth-5em][s]{\hyperlink{paper:K0117}{\wuhao 基于频带感知的噪声环境电力变压器声学异常检测}\dotfill\makebox[2.5em][r]{\wuhao\pageref{paper:K0117}}}
\par
\noindent\hspace*{5em}{\kaishu\wuhao 杨俊采}
\par\vspace{5pt}
\noindent\makebox[5em][l]{\wuhao K0135}\makebox[\dimexpr\linewidth-5em][s]{\hyperlink{paper:K0135}{\wuhao FSC-Net：融合快速傅里叶卷积与渐进式学习的语音带宽扩展网络}\dotfill\makebox[2.5em][r]{\wuhao\pageref{paper:K0135}}}
\par
\noindent\hspace*{5em}{\kaishu\wuhao 陈欣安, 容晓彬, 胡沁雯, 卢晶, 陈锴}
\par\vspace{5pt}
\noindent\makebox[5em][l]{\wuhao K0178}\makebox[\dimexpr\linewidth-5em][s]{\hyperlink{paper:K0178}{\wuhao 涡旋参量阵扬声器}\dotfill\makebox[2.5em][r]{\wuhao\pageref{paper:K0178}}}
\par
\noindent\hspace*{5em}{\kaishu\wuhao 李绍哲, 钟家鑫, 卢晶, Yun Jing}
\par\vspace{5pt}
\noindent\makebox[5em][l]{\wuhao K0249}\makebox[\dimexpr\linewidth-5em][s]{\hyperlink{paper:K0249}{\wuhao 刚性障板圆形扬声器阵列的差分波束形成方法}\dotfill\makebox[2.5em][r]{\wuhao\pageref{paper:K0249}}}
\par
\noindent\hspace*{5em}{\kaishu\wuhao 张燕凯}
\par\vspace{5pt}
\noindent\makebox[5em][l]{\wuhao K0264}\makebox[\dimexpr\linewidth-5em][s]{\hyperlink{paper:K0264}{\wuhao 一种基于CNN-BiLSTM-Attention混合架构的乐器识别方法}\dotfill\makebox[2.5em][r]{\wuhao\pageref{paper:K0264}}}
\par
\noindent\hspace*{5em}{\kaishu\wuhao 王喆, 陈奕涵, 陈静}
\par\vspace{5pt}
\noindent\makebox[5em][l]{\wuhao K0272}\makebox[\dimexpr\linewidth-5em][s]{\hyperlink{paper:K0272}{\wuhao 相控参量阵扬声器阵列的波束调控研究}\dotfill\makebox[2.5em][r]{\wuhao\pageref{paper:K0272}}}
\par
\noindent\hspace*{5em}{\kaishu\wuhao 吴旭祥, 李绍哲, 李梦同, 钟家鑫, 卢晶}
\par\vspace{5pt}
\noindent\makebox[5em][l]{\wuhao K0305}\makebox[\dimexpr\linewidth-5em][s]{\hyperlink{paper:K0305}{\wuhao 面向水声通信组网的码分多址接入干扰抵消技术}\dotfill\makebox[2.5em][r]{\wuhao\pageref{paper:K0305}}}
\par
\noindent\hspace*{5em}{\kaishu\wuhao 温梦华, 熊省军, 李栋}
\par\vspace{5pt}
\noindent\makebox[5em][l]{\wuhao K0370}\makebox[\dimexpr\linewidth-5em][s]{\hyperlink{paper:K0370}{\wuhao 基于训练序列相位匹配的水声单载波通信系统残余相偏消除方法}\dotfill\makebox[2.5em][r]{\wuhao\pageref{paper:K0370}}}
\par
\noindent\hspace*{5em}{\kaishu\wuhao 谢哲, 王桢铎}
\par\vspace{5pt}
\noindent\makebox[5em][l]{\wuhao K0381}\makebox[\dimexpr\linewidth-5em][s]{\hyperlink{paper:K0381}{\wuhao 面向助听器声反馈消除的深度步长控制与去相关策略}\dotfill\makebox[2.5em][r]{\wuhao\pageref{paper:K0381}}}
\par
\noindent\hspace*{5em}{\kaishu\wuhao 詹小凡, 郝逢源, 李晓东, 郑成诗}
\par\vspace{5pt}
\noindent\makebox[5em][l]{\wuhao K0400}\makebox[\dimexpr\linewidth-5em][s]{\hyperlink{paper:K0400}{\wuhao 复杂量测环境下的改进 SMC-PHD 多目标跟踪算法}\dotfill\makebox[2.5em][r]{\wuhao\pageref{paper:K0400}}}
\par
\noindent\hspace*{5em}{\kaishu\wuhao 韩帅伊, 王嘉伟, 张舒杨, 胡小青}
\par\vspace{5pt}
\noindent\makebox[5em][l]{\wuhao K0453}\makebox[\dimexpr\linewidth-5em][s]{\hyperlink{paper:K0453}{\wuhao 强干扰下基于稀疏表示的声源分离与定位方法}\dotfill\makebox[2.5em][r]{\wuhao\pageref{paper:K0453}}}
\par
\noindent\hspace*{5em}{\kaishu\wuhao 崔贞鸣, 石同阳, 吴鸣, 杨军}
\par\vspace{5pt}
\vspace{6pt}

\begin{center}{\heiti\xiaosi L.语言声学与语音信号处理}\end{center}
\par\vspace{6pt}
\noindent\makebox[5em][l]{\wuhao L0115}\makebox[\dimexpr\linewidth-5em][s]{\hyperlink{paper:L0115}{\wuhao 基于音节融合模型的连读语音识别研究}\dotfill\makebox[2.5em][r]{\wuhao\pageref{paper:L0115}}}
\par
\noindent\hspace*{5em}{\kaishu\wuhao 廖玮}
\par\vspace{5pt}
\noindent\makebox[5em][l]{\wuhao L0123}\makebox[\dimexpr\linewidth-5em][s]{\hyperlink{paper:L0123}{\wuhao 基于多尺度自监督学习的目标说话人识别方法研究​}\dotfill\makebox[2.5em][r]{\wuhao\pageref{paper:L0123}}}
\par
\noindent\hspace*{5em}{\kaishu\wuhao 还好, 林玮}
\par\vspace{5pt}
\noindent\makebox[5em][l]{\wuhao L0240}\makebox[\dimexpr\linewidth-5em][s]{\hyperlink{paper:L0240}{\wuhao 基于状态空间模型引导残差扩散的实时语音增强方法研究}\dotfill\makebox[2.5em][r]{\wuhao\pageref{paper:L0240}}}
\par
\noindent\hspace*{5em}{\kaishu\wuhao 冯帅铧, 林玮}
\par\vspace{5pt}
\noindent\makebox[5em][l]{\wuhao L0320}\makebox[\dimexpr\linewidth-5em][s]{\hyperlink{paper:L0320}{\wuhao 融合小波变换的骨导语音盲增强方法}\dotfill\makebox[2.5em][r]{\wuhao\pageref{paper:L0320}}}
\par
\noindent\hspace*{5em}{\kaishu\wuhao 岳昭新, 郭剑锋}
\par\vspace{5pt}
\noindent\makebox[5em][l]{\wuhao L0336}\makebox[\dimexpr\linewidth-5em][s]{\hyperlink{paper:L0336}{\wuhao 物理启发的频域约束与幅-相解耦伪标签学习语音增强方法}\dotfill\makebox[2.5em][r]{\wuhao\pageref{paper:L0336}}}
\par
\noindent\hspace*{5em}{\kaishu\wuhao 李扬}
\par\vspace{5pt}
\vspace{6pt}

\begin{center}{\heiti\xiaosi M. 结构与建筑声学}\end{center}
\par\vspace{6pt}
\noindent\hangindent=5em\hangafter=1\makebox[5em][l]{\wuhao M0019}\hyperlink{paper:M0019}{\wuhao Machine Learning-Driven Prediction and Reverse Design Optimization of Acoustic Properties in Nonwoven Materials}
\par\noindent\hspace*{5em}\dotfill\makebox[2.5em][r]{\wuhao\pageref{paper:M0019}}
\par\vspace{2pt}
\noindent\hspace*{5em}{\kaishu\wuhao 陈猛, 陈江龙, 支超}
\par\vspace{5pt}
\noindent\makebox[5em][l]{\wuhao M0065}\makebox[\dimexpr\linewidth-5em][s]{\hyperlink{paper:M0065}{\wuhao 大众羽毛球馆的声场特性与言语可懂度研究}\dotfill\makebox[2.5em][r]{\wuhao\pageref{paper:M0065}}}
\par
\noindent\hspace*{5em}{\kaishu\wuhao 李嘉奇, 梁林达}
\par\vspace{5pt}
\noindent\makebox[5em][l]{\wuhao M0072}\makebox[\dimexpr\linewidth-5em][s]{\hyperlink{paper:M0072}{\wuhao 工人视角下工业建筑健康声环境评价研究：混合方法}\dotfill\makebox[2.5em][r]{\wuhao\pageref{paper:M0072}}}
\par
\noindent\hspace*{5em}{\kaishu\wuhao 张宇轩, 秦金辉}
\par\vspace{5pt}
\noindent\makebox[5em][l]{\wuhao M0105}\makebox[\dimexpr\linewidth-5em][s]{\hyperlink{paper:M0105}{\wuhao 恒定吸声量地下大厅音质设计研究}\dotfill\makebox[2.5em][r]{\wuhao\pageref{paper:M0105}}}
\par
\noindent\hspace*{5em}{\kaishu\wuhao 卢亦林, 燕翔, 李卉, 杜春雨, 姜建中, 蔡浩}
\par\vspace{5pt}
\noindent\makebox[5em][l]{\wuhao M0130}\makebox[\dimexpr\linewidth-5em][s]{\hyperlink{paper:M0130}{\wuhao 高阻尼夹胶玻璃的均质化建模与宽频减振特性分析}\dotfill\makebox[2.5em][r]{\wuhao\pageref{paper:M0130}}}
\par
\noindent\hspace*{5em}{\kaishu\wuhao 王博, 程可欣}
\par\vspace{5pt}
\noindent\hangindent=5em\hangafter=1\makebox[5em][l]{\wuhao M0168}\hyperlink{paper:M0168}{\wuhao 基于S-O-R模型的地铁站域视听环境对商业空间活力链式传导路径研究}
\par\noindent\hspace*{5em}\dotfill\makebox[2.5em][r]{\wuhao\pageref{paper:M0168}}
\par\vspace{2pt}
\noindent\hspace*{5em}{\kaishu\wuhao 胡赵敏, 毛琳箐, 角楚云, 贾怡红, 宋恒玲}
\par\vspace{5pt}
\noindent\makebox[5em][l]{\wuhao M0179}\makebox[\dimexpr\linewidth-5em][s]{\hyperlink{paper:M0179}{\wuhao 可听化声景实验在大型商业空间混响时间主观评价中的应用}\dotfill\makebox[2.5em][r]{\wuhao\pageref{paper:M0179}}}
\par
\noindent\hspace*{5em}{\kaishu\wuhao 杜春雨, 王鹏, 燕翔, 卢亦林, 李卉}
\par\vspace{5pt}
\noindent\makebox[5em][l]{\wuhao M0186}\makebox[\dimexpr\linewidth-5em][s]{\hyperlink{paper:M0186}{\wuhao 可调吸声构造设计及其在可变混响听音室中的应用}\dotfill\makebox[2.5em][r]{\wuhao\pageref{paper:M0186}}}
\par
\noindent\hspace*{5em}{\kaishu\wuhao 查颖杰, 毛东兴}
\par\vspace{5pt}
\noindent\makebox[5em][l]{\wuhao M0230}\makebox[\dimexpr\linewidth-5em][s]{\hyperlink{paper:M0230}{\wuhao 关于浮筑楼板弹性材料的动刚度实验研究}\dotfill\makebox[2.5em][r]{\wuhao\pageref{paper:M0230}}}
\par
\noindent\hspace*{5em}{\kaishu\wuhao 范洪旺}
\par\vspace{5pt}
\noindent\makebox[5em][l]{\wuhao M0253}\makebox[\dimexpr\linewidth-5em][s]{\hyperlink{paper:M0253}{\wuhao 基于问卷调查的住宅夜间室内噪声源及睡眠声环境现状}\dotfill\makebox[2.5em][r]{\wuhao\pageref{paper:M0253}}}
\par
\noindent\hspace*{5em}{\kaishu\wuhao 霍泽宇, 周杨坤, 武宝华, 陈静, 蔡超}
\par\vspace{5pt}
\noindent\makebox[5em][l]{\wuhao M0300}\makebox[\dimexpr\linewidth-5em][s]{\hyperlink{paper:M0300}{\wuhao 宁海古戏台历史声景复原与保护研究}\dotfill\makebox[2.5em][r]{\wuhao\pageref{paper:M0300}}}
\par
\noindent\hspace*{5em}{\kaishu\wuhao 黄姗, 冯怡涵}
\par\vspace{5pt}
\noindent\makebox[5em][l]{\wuhao M0333}\makebox[\dimexpr\linewidth-5em][s]{\hyperlink{paper:M0333}{\wuhao 高层住宅空气声与撞击声隔声性能及垂向传播规律实测研究}\dotfill\makebox[2.5em][r]{\wuhao\pageref{paper:M0333}}}
\par
\noindent\hspace*{5em}{\kaishu\wuhao 乔木薇, 宋阳}
\par\vspace{5pt}
\noindent\makebox[5em][l]{\wuhao M0338}\makebox[\dimexpr\linewidth-5em][s]{\hyperlink{paper:M0338}{\wuhao 地铁沿线居住区室内二次辐射噪声感知特征研究——以石家庄市为例}\dotfill\makebox[2.5em][r]{\wuhao\pageref{paper:M0338}}}
\par
\noindent\hspace*{5em}{\kaishu\wuhao 角楚云, 毛琳箐, 胡赵敏, 贾怡红, 宋恒玲}
\par\vspace{5pt}
\noindent\makebox[5em][l]{\wuhao M0362}\makebox[\dimexpr\linewidth-5em][s]{\hyperlink{paper:M0362}{\wuhao 社区水景场景参数、声学特征与老年人恢复效应的关联研究}\dotfill\makebox[2.5em][r]{\wuhao\pageref{paper:M0362}}}
\par
\noindent\hspace*{5em}{\kaishu\wuhao 张心茹, 龙晓婕, 雷雨亮, Nazli Che Din}
\par\vspace{5pt}
\noindent\makebox[5em][l]{\wuhao M0445}\makebox[\dimexpr\linewidth-5em][s]{\hyperlink{paper:M0445}{\wuhao 噪声对睡眠干扰核心影响因素：最大声压级还是相对声压级？}\dotfill\makebox[2.5em][r]{\wuhao\pageref{paper:M0445}}}
\par
\noindent\hspace*{5em}{\kaishu\wuhao 陈静, 霍泽宇, 武宝华, 周杨坤, 蔡超}
\par\vspace{5pt}
\noindent\hangindent=5em\hangafter=1\makebox[5em][l]{\wuhao M0457}\hyperlink{paper:M0457}{\wuhao 加气混凝土砌块墙体隔声性能浅析摘要：本文从现行国家规范对隔墙的隔声要求出发，分析加气混凝土砌块隔墙隔声量的影响因素，在实验室测试6种不同厚度与抹灰工况的B06加气混凝土砌块墙体的隔声量。结果说明了不同厚度砌块墙体的隔声量，不同抹灰条件对墙体隔声量的影响情况。200mm厚B06砌块墙仍无法满足新国标分户墙隔声要求。本文给出了满足现行标准与宁静住宅得分项的墙体构造方案。}
\par\noindent\hspace*{5em}\dotfill\makebox[2.5em][r]{\wuhao\pageref{paper:M0457}}
\par\vspace{2pt}
\noindent\hspace*{5em}{\kaishu\wuhao 金迪锋}
\par\vspace{5pt}
\vspace{6pt}

\begin{center}{\heiti\xiaosi O. 机械振动与冲击}\end{center}
\par\vspace{6pt}
\noindent\makebox[5em][l]{\wuhao O0028}\makebox[\dimexpr\linewidth-5em][s]{\hyperlink{paper:O0028}{\wuhao U型波纹夹芯板入水冲击载荷及动态响应数值研究}\dotfill\makebox[2.5em][r]{\wuhao\pageref{paper:O0028}}}
\par
\noindent\hspace*{5em}{\kaishu\wuhao 廖永, 胡勇军, 陈昭宇}
\par\vspace{5pt}
\noindent\makebox[5em][l]{\wuhao O0110}\makebox[\dimexpr\linewidth-5em][s]{\hyperlink{paper:O0110}{\wuhao 集成 PZT 作动器的空间索网反射面天线主动形面调整与振动控制}\dotfill\makebox[2.5em][r]{\wuhao\pageref{paper:O0110}}}
\par
\noindent\hspace*{5em}{\kaishu\wuhao 张知行}
\par\vspace{5pt}
\noindent\makebox[5em][l]{\wuhao O0120}\makebox[\dimexpr\linewidth-5em][s]{\hyperlink{paper:O0120}{\wuhao 基于神经网络代理模型的管路约束参数反演及位置优化}\dotfill\makebox[2.5em][r]{\wuhao\pageref{paper:O0120}}}
\par
\noindent\hspace*{5em}{\kaishu\wuhao 彭科, 赵晓臣, 张新玉, 柳贡民}
\par\vspace{5pt}
\noindent\hangindent=5em\hangafter=1\makebox[5em][l]{\wuhao O0382}\hyperlink{paper:O0382}{\wuhao 针对超声振动辅助加工的声表面波多模态深度融合及加工状态预测系统}
\par\noindent\hspace*{5em}\dotfill\makebox[2.5em][r]{\wuhao\pageref{paper:O0382}}
\par\vspace{2pt}
\noindent\hspace*{5em}{\kaishu\wuhao 邵逸之, 龚涛, 李思慧, 孙可}
\par\vspace{5pt}
\noindent\makebox[5em][l]{\wuhao O0415}\makebox[\dimexpr\linewidth-5em][s]{\hyperlink{paper:O0415}{\wuhao 纯电动汽车空气压缩机系统振动噪声性能改进}\dotfill\makebox[2.5em][r]{\wuhao\pageref{paper:O0415}}}
\par
\noindent\hspace*{5em}{\kaishu\wuhao 程志伟}
\par\vspace{5pt}
\vspace{6pt}

\begin{center}{\heiti\xiaosi P.生理声学、心理声学、音乐声学}\end{center}
\par\vspace{6pt}
\noindent\makebox[5em][l]{\wuhao P0082}\makebox[\dimexpr\linewidth-5em][s]{\hyperlink{paper:P0082}{\wuhao 基于音频特征的音乐情感识别算法研究}\dotfill\makebox[2.5em][r]{\wuhao\pageref{paper:P0082}}}
\par
\noindent\hspace*{5em}{\kaishu\wuhao 何宇扬, 林玮}
\par\vspace{5pt}
\noindent\makebox[5em][l]{\wuhao P0187}\makebox[\dimexpr\linewidth-5em][s]{\hyperlink{paper:P0187}{\wuhao 助听器使用对汉语辅音语音特征信息传递的影响研究}\dotfill\makebox[2.5em][r]{\wuhao\pageref{paper:P0187}}}
\par
\noindent\hspace*{5em}{\kaishu\wuhao 周华莉, 贝先明, 郭振宇, 曾宪海, 魏朝刚, 郑能恒, 李娟娟, 孟庆林}
\par\vspace{5pt}
\noindent\hangindent=5em\hangafter=1\makebox[5em][l]{\wuhao P0220}\hyperlink{paper:P0220}{\wuhao Sottek听觉模型驱动的心理声学参数标准化体系搭建与人工智能应用实践}
\par\noindent\hspace*{5em}\dotfill\makebox[2.5em][r]{\wuhao\pageref{paper:P0220}}
\par\vspace{2pt}
\noindent\hspace*{5em}{\kaishu\wuhao 毛阳, 禹丹江}
\par\vspace{5pt}
\noindent\makebox[5em][l]{\wuhao P0257}\makebox[\dimexpr\linewidth-5em][s]{\hyperlink{paper:P0257}{\wuhao 基于声学物理特征的双目标听觉感知脑电机制研究}\dotfill\makebox[2.5em][r]{\wuhao\pageref{paper:P0257}}}
\par
\noindent\hspace*{5em}{\kaishu\wuhao 刘诗韵, 王嘉伟, 李忠杰, 甘霖}
\par\vspace{5pt}
\noindent\makebox[5em][l]{\wuhao P0274}\makebox[\dimexpr\linewidth-5em][s]{\hyperlink{paper:P0274}{\wuhao 基于“音色-响度”交互特征的听觉响度感知层级化脑电表征研究}\dotfill\makebox[2.5em][r]{\wuhao\pageref{paper:P0274}}}
\par
\noindent\hspace*{5em}{\kaishu\wuhao 周嘉诚, 黄俊伟, 李忠杰, 陈思宇, 甘霖}
\par\vspace{5pt}
\noindent\hangindent=5em\hangafter=1\makebox[5em][l]{\wuhao P0276}\hyperlink{paper:P0276}{\wuhao 从噪声控制到生成式设计：基于机器学习文献计量学绘制声景研究的演进脉络}
\par\noindent\hspace*{5em}\dotfill\makebox[2.5em][r]{\wuhao\pageref{paper:P0276}}
\par\vspace{2pt}
\noindent\hspace*{5em}{\kaishu\wuhao 陈骏博, 屠娟}
\par\vspace{5pt}
\noindent\makebox[5em][l]{\wuhao P0278}\makebox[\dimexpr\linewidth-5em][s]{\hyperlink{paper:P0278}{\wuhao 基于脑电响应的音乐调性与和弦感知的神经机制研究}\dotfill\makebox[2.5em][r]{\wuhao\pageref{paper:P0278}}}
\par
\noindent\hspace*{5em}{\kaishu\wuhao 王嘉伟, 李忠杰, 秦川, 赵政越, 钱丞煜, 甘霖}
\par\vspace{5pt}
\noindent\makebox[5em][l]{\wuhao P0301}\makebox[\dimexpr\linewidth-5em][s]{\hyperlink{paper:P0301}{\wuhao 基于数字听诊的人体肠鸣音声学特征分析与分级评价指标体系构建}\dotfill\makebox[2.5em][r]{\wuhao\pageref{paper:P0301}}}
\par
\noindent\hspace*{5em}{\kaishu\wuhao 方宇辰, 梁雍, 闫靓}
\par\vspace{5pt}
\noindent\makebox[5em][l]{\wuhao P0341}\makebox[\dimexpr\linewidth-5em][s]{\hyperlink{paper:P0341}{\wuhao 基于驾驶环境噪声与声学干预的疲劳状态检测与调控研究}\dotfill\makebox[2.5em][r]{\wuhao\pageref{paper:P0341}}}
\par
\noindent\hspace*{5em}{\kaishu\wuhao 赖展翔, 李忠杰, 郝恬然, 曹浩然, 甘霖}
\par\vspace{5pt}
\noindent\makebox[5em][l]{\wuhao P0342}\makebox[\dimexpr\linewidth-5em][s]{\hyperlink{paper:P0342}{\wuhao 基于音乐听觉干预的精神疲劳缓解研究}\dotfill\makebox[2.5em][r]{\wuhao\pageref{paper:P0342}}}
\par
\noindent\hspace*{5em}{\kaishu\wuhao 上官世宝, 李忠杰, 牛超文, 丁祥祥, 秦艺卓, 甘霖}
\par\vspace{5pt}
\noindent\makebox[5em][l]{\wuhao P0345}\makebox[\dimexpr\linewidth-5em][s]{\hyperlink{paper:P0345}{\wuhao 基于脑电信号源定位技术的调性音乐特征诱发机制研究}\dotfill\makebox[2.5em][r]{\wuhao\pageref{paper:P0345}}}
\par
\noindent\hspace*{5em}{\kaishu\wuhao 周璟轩, 黄艳, 李忠杰, 甘霖}
\par\vspace{5pt}
\noindent\makebox[5em][l]{\wuhao P0375}\makebox[\dimexpr\linewidth-5em][s]{\hyperlink{paper:P0375}{\wuhao 基于头戴式传声器阵列的不同距离双耳重建算法感知评估}\dotfill\makebox[2.5em][r]{\wuhao\pageref{paper:P0375}}}
\par
\noindent\hspace*{5em}{\kaishu\wuhao 李笑宇, 柯雨璇, 蔡娟娟, 吴彦琴, 郑成诗}
\par\vspace{5pt}
\noindent\hangindent=5em\hangafter=1\makebox[5em][l]{\wuhao P0383}\hyperlink{paper:P0383}{\wuhao 基于GAM-神经网络混合模型的全人群感音神经性听力损失听阈轨迹预测研究}
\par\noindent\hspace*{5em}\dotfill\makebox[2.5em][r]{\wuhao\pageref{paper:P0383}}
\par\vspace{2pt}
\noindent\hspace*{5em}{\kaishu\wuhao 冀飞, 韩璐檬}
\par\vspace{5pt}
\noindent\makebox[5em][l]{\wuhao P0391}\makebox[\dimexpr\linewidth-5em][s]{\hyperlink{paper:P0391}{\wuhao 年龄相关性听力损失的多影响因素分析}\dotfill\makebox[2.5em][r]{\wuhao\pageref{paper:P0391}}}
\par
\noindent\hspace*{5em}{\kaishu\wuhao 冀飞, 韩璐檬}
\par\vspace{5pt}
\noindent\makebox[5em][l]{\wuhao P0405}\makebox[\dimexpr\linewidth-5em][s]{\hyperlink{paper:P0405}{\wuhao 基于被动声刺激的昏迷状态脑认知能力分析}\dotfill\makebox[2.5em][r]{\wuhao\pageref{paper:P0405}}}
\par
\noindent\hspace*{5em}{\kaishu\wuhao 焦皓冉, 柴鹏, 李忠杰, 甘霖}
\par\vspace{5pt}
\noindent\makebox[5em][l]{\wuhao P0409}\makebox[\dimexpr\linewidth-5em][s]{\hyperlink{paper:P0409}{\wuhao 基于声景大数据的汉传佛教寺庙音乐审美认知研究}\dotfill\makebox[2.5em][r]{\wuhao\pageref{paper:P0409}}}
\par
\noindent\hspace*{5em}{\kaishu\wuhao 李昀昊, 任欣欣}
\par\vspace{5pt}
\noindent\makebox[5em][l]{\wuhao P0433}\makebox[\dimexpr\linewidth-5em][s]{\hyperlink{paper:P0433}{\wuhao 自动纯音测听中受试者反应时间的声级与频率依赖性研究}\dotfill\makebox[2.5em][r]{\wuhao\pageref{paper:P0433}}}
\par
\noindent\hspace*{5em}{\kaishu\wuhao 郭振宇, 李秀霞, 周华莉, 李娟娟, 冀飞, 孟庆林}
\par\vspace{5pt}
\noindent\makebox[5em][l]{\wuhao P0452}\makebox[\dimexpr\linewidth-5em][s]{\hyperlink{paper:P0452}{\wuhao 一种基于双时间尺度控制的自适应复合主动声品质控制架构}\dotfill\makebox[2.5em][r]{\wuhao\pageref{paper:P0452}}}
\par
\noindent\hspace*{5em}{\kaishu\wuhao 罗剑峰, 陈克安}
\par\vspace{5pt}
\vspace{6pt}

\begin{center}{\heiti\xiaosi Q. 声学测量与仪器}\end{center}
\par\vspace{6pt}
\noindent\makebox[5em][l]{\wuhao Q0035}\makebox[\dimexpr\linewidth-5em][s]{\hyperlink{paper:Q0035}{\wuhao 基于压电嵌套盘环结构的高灵敏度双谐振宽频带声发射传感器设计}\dotfill\makebox[2.5em][r]{\wuhao\pageref{paper:Q0035}}}
\par
\noindent\hspace*{5em}{\kaishu\wuhao 韩庆, 卿小凤, 熊堂虹, 钟泽锋, 魏贤华}
\par\vspace{5pt}
\noindent\hangindent=5em\hangafter=1\makebox[5em][l]{\wuhao Q0095}\hyperlink{paper:Q0095}{\wuhao 扩散声场阵列合成通常依赖扬声器与控制点之间传递矩阵的准确获取。然而，通道数量的增加往往伴随着较高的测量代价。由于实际扬声器的辐射指向性通常较为复杂，单一声中心模型难以准确表征声源在不同空间方向上的局部传播特征。针对上述问题，本文提出一种二维分区声中心模型及其传递矩阵预测方法。首先，将目标合成区域划分为若干局部子区域，并在各子区域内引入局部等效声中心表示。其次，分别建立幅值型与相位型局部声中心模型，以刻画局部传递函数的幅值和相位特性。在声中心参数估计过程中，利用全区域已选控制点构建全局参考参数，并结合归一化正则化非线性最小二乘方法，实现局部声中心参数的稳健估计。最后，基于所估计的参数预测目标区域内各控制点的传递函数，并组装得到预测传递矩阵。结果表明，所提方法在减少测量点数量的条件下，仍能较准确地重构传递矩阵，并保持良好的扩散声场合成性能。}
\par\noindent\hspace*{5em}\dotfill\makebox[2.5em][r]{\wuhao\pageref{paper:Q0095}}
\par\vspace{2pt}
\noindent\hspace*{5em}{\kaishu\wuhao 刘好胜, 刘碧龙, 秦超}
\par\vspace{5pt}
\noindent\makebox[5em][l]{\wuhao Q0181}\makebox[\dimexpr\linewidth-5em][s]{\hyperlink{paper:Q0181}{\wuhao 基于真耳内置法的护听器个体适应性测量系统研究}\dotfill\makebox[2.5em][r]{\wuhao\pageref{paper:Q0181}}}
\par
\noindent\hspace*{5em}{\kaishu\wuhao 张永军}
\par\vspace{5pt}
\noindent\makebox[5em][l]{\wuhao Q0202}\makebox[\dimexpr\linewidth-5em][s]{\hyperlink{paper:Q0202}{\wuhao 一种基于同源标志触发信号的数字水听器指向性测量方法}\dotfill\makebox[2.5em][r]{\wuhao\pageref{paper:Q0202}}}
\par
\noindent\hspace*{5em}{\kaishu\wuhao 曹猛, 佟昊阳, 杨柳青, 李文静}
\par\vspace{5pt}
\noindent\makebox[5em][l]{\wuhao Q0233}\makebox[\dimexpr\linewidth-5em][s]{\hyperlink{paper:Q0233}{\wuhao 机场噪声测量技术综述}\dotfill\makebox[2.5em][r]{\wuhao\pageref{paper:Q0233}}}
\par
\noindent\hspace*{5em}{\kaishu\wuhao 熊文波}
\par\vspace{5pt}
\noindent\makebox[5em][l]{\wuhao Q0250}\makebox[\dimexpr\linewidth-5em][s]{\hyperlink{paper:Q0250}{\wuhao 非侵入式测量管道内壁流体流速与压力的流量计}\dotfill\makebox[2.5em][r]{\wuhao\pageref{paper:Q0250}}}
\par
\noindent\hspace*{5em}{\kaishu\wuhao 薛若珂}
\par\vspace{5pt}
\noindent\makebox[5em][l]{\wuhao Q0293}\makebox[\dimexpr\linewidth-5em][s]{\hyperlink{paper:Q0293}{\wuhao 面向复杂工况的超声线阵换能器液面形态检测方法研究}\dotfill\makebox[2.5em][r]{\wuhao\pageref{paper:Q0293}}}
\par
\noindent\hspace*{5em}{\kaishu\wuhao 高琬佳, 何常德, 戴文舒}
\par\vspace{5pt}
\noindent\makebox[5em][l]{\wuhao Q0356}\makebox[\dimexpr\linewidth-5em][s]{\hyperlink{paper:Q0356}{\wuhao 车内声压测量的重复性不确定度研究}\dotfill\makebox[2.5em][r]{\wuhao\pageref{paper:Q0356}}}
\par
\noindent\hspace*{5em}{\kaishu\wuhao 周舟, 陈若彦, 陶建成, 邱小军}
\par\vspace{5pt}
\noindent\makebox[5em][l]{\wuhao Q0371}\makebox[\dimexpr\linewidth-5em][s]{\hyperlink{paper:Q0371}{\wuhao 面向高坝深孔的微流场感测信号处理与渗漏点反演定位}\dotfill\makebox[2.5em][r]{\wuhao\pageref{paper:Q0371}}}
\par
\noindent\hspace*{5em}{\kaishu\wuhao 梅宇健}
\par\vspace{5pt}
\noindent\makebox[5em][l]{\wuhao Q0442}\makebox[\dimexpr\linewidth-5em][s]{\hyperlink{paper:Q0442}{\wuhao 低频低噪声水听器噪声测量方法及低频压差式矢量水听器探索}\dotfill\makebox[2.5em][r]{\wuhao\pageref{paper:Q0442}}}
\par
\noindent\hspace*{5em}{\kaishu\wuhao 姜亚浩}
\par\vspace{5pt}
\noindent\makebox[5em][l]{\wuhao Q0451}\makebox[\dimexpr\linewidth-5em][s]{\hyperlink{paper:Q0451}{\wuhao 基于五传声器阻抗管的吸隔声同时测量方法}\dotfill\makebox[2.5em][r]{\wuhao\pageref{paper:Q0451}}}
\par
\noindent\hspace*{5em}{\kaishu\wuhao 邓鹏, 梅中建, 石同阳}
\par\vspace{5pt}
\vspace{6pt}

\begin{center}{\heiti\xiaosi R. 声学换能器}\end{center}
\par\vspace{6pt}
\noindent\makebox[5em][l]{\wuhao R0036}\makebox[\dimexpr\linewidth-5em][s]{\hyperlink{paper:R0036}{\wuhao 水下航行器舷侧阵自减振模块研究}\dotfill\makebox[2.5em][r]{\wuhao\pageref{paper:R0036}}}
\par
\noindent\hspace*{5em}{\kaishu\wuhao 华逸能}
\par\vspace{5pt}
\noindent\makebox[5em][l]{\wuhao R0044}\makebox[\dimexpr\linewidth-5em][s]{\hyperlink{paper:R0044}{\wuhao 基于摩擦发电技术的智能声学传感器设计}\dotfill\makebox[2.5em][r]{\wuhao\pageref{paper:R0044}}}
\par
\noindent\hspace*{5em}{\kaishu\wuhao 张笑天, 赵波}
\par\vspace{5pt}
\noindent\makebox[5em][l]{\wuhao R0066}\makebox[\dimexpr\linewidth-5em][s]{\hyperlink{paper:R0066}{\wuhao 一种双端对称四线圈驱动宽带超低频水声换能器的研究}\dotfill\makebox[2.5em][r]{\wuhao\pageref{paper:R0066}}}
\par
\noindent\hspace*{5em}{\kaishu\wuhao 孟肯, 胡建强, 刘显锋, 范进良, 张峰海, 陈玲}
\par\vspace{5pt}
\noindent\makebox[5em][l]{\wuhao R0067}\makebox[\dimexpr\linewidth-5em][s]{\hyperlink{paper:R0067}{\wuhao 不依赖气体压力补偿的电动式换能器}\dotfill\makebox[2.5em][r]{\wuhao\pageref{paper:R0067}}}
\par
\noindent\hspace*{5em}{\kaishu\wuhao 张羿双, 桑永杰}
\par\vspace{5pt}
\noindent\makebox[5em][l]{\wuhao R0070}\makebox[\dimexpr\linewidth-5em][s]{\hyperlink{paper:R0070}{\wuhao 油冷式超磁致伸缩锥形电声换能器设计​}\dotfill\makebox[2.5em][r]{\wuhao\pageref{paper:R0070}}}
\par
\noindent\hspace*{5em}{\kaishu\wuhao 李春晨, 赵能桐}
\par\vspace{5pt}
\noindent\makebox[5em][l]{\wuhao R0185}\makebox[\dimexpr\linewidth-5em][s]{\hyperlink{paper:R0185}{\wuhao 基于螺旋菲涅尔波带片结构的聚焦涡旋声换能器的研究}\dotfill\makebox[2.5em][r]{\wuhao\pageref{paper:R0185}}}
\par
\noindent\hspace*{5em}{\kaishu\wuhao 毛丹阳, 张小凤}
\par\vspace{5pt}
\noindent\makebox[5em][l]{\wuhao R0214}\makebox[\dimexpr\linewidth-5em][s]{\hyperlink{paper:R0214}{\wuhao 一种拱桥形深水换能器研究}\dotfill\makebox[2.5em][r]{\wuhao\pageref{paper:R0214}}}
\par
\noindent\hspace*{5em}{\kaishu\wuhao 马振, 范进良}
\par\vspace{5pt}
\noindent\makebox[5em][l]{\wuhao R0238}\makebox[\dimexpr\linewidth-5em][s]{\hyperlink{paper:R0238}{\wuhao 一种压电复合Cymbal结构的低频矢量水听器设计}\dotfill\makebox[2.5em][r]{\wuhao\pageref{paper:R0238}}}
\par
\noindent\hspace*{5em}{\kaishu\wuhao 兆梦荒, 张强}
\par\vspace{5pt}
\noindent\makebox[5em][l]{\wuhao R0243}\makebox[\dimexpr\linewidth-5em][s]{\hyperlink{paper:R0243}{\wuhao 一种基于机电阻抗的压电超声导波基线漂移校准方法}\dotfill\makebox[2.5em][r]{\wuhao\pageref{paper:R0243}}}
\par
\noindent\hspace*{5em}{\kaishu\wuhao 李明远, 胡正, 项延训}
\par\vspace{5pt}
\noindent\makebox[5em][l]{\wuhao R0244}\makebox[\dimexpr\linewidth-5em][s]{\hyperlink{paper:R0244}{\wuhao 一种基于零极点反馈法的亚赫兹磁电式传感器补偿网络}\dotfill\makebox[2.5em][r]{\wuhao\pageref{paper:R0244}}}
\par
\noindent\hspace*{5em}{\kaishu\wuhao 韩品, 张强}
\par\vspace{5pt}
\noindent\makebox[5em][l]{\wuhao R0248}\makebox[\dimexpr\linewidth-5em][s]{\hyperlink{paper:R0248}{\wuhao 基于AlN薄膜阵列式压电MEMS水听器的研究}\dotfill\makebox[2.5em][r]{\wuhao\pageref{paper:R0248}}}
\par
\noindent\hspace*{5em}{\kaishu\wuhao 方昊, 张强}
\par\vspace{5pt}
\noindent\makebox[5em][l]{\wuhao R0256}\makebox[\dimexpr\linewidth-5em][s]{\hyperlink{paper:R0256}{\wuhao 一种宽频带发射阵的设计与试验}\dotfill\makebox[2.5em][r]{\wuhao\pageref{paper:R0256}}}
\par
\noindent\hspace*{5em}{\kaishu\wuhao 张怡珺}
\par\vspace{5pt}
\noindent\makebox[5em][l]{\wuhao R0304}\makebox[\dimexpr\linewidth-5em][s]{\hyperlink{paper:R0304}{\wuhao 深水发射换能器性能影响研究}\dotfill\makebox[2.5em][r]{\wuhao\pageref{paper:R0304}}}
\par
\noindent\hspace*{5em}{\kaishu\wuhao 范进良}
\par\vspace{5pt}
\noindent\makebox[5em][l]{\wuhao R0312}\makebox[\dimexpr\linewidth-5em][s]{\hyperlink{paper:R0312}{\wuhao 主动声纳拖体内部结构跨尺度耦合规律及其结构优化研究}\dotfill\makebox[2.5em][r]{\wuhao\pageref{paper:R0312}}}
\par
\noindent\hspace*{5em}{\kaishu\wuhao 王世博, 吴迪, 张森}
\par\vspace{5pt}
\noindent\makebox[5em][l]{\wuhao R0325}\makebox[\dimexpr\linewidth-5em][s]{\hyperlink{paper:R0325}{\wuhao 深水换能器性能评估方法研究}\dotfill\makebox[2.5em][r]{\wuhao\pageref{paper:R0325}}}
\par
\noindent\hspace*{5em}{\kaishu\wuhao 郝浩琦, 纪京召, 马振}
\par\vspace{5pt}
\noindent\makebox[5em][l]{\wuhao R0349}\makebox[\dimexpr\linewidth-5em][s]{\hyperlink{paper:R0349}{\wuhao 多层压电薄膜构型水听器的仿真研究}\dotfill\makebox[2.5em][r]{\wuhao\pageref{paper:R0349}}}
\par
\noindent\hspace*{5em}{\kaishu\wuhao 严旭}
\par\vspace{5pt}
\noindent\makebox[5em][l]{\wuhao R0406}\makebox[\dimexpr\linewidth-5em][s]{\hyperlink{paper:R0406}{\wuhao 海底沉积物对水声换能器性能的影响仿真研究}\dotfill\makebox[2.5em][r]{\wuhao\pageref{paper:R0406}}}
\par
\noindent\hspace*{5em}{\kaishu\wuhao 常坤鹏, 潘耀宗}
\par\vspace{5pt}
\noindent\makebox[5em][l]{\wuhao R0418}\makebox[\dimexpr\linewidth-5em][s]{\hyperlink{paper:R0418}{\wuhao 两端开口圆柱形液腔低频耦合振动建模}\dotfill\makebox[2.5em][r]{\wuhao\pageref{paper:R0418}}}
\par
\noindent\hspace*{5em}{\kaishu\wuhao 刘文钊, 柴勇, 王厚琦, 张秀侦, 王秋木, 吴彤}
\par\vspace{5pt}
\noindent\makebox[5em][l]{\wuhao R0424}\makebox[\dimexpr\linewidth-5em][s]{\hyperlink{paper:R0424}{\wuhao 微型超声换能器的模拟前端接收芯片研制}\dotfill\makebox[2.5em][r]{\wuhao\pageref{paper:R0424}}}
\par
\noindent\hspace*{5em}{\kaishu\wuhao 余悦扬, 李照希, 杨银堂}
\par\vspace{5pt}
\noindent\makebox[5em][l]{\wuhao R0431}\makebox[\dimexpr\linewidth-5em][s]{\hyperlink{paper:R0431}{\wuhao 变截面超声振动系统纵-径耦合振动建模与动力学响应分析}\dotfill\makebox[2.5em][r]{\wuhao\pageref{paper:R0431}}}
\par
\noindent\hspace*{5em}{\kaishu\wuhao 丁文翔, 梁召峰}
\par\vspace{5pt}
\vspace{6pt}

\begin{center}{\heiti\xiaosi S. 智能媒体信息与网络}\end{center}
\par\vspace{6pt}
\noindent\makebox[5em][l]{\wuhao S0196}\makebox[\dimexpr\linewidth-5em][s]{\hyperlink{paper:S0196}{\wuhao 面向低算力水下节点的轻量化混合时序网络异常检测方法}\dotfill\makebox[2.5em][r]{\wuhao\pageref{paper:S0196}}}
\par
\noindent\hspace*{5em}{\kaishu\wuhao 薛志召, 刘玥帆, 张虎}
\par\vspace{5pt}
\vspace{6pt}

\begin{center}{\heiti\xiaosi U. 计算声学}\end{center}
\par\vspace{6pt}
\noindent\makebox[5em][l]{\wuhao U0037}\makebox[\dimexpr\linewidth-5em][s]{\hyperlink{paper:U0037}{\wuhao 基于算子分裂方法的参量声场数值建模研究}\dotfill\makebox[2.5em][r]{\wuhao\pageref{paper:U0037}}}
\par
\noindent\hspace*{5em}{\kaishu\wuhao 杨焕然, 杨德森, 施浩康, 时洁}
\par\vspace{5pt}
\noindent\makebox[5em][l]{\wuhao U0038}\makebox[\dimexpr\linewidth-5em][s]{\hyperlink{paper:U0038}{\wuhao 运动介质中声传播问题的快速边界元法研究}\dotfill\makebox[2.5em][r]{\wuhao\pageref{paper:U0038}}}
\par
\noindent\hspace*{5em}{\kaishu\wuhao 刘学良, 文新}
\par\vspace{5pt}
\noindent\makebox[5em][l]{\wuhao U0137}\makebox[\dimexpr\linewidth-5em][s]{\hyperlink{paper:U0137}{\wuhao 基于融合物理先验的多输入傅里叶神经算子开展浅海声场端到端预测}\dotfill\makebox[2.5em][r]{\wuhao\pageref{paper:U0137}}}
\par
\noindent\hspace*{5em}{\kaishu\wuhao 张运祥, 王勇献, 屠厚旺, 孙昭, 王泽宇}
\par\vspace{5pt}
\noindent\makebox[5em][l]{\wuhao U0151}\makebox[\dimexpr\linewidth-5em][s]{\hyperlink{paper:U0151}{\wuhao 水下复杂壳体外部声振耦合分析的 Chebyshev 谱元-谱边界元方法}\dotfill\makebox[2.5em][r]{\wuhao\pageref{paper:U0151}}}
\par
\noindent\hspace*{5em}{\kaishu\wuhao 赵天通, 叶天贵, 陈玉坤, 谢祥, 靳国永, 孙谦}
\par\vspace{5pt}
\noindent\makebox[5em][l]{\wuhao U0183}\makebox[\dimexpr\linewidth-5em][s]{\hyperlink{paper:U0183}{\wuhao 基于物理信息核函数神经网络的海洋结构声振耦合研究}\dotfill\makebox[2.5em][r]{\wuhao\pageref{paper:U0183}}}
\par
\noindent\hspace*{5em}{\kaishu\wuhao 习强}
\par\vspace{5pt}
\noindent\makebox[5em][l]{\wuhao U0200}\makebox[\dimexpr\linewidth-5em][s]{\hyperlink{paper:U0200}{\wuhao 夹层结构隔声性能预测的高效声-固耦合有限元程序}\dotfill\makebox[2.5em][r]{\wuhao\pageref{paper:U0200}}}
\par
\noindent\hspace*{5em}{\kaishu\wuhao 许英浩, 范一良, 王普缘, 季振林}
\par\vspace{5pt}
\noindent\makebox[5em][l]{\wuhao U0245}\makebox[\dimexpr\linewidth-5em][s]{\hyperlink{paper:U0245}{\wuhao 有流状态下消声器宽频声学性能预测的自适应有限元方法}\dotfill\makebox[2.5em][r]{\wuhao\pageref{paper:U0245}}}
\par
\noindent\hspace*{5em}{\kaishu\wuhao 范一良, 季振林}
\par\vspace{5pt}
\noindent\makebox[5em][l]{\wuhao U0252}\makebox[\dimexpr\linewidth-5em][s]{\hyperlink{paper:U0252}{\wuhao 自动频率无关完美匹配层}\dotfill\makebox[2.5em][r]{\wuhao\pageref{paper:U0252}}}
\par
\noindent\hspace*{5em}{\kaishu\wuhao 谢祥, 靳国永}
\par\vspace{5pt}
\noindent\makebox[5em][l]{\wuhao U0318}\makebox[\dimexpr\linewidth-5em][s]{\hyperlink{paper:U0318}{\wuhao 参数可编程完美声涡旋的声全息构建及其微粒操控}\dotfill\makebox[2.5em][r]{\wuhao\pageref{paper:U0318}}}
\par
\noindent\hspace*{5em}{\kaishu\wuhao 全浩, 李昕珈, 孟龙, 刘秀芳}
\par\vspace{5pt}
\noindent\hangindent=5em\hangafter=1\makebox[5em][l]{\wuhao U0335}\hyperlink{paper:U0335}{\wuhao Fluid Structure Interaction Simulation and Forward Looking Sonar Detection of Moving Jellyfish: A SPIM-IB-LBM Approach}
\par\noindent\hspace*{5em}\dotfill\makebox[2.5em][r]{\wuhao\pageref{paper:U0335}}
\par\vspace{2pt}
\noindent\hspace*{5em}{\kaishu\wuhao Irfan Ali, 肖永锟, 鲁建华}
\par\vspace{5pt}
\noindent\makebox[5em][l]{\wuhao U0378}\makebox[\dimexpr\linewidth-5em][s]{\hyperlink{paper:U0378}{\wuhao 基于深度神经网络的水中大尺度圆柱壳声目标强度快速预测}\dotfill\makebox[2.5em][r]{\wuhao\pageref{paper:U0378}}}
\par
\noindent\hspace*{5em}{\kaishu\wuhao 刘子豪, 赵宏刚}
\par\vspace{5pt}
\noindent\makebox[5em][l]{\wuhao U0390}\makebox[\dimexpr\linewidth-5em][s]{\hyperlink{paper:U0390}{\wuhao 基于虚单元法的三维声腔模态与频响特性分析}\dotfill\makebox[2.5em][r]{\wuhao\pageref{paper:U0390}}}
\par
\noindent\hspace*{5em}{\kaishu\wuhao 倪帅龙, 谢祥, 叶天贵, 靳国永}
\par\vspace{5pt}
\noindent\makebox[5em][l]{\wuhao U0447}\makebox[\dimexpr\linewidth-5em][s]{\hyperlink{paper:U0447}{\wuhao 基于Transformer架构的分段式管束穿孔板逆向设计}\dotfill\makebox[2.5em][r]{\wuhao\pageref{paper:U0447}}}
\par
\noindent\hspace*{5em}{\kaishu\wuhao 相同林, 梅中建, 石同阳}
\par\vspace{5pt}
\noindent\makebox[5em][l]{\wuhao U0468}\makebox[\dimexpr\linewidth-5em][s]{\hyperlink{paper:U0468}{\wuhao 声学快速多极波叠加法的动态稀疏近似逆预处理技术}\dotfill\makebox[2.5em][r]{\wuhao\pageref{paper:U0468}}}
\par
\noindent\hspace*{5em}{\kaishu\wuhao 刘烨崴, 向宇, 李鑫, 琚航}
\par\vspace{5pt}
\vspace{6pt}

\begin{center}{\heiti\xiaosi V. 生物声学}\end{center}
\par\vspace{6pt}
\noindent\makebox[5em][l]{\wuhao V0001}\makebox[\dimexpr\linewidth-5em][s]{\hyperlink{paper:V0001}{\wuhao 测试}\dotfill\makebox[2.5em][r]{\wuhao\pageref{paper:V0001}}}
\par
\noindent\hspace*{5em}{\kaishu\wuhao 胡同}
\par\vspace{5pt}
\noindent\makebox[5em][l]{\wuhao V0093}\makebox[\dimexpr\linewidth-5em][s]{\hyperlink{paper:V0093}{\wuhao 基于海豚Click的仿生隐蔽探通一体化研究}\dotfill\makebox[2.5em][r]{\wuhao\pageref{paper:V0093}}}
\par
\noindent\hspace*{5em}{\kaishu\wuhao 华志远, 黄子豪, 彭旭名, 夏源, 张宇, 宋忠长}
\par\vspace{5pt}
\noindent\makebox[5em][l]{\wuhao V0094}\makebox[\dimexpr\linewidth-5em][s]{\hyperlink{paper:V0094}{\wuhao 中华白海豚（Sousa chinensis）的signature whistle指向性研究}\dotfill\makebox[2.5em][r]{\wuhao\pageref{paper:V0094}}}
\par
\noindent\hspace*{5em}{\kaishu\wuhao 李玉鹏, 吴福星, 张洁瑜, 彭旭名, 孟凡祎, 傅伟杰, 张飞, 张闯, 项文杰, 欧文湛, 张宇, 宋忠长}
\par\vspace{5pt}
\noindent\makebox[5em][l]{\wuhao V0146}\makebox[\dimexpr\linewidth-5em][s]{\hyperlink{paper:V0146}{\wuhao 大黄鱼发声机制建模与影响因素研究}\dotfill\makebox[2.5em][r]{\wuhao\pageref{paper:V0146}}}
\par
\noindent\hspace*{5em}{\kaishu\wuhao 彭学舟, 惠健刚, 欧文湛, 李宁, 宋忠长}
\par\vspace{5pt}
\noindent\makebox[5em][l]{\wuhao V0182}\makebox[\dimexpr\linewidth-5em][s]{\hyperlink{paper:V0182}{\wuhao 面向肠梗阻识别的肠鸣音稳定声学特征筛选}\dotfill\makebox[2.5em][r]{\wuhao\pageref{paper:V0182}}}
\par
\noindent\hspace*{5em}{\kaishu\wuhao 张嘉怡, 闫靓}
\par\vspace{5pt}
\noindent\makebox[5em][l]{\wuhao V0188}\makebox[\dimexpr\linewidth-5em][s]{\hyperlink{paper:V0188}{\wuhao 昆虫声振通讯的行为特征与神经机制及其在害虫防控中的应用思路}\dotfill\makebox[2.5em][r]{\wuhao\pageref{paper:V0188}}}
\par
\noindent\hspace*{5em}{\kaishu\wuhao 张欣杨, 熊时姣}
\par\vspace{5pt}
\noindent\makebox[5em][l]{\wuhao V0204}\makebox[\dimexpr\linewidth-5em][s]{\hyperlink{paper:V0204}{\wuhao 仿江豚水声接收器件设计}\dotfill\makebox[2.5em][r]{\wuhao\pageref{paper:V0204}}}
\par
\noindent\hspace*{5em}{\kaishu\wuhao 欧文湛, 宋忠长, 叶馨, 张金虎, 张闯, 张宇}
\par\vspace{5pt}
\noindent\makebox[5em][l]{\wuhao V0212}\makebox[\dimexpr\linewidth-5em][s]{\hyperlink{paper:V0212}{\wuhao 空间混叠条件下麦克风阵列并发多鸟鸣无模糊3D测向方法}\dotfill\makebox[2.5em][r]{\wuhao\pageref{paper:V0212}}}
\par
\noindent\hspace*{5em}{\kaishu\wuhao 刘力铭, 肖越, 赵兆, 许志勇}
\par\vspace{5pt}
\noindent\makebox[5em][l]{\wuhao V0235}\makebox[\dimexpr\linewidth-5em][s]{\hyperlink{paper:V0235}{\wuhao 基于多模态检测与生态声学的生物多样性平台开发}\dotfill\makebox[2.5em][r]{\wuhao\pageref{paper:V0235}}}
\par
\noindent\hspace*{5em}{\kaishu\wuhao 张汇宇}
\par\vspace{5pt}
\noindent\makebox[5em][l]{\wuhao V0246}\makebox[\dimexpr\linewidth-5em][s]{\hyperlink{paper:V0246}{\wuhao 基于Click信号复杂度特征的混响背景下目标检测方法研究}\dotfill\makebox[2.5em][r]{\wuhao\pageref{paper:V0246}}}
\par
\noindent\hspace*{5em}{\kaishu\wuhao 何书婷, 宋忠长, 傅伟杰, 张金虎, 张宇}
\par\vspace{5pt}
\noindent\makebox[5em][l]{\wuhao V0279}\makebox[\dimexpr\linewidth-5em][s]{\hyperlink{paper:V0279}{\wuhao 厦门湾中华白海豚click脉冲串的提取与动态特征分析}\dotfill\makebox[2.5em][r]{\wuhao\pageref{paper:V0279}}}
\par
\noindent\hspace*{5em}{\kaishu\wuhao 吴景怡, 项文杰, 吴福星, 宋忠长}
\par\vspace{5pt}
\noindent\makebox[5em][l]{\wuhao V0282}\makebox[\dimexpr\linewidth-5em][s]{\hyperlink{paper:V0282}{\wuhao 脉冲数量对仿生多脉冲目标识别性能的影响研究}\dotfill\makebox[2.5em][r]{\wuhao\pageref{paper:V0282}}}
\par
\noindent\hspace*{5em}{\kaishu\wuhao 彭旭名, 欧文湛, 项文杰, 石明杰, 张博宇, 宋忠长, 张宇}
\par\vspace{5pt}
\noindent\makebox[5em][l]{\wuhao V0290}\makebox[\dimexpr\linewidth-5em][s]{\hyperlink{paper:V0290}{\wuhao 厦门湾齿鲸click的实时检测与分类算法}\dotfill\makebox[2.5em][r]{\wuhao\pageref{paper:V0290}}}
\par
\noindent\hspace*{5em}{\kaishu\wuhao 项文杰, 彭旭名, 吴福星, 宋忠长}
\par\vspace{5pt}
\noindent\makebox[5em][l]{\wuhao V0307}\makebox[\dimexpr\linewidth-5em][s]{\hyperlink{paper:V0307}{\wuhao 香港水域枪虾声学活动的时空模式研究}\dotfill\makebox[2.5em][r]{\wuhao\pageref{paper:V0307}}}
\par
\noindent\hspace*{5em}{\kaishu\wuhao 黎小琪, Lindsay Porte, Naomi Brannan, 欧文湛, 宋忠长}
\par\vspace{5pt}
\noindent\makebox[5em][l]{\wuhao V0332}\makebox[\dimexpr\linewidth-5em][s]{\hyperlink{paper:V0332}{\wuhao 真实几何下弯曲多层动脉粥样硬化血管的光声成像仿真研究}\dotfill\makebox[2.5em][r]{\wuhao\pageref{paper:V0332}}}
\par
\noindent\hspace*{5em}{\kaishu\wuhao 任妤晴, 曹辉}
\par\vspace{5pt}
\noindent\makebox[5em][l]{\wuhao V0343}\makebox[\dimexpr\linewidth-5em][s]{\hyperlink{paper:V0343}{\wuhao 基于脑电和眼电特征的警觉度研究}\dotfill\makebox[2.5em][r]{\wuhao\pageref{paper:V0343}}}
\par
\noindent\hspace*{5em}{\kaishu\wuhao 张天琼, 伍南西, 李忠杰, 李佳蔚, 甘霖}
\par\vspace{5pt}
\noindent\makebox[5em][l]{\wuhao V0368}\makebox[\dimexpr\linewidth-5em][s]{\hyperlink{paper:V0368}{\wuhao 长江江豚下颌发育对声接收的影响}\dotfill\makebox[2.5em][r]{\wuhao\pageref{paper:V0368}}}
\par
\noindent\hspace*{5em}{\kaishu\wuhao 张闯, 陈宇维, 欧文湛, 余力涵, 宋忠长, 郝玉江, 郑劲松, 王克雄}
\par\vspace{5pt}
\noindent\makebox[5em][l]{\wuhao V0387}\makebox[\dimexpr\linewidth-5em][s]{\hyperlink{paper:V0387}{\wuhao 繁殖季节的共存策略——新疆中亚侧褶蛙与黑斑侧褶蛙的声学生态位分化}\dotfill\makebox[2.5em][r]{\wuhao\pageref{paper:V0387}}}
\par
\noindent\hspace*{5em}{\kaishu\wuhao 龚蒙蒙, 周鹿}
\par\vspace{5pt}
\noindent\hangindent=5em\hangafter=1\makebox[5em][l]{\wuhao V0410}\hyperlink{paper:V0410}{\wuhao 厦门湾船舶噪声特性及其对中华白海豚（Sousa chinensis）潜在影响研究}
\par\noindent\hspace*{5em}\dotfill\makebox[2.5em][r]{\wuhao\pageref{paper:V0410}}
\par\vspace{2pt}
\noindent\hspace*{5em}{\kaishu\wuhao 傅伟杰, 宋忠长, 彭旭名, 吴福星, 张飞, 欧文湛, 张宇}
\par\vspace{5pt}
\noindent\makebox[5em][l]{\wuhao V0414}\makebox[\dimexpr\linewidth-5em][s]{\hyperlink{paper:V0414}{\wuhao 福建东山岛珊瑚礁光合气泡声学信号特征及其时空变化}\dotfill\makebox[2.5em][r]{\wuhao\pageref{paper:V0414}}}
\par
\noindent\hspace*{5em}{\kaishu\wuhao 张飞, 宋忠长}
\par\vspace{5pt}
\noindent\makebox[5em][l]{\wuhao V0467}\makebox[\dimexpr\linewidth-5em][s]{\hyperlink{paper:V0467}{\wuhao 低强度脉冲超声干预对后肢去负荷小鼠肌肉萎缩和纤维化的影响}\dotfill\makebox[2.5em][r]{\wuhao\pageref{paper:V0467}}}
\par
\noindent\hspace*{5em}{\kaishu\wuhao 禹娅南, 唐量}
\par\vspace{5pt}
\vspace{6pt}

\begin{center}{\heiti\xiaosi W.声学材料（声超构材料、压电陶瓷材料、单晶材料等）}\end{center}
\par\vspace{6pt}
\noindent\makebox[5em][l]{\wuhao W0023}\makebox[\dimexpr\linewidth-5em][s]{\hyperlink{paper:W0023}{\wuhao 具有超高εr、d33与高Tm的PNN-PMN-PZT四元压电陶瓷}\dotfill\makebox[2.5em][r]{\wuhao\pageref{paper:W0023}}}
\par
\noindent\hspace*{5em}{\kaishu\wuhao 项光磊}
\par\vspace{5pt}
\noindent\makebox[5em][l]{\wuhao W0026}\makebox[\dimexpr\linewidth-5em][s]{\hyperlink{paper:W0026}{\wuhao 声学黑洞筏架系统宽频抑振特性研究}\dotfill\makebox[2.5em][r]{\wuhao\pageref{paper:W0026}}}
\par
\noindent\hspace*{5em}{\kaishu\wuhao 王帅, 刘璞, 朱金晏, 马德树, 夏梦, 童宗鹏}
\par\vspace{5pt}
\noindent\hangindent=5em\hangafter=1\makebox[5em][l]{\wuhao W0050}\hyperlink{paper:W0050}{\wuhao Hafnium-doped (Ba, Ca) (Zr, Ti) O3 lead-free focused ultrasound transducer and quantitative neuromodulation in Parkinson’s disease}
\par\noindent\hspace*{5em}\dotfill\makebox[2.5em][r]{\wuhao\pageref{paper:W0050}}
\par\vspace{2pt}
\noindent\hspace*{5em}{\kaishu\wuhao Junkai Tong, Renjie Wu}
\par\vspace{5pt}
\noindent\makebox[5em][l]{\wuhao W0059}\makebox[\dimexpr\linewidth-5em][s]{\hyperlink{paper:W0059}{\wuhao 面向弯曲波轨道角动量感知与通信的弹性超表面设计}\dotfill\makebox[2.5em][r]{\wuhao\pageref{paper:W0059}}}
\par
\noindent\hspace*{5em}{\kaishu\wuhao 赵流现}
\par\vspace{5pt}
\noindent\makebox[5em][l]{\wuhao W0085}\makebox[\dimexpr\linewidth-5em][s]{\hyperlink{paper:W0085}{\wuhao 基于声学黑洞吸声体的声能量采集结构的设计与实验研究}\dotfill\makebox[2.5em][r]{\wuhao\pageref{paper:W0085}}}
\par
\noindent\hspace*{5em}{\kaishu\wuhao 彭丽华, 毛崎波}
\par\vspace{5pt}
\noindent\makebox[5em][l]{\wuhao W0087}\makebox[\dimexpr\linewidth-5em][s]{\hyperlink{paper:W0087}{\wuhao 双孔隙率声学超构材料的低频吸声特性研究}\dotfill\makebox[2.5em][r]{\wuhao\pageref{paper:W0087}}}
\par
\noindent\hspace*{5em}{\kaishu\wuhao 楼华鼎, 顾震宇}
\par\vspace{5pt}
\noindent\makebox[5em][l]{\wuhao W0140}\makebox[\dimexpr\linewidth-5em][s]{\hyperlink{paper:W0140}{\wuhao 基于数据驱动方法的超构材料吸声特性研究}\dotfill\makebox[2.5em][r]{\wuhao\pageref{paper:W0140}}}
\par
\noindent\hspace*{5em}{\kaishu\wuhao 赵赫辉, 杨金水, 李鸿泽}
\par\vspace{5pt}
\noindent\makebox[5em][l]{\wuhao W0158}\makebox[\dimexpr\linewidth-5em][s]{\hyperlink{paper:W0158}{\wuhao 吸声材料布置方式对声学黑洞吸声体性能的影响}\dotfill\makebox[2.5em][r]{\wuhao\pageref{paper:W0158}}}
\par
\noindent\hspace*{5em}{\kaishu\wuhao 胡肖杭, 毛崎波}
\par\vspace{5pt}
\noindent\makebox[5em][l]{\wuhao W0413}\makebox[\dimexpr\linewidth-5em][s]{\hyperlink{paper:W0413}{\wuhao 基于柔性智能材料的可编程低频吸声超材料}\dotfill\makebox[2.5em][r]{\wuhao\pageref{paper:W0413}}}
\par
\noindent\hspace*{5em}{\kaishu\wuhao 魏仙凤, 王震, 韩帅, 常龙飞}
\par\vspace{5pt}
\vspace{6pt}

\clearpage

\printpaper
{A. 水下声学}
{A0094}
{浅海环境中低频声传播的三维耦合简正波建模与实验验证}
{%
  张三\ns{1},
  李四\ns{2}\ca,
  王五\ns{1}
}
{%
  (1) 中国科学院声学研究所，北京 100190\\
  (2) 哈尔滨工程大学水声技术重点实验室，哈尔滨 150001
}
{%
  \textsuperscript{*}lisi@example.com
}
{%
  浅海声传播受海底地形、声速剖面和海面波浪等多重因素耦合影响，传统的二维传播模型难以准确描述水平折射效应。本文提出一种三维耦合简正波方法，将水平变化环境离散为多个距离无关段，通过段间模式耦合矩阵实现全局声场求解。数值仿真表明，该方法在内波环境和大陆坡海域均能有效捕获水平折射引起的声能水平偏转现象，与三维抛物方程方法计算结果一致，计算效率提高约5倍。在南海北部海域开展了低频宽带声传播实验，实验数据与模型预测的传播损失在500\,Hz至2\,kHz频段内均方误差小于2.5\,dB，验证了模型的有效性。
}
{%
  三维声传播；耦合简正波；水平折射；浅海波导；实验验证
}
{A0094}

\end{document}
